% 第4章 LaTeX 章节片段。
% 建议主文档使用 XeLaTeX 或 LuaLaTeX 编译中文，并加载：
% \usepackage{amsmath,amssymb,booktabs,tabularx,enumitem,hyperref}
% 本文件不含 \documentclass、\begin{document} 或参考文献配置，
% 请在主文档中使用 % 第4章 LaTeX 章节片段。
% 建议主文档使用 XeLaTeX 或 LuaLaTeX 编译中文，并加载：
% \usepackage{amsmath,amssymb,booktabs,tabularx,enumitem,hyperref}
% 本文件不含 \documentclass、\begin{document} 或参考文献配置，
% 请在主文档中使用 % 第4章 LaTeX 章节片段。
% 建议主文档使用 XeLaTeX 或 LuaLaTeX 编译中文，并加载：
% \usepackage{amsmath,amssymb,booktabs,tabularx,enumitem,hyperref}
% 本文件不含 \documentclass、\begin{document} 或参考文献配置，
% 请在主文档中使用 % 第4章 LaTeX 章节片段。
% 建议主文档使用 XeLaTeX 或 LuaLaTeX 编译中文，并加载：
% \usepackage{amsmath,amssymb,booktabs,tabularx,enumitem,hyperref}
% 本文件不含 \documentclass、\begin{document} 或参考文献配置，
% 请在主文档中使用 \input{thesis_chapter4_draft_zh.tex}。

\chapter{自适应标定与人体手臂姿态估计}
\label{chap:adaptive_calibration_arm_estimation}

\section{安装标定与姿态估计相关研究}
\label{sec:ch4_related_calibration}

可穿戴传感器报告的是传感器外壳相对参考坐标系的位姿，而人体运动分析需要的是肢段坐标系或
关节变量。两者之间的传感器--肢段关系会随佩戴位置和朝向改变，是可穿戴上肢测量中的核心
误差来源。国际生物力学学会（ISB）给出了肩、肘、腕和手部坐标系及关节运动报告的推荐
定义，为不同测量系统之间统一轴向和角度语义提供了基础~\cite{wu2005isb}。Picerno 等通过
解剖标志点建立上肢惯性传感器与肢段坐标系之间的关系，并使用光学
系统验证肩、肘运动学~\cite{picerno2019upper}。Seel 等利用运动学约束从功能动作中识别关节轴，
说明不必要求传感器外壳与解剖轴精确平行~\cite{seel2014imu}。Zabat 等进一步表明，相比单一
静态姿势，多姿态信息可以减小安装偏差和软组织运动引起的角度误差~\cite{zabat2019multiple}。
近期研究还指出，待测系统与参考系统若采用不同的传感器--肢段定义，会形成明显的系统性关节角
偏差~\cite{adjel2025influence}。

现有标定方法大体可分为静态、功能性和多姿态方法。静态方法操作简单，但对参考姿势执行误差
敏感；功能性方法通过重复运动拟合关节轴，降低了对安装朝向的要求，但依赖充分且近似单轴的
动作激励；多姿态方法利用若干互补方向联合估计安装关系，并可通过各姿态残差发现错误动作或
绑带移动。对肩部三自由度、肘部屈伸和前臂轴向旋转同时估计时，单一姿态不能提供足够约束，
尤其在大臂与前臂接近共线时，肘角和轴向旋转的可观测性会降低。

Cutti 等使用惯性与磁传感器进行肩、肘动态测量的研究也表明，统一的肢段定义、安装标定和
磁环境检查是获得可解释上肢角度的共同前提~\cite{cutti2008ambulatory}。

电磁设备还存在“设备场标定”与“人体安装标定”两个不同层次。前者从已知空间位姿建立电磁
观测与接收器位姿之间的映射，并处理工作空间内的位置相关误差；相关综述与标准化评估研究
均指出，金属畸变、工作距离和空间位置会共同影响电磁跟踪误差
~\cite{kindratenko2000survey,frantz2003accuracy,hummel2006evaluation,franz2014electromagnetic,
maclachlan2025calibration}。后者在设备位姿已经
可用的前提下，估计接收器外壳与人体肢段之间的关系。本文不修改设备固件或制造商场模型，而是
针对每次佩戴执行多姿态安装级标定，并把姿态残差、位置模型残差和重复性作为进入在线估计前的
质量门控。

\section{问题描述与总体思路}
\label{sec:ch4_problem_overview}

本研究使用一套电磁位姿跟踪系统获取操作者左臂运动。电磁发射器固定于操作者躯干，两个
接收器分别固定于左大臂和左前臂。系统以发射器坐标系为共同参考，以约 $50\,\mathrm{Hz}$
的频率输出两个接收器的三维位置和姿态。与基于人体外观的视觉方法相比，电磁测量不依赖
纹理和光照，也不存在人体自身造成的直接视觉遮挡。然而，接收器在每次佩戴后的安装朝向
并不相同，绑带和软组织还会使接收器运动偏离理想的刚体关节模型。因此，原始接收器位姿
不能直接作为机器人关节角使用，必须先完成针对本次安装的标定。这里的“原始位姿”专指
设备直接报告的接收器外壳相对发射器的位置和旋转；它尚未去除外壳安装角，也尚未转换为
人体大臂、前臂方向或关节角。

本章所研究的问题可表述为：给定大臂接收器和前臂接收器在电磁发射器坐标系中的同步位姿，
估计人体左臂的五个运动量
\begin{equation}
  \boldsymbol q_h =
  \begin{bmatrix}
    q_p & q_r & q_y & q_e & q_w
  \end{bmatrix}^{\mathrm T},
  \label{eq:ch4_human_arm_state}
\end{equation}
其中，$q_p$、$q_r$ 和 $q_y$ 分别表示肩部前举/后摆、外展/内收和轴向旋转，$q_e$
表示肘部由约 $90^\circ$ 屈曲姿态向伸直方向的变化量，$q_w$ 表示前臂旋前/旋后。
上述变量沿用 Unitree G1 左臂的关节语义，但其数值首先描述人体动作，随后才在重定向和
安全层中被限制为机器人可执行目标。

需要强调的是，式~\eqref{eq:ch4_human_arm_state} 中的五个量是本研究所选运动学模型的
广义坐标，而不是由传感器直接测得的五个解剖学关节角。人体肩胛骨、锁骨、腕部屈伸以及
关节中心的平移均未被单独建模。本文的目标是构造一组与 G1 左臂关节语义一致、能够由当前
两接收器配置稳定估计的低维运动描述，而不是恢复完整的人体上肢骨骼模型。

系统采用“安装级自适应标定”思路。这里的“自适应”并非在线学习或自适应控制，而是指每次
重新佩戴设备后，系统根据本次安装采集的参考动作重新辨识：
\begin{enumerate}[label=\arabic*)]
  \item 电磁发射器坐标系到人体标准坐标系的映射；
  \item 接收器局部坐标系中的肢段纵轴；
  \item 大臂接收器位置与人体大臂方向之间的映射；
  \item 前臂旋前/旋后的局部旋转轴和零位；
  \item 肘角几何估计的增益和偏置。
\end{enumerate}

在线阶段只使用本次安装生成的参数。任一接收器或发射器被拆卸、重新绑扎或在支架中移动
后，原标定不再有效，必须重新执行标定流程。
换言之，“安装级”表示参数在一次佩戴完成后被辨识，并在该次运行期间保持固定；它不表示
参数随每一帧在线更新。这样既允许系统适应重复佩戴造成的安装差异，又避免在线更新参数时
把操作者的真实动作错误地吸收到标定模型中。

\subsection{坐标系与符号}
\label{subsec:ch4_frames_notation}

本章使用表~\ref{tab:ch4_coordinate_frames} 所列坐标系。
\begin{table}[htbp]
  \centering
  \caption{人体手臂估计使用的坐标系}
  \label{tab:ch4_coordinate_frames}
  \begin{tabularx}{\linewidth}{@{}llX@{}}
    \toprule
    符号 & 英文含义 & 物理含义 \\
    \midrule
    $\{T\}$ & Transmitter & 电磁发射器坐标系 \\
    $\{B\}$ & Body & 人体标准坐标系，与机器人约定一致：$+x$ 向前、$+y$ 向左、$+z$ 向上 \\
    $\{U\}$ & Upper arm & 大臂接收器局部坐标系 \\
    $\{F\}$ & Forearm & 前臂接收器局部坐标系 \\
    \bottomrule
  \end{tabularx}
\end{table}

记 ${}^{T}\boldsymbol p_U$ 和 ${}^{T}\boldsymbol p_F$ 为两个接收器在 $\{T\}$ 中的位置，
${}^{T}\boldsymbol R_U$ 和 ${}^{T}\boldsymbol R_F$ 为将接收器局部向量旋转到 $\{T\}$ 的
旋转矩阵。上标表示量所在的表达坐标系，下标表示被描述的对象；例如
${}^{T}\boldsymbol p_U$ 表示“大臂接收器的位置，用发射器坐标系表达”。
${}^{B}\boldsymbol u$ 和 ${}^{B}\boldsymbol f$ 分别表示人体大臂和前臂在
$\{B\}$ 中的单位方向。本文所称“肢段”是由相邻关节界定、在当前模型中近似视为刚体的
身体部分：大臂肢段连接肩与肘，前臂肢段连接肘与腕。所谓单位方向，是指只保留肢段的朝向
而忽略其长度，即
\begin{equation}
  \lVert{}^{B}\boldsymbol u\rVert_2=1,
  \qquad
  \lVert{}^{B}\boldsymbol f\rVert_2=1.
  \label{eq:ch4_unit_segment_directions}
\end{equation}
将单位向量平移到 $\{B\}$ 的数学原点后，其终点位于单位球面上。向量是自由向量，平移
起点不会改变其方向；因此这里的单位球面只是统一表示所有三维方向的数学空间，并不要求
人体肩关节位于 $\{B\}$ 原点，也不表示接收器在物理上沿半径为 1 的球面运动。

为与 G1 关节模型保持一致，本研究将零位大臂方向和零位前臂方向定义为
\begin{equation}
  \boldsymbol a_{u,0}=\begin{bmatrix}0&0&-1\end{bmatrix}^{\mathrm T},
  \qquad
  \boldsymbol a_{f,0}=\begin{bmatrix}1&0&0\end{bmatrix}^{\mathrm T}.
  \label{eq:ch4_zero_segment_axes}
\end{equation}
因此，零位对应“大臂自然下垂、肘部弯曲约 $90^\circ$、前臂水平向前”的姿态。
这里的“零位方向”是运动学模型在所有广义坐标均为零时规定的参考方向，并非接收器局部
坐标轴，也不等同于医学上的
标准解剖姿势。当所有模型关节变量为零时，大臂和前臂方向分别等于
$\boldsymbol a_{u,0}$ 和 $\boldsymbol a_{f,0}$；其他姿态均由关节旋转作用于这两个参考
方向得到。图~\ref{fig:ch4_frames_zero_pose} 汇总了坐标轴、接收器和零位方向之间的关系。

\begin{figure}[htbp]
  \centering
  \fbox{\parbox{0.90\linewidth}{\centering
    示意图内容：人体左臂处于 DOWN 零位；标出 $\{B\}$ 的 $+x$ 前、$+y$ 左、$+z$ 上，
    发射器坐标系 $\{T\}$，大臂与前臂接收器坐标系 $\{U\}$、$\{F\}$，以及
    $\boldsymbol a_{u,0}=[0,0,-1]^{\mathrm T}$ 和
    $\boldsymbol a_{f,0}=[1,0,0]^{\mathrm T}$。}}
  \caption{双接收器安装、坐标系与运动学零位示意（待替换为正式绘图）}
  \label{fig:ch4_frames_zero_pose}
\end{figure}

\subsection{五自由度运动学模型}
\label{subsec:ch4_five_dof_model}

本研究用肩部三个旋转、肘部一个旋转以及前臂一个轴向旋转描述左臂运动，即
\begin{equation}
  3\ \text{(shoulder)}+1\ \text{(elbow)}+1\ \text{(forearm twist)}
  =5\ \text{DOF}.
  \label{eq:ch4_five_dof_count}
\end{equation}
肩部三个变量 $q_p$（pitch，前举/后摆）、$q_r$（roll，外展/内收）和 $q_y$（yaw，绕
大臂纵轴旋转）共同确定大臂的空间方向及肘平面朝向；$q_e$ 描述肘部由弯曲向伸直的变化，
$q_w$ 描述前臂旋前/旋后。本模型没有包含 G1 腕部的 pitch 和 yaw，也没有单独描述肩胛带
运动，因此应理解为面向动作重定向的五自由度表示，而非完整的人体上肢解剖模型。

肩部和肘部的正向运动学采用
\begin{align}
  {}^{B}\boldsymbol R_S
  &=\boldsymbol R_y(q_p)\boldsymbol R_x(q_r)\boldsymbol R_z(q_y),
  \label{eq:ch4_shoulder_rotation}\\
  {}^{B}\boldsymbol R_E
  &={}^B\boldsymbol R_S\boldsymbol R_y(q_e),
  \label{eq:ch4_forearm_rotation}
\end{align}
其中，${}^{B}\boldsymbol R_S$ 表示肩部旋转，${}^{B}\boldsymbol R_E$ 表示不包含腕部
旋转的前臂坐标系。相应的肢段方向为
\begin{equation}
  {}^{B}\boldsymbol u={}^B\boldsymbol R_S\boldsymbol a_{u,0},
  \qquad
  {}^{B}\boldsymbol f={}^B\boldsymbol R_E\boldsymbol a_{f,0}.
  \label{eq:ch4_segment_directions}
\end{equation}

式~\eqref{eq:ch4_shoulder_rotation} 的乘法顺序是模型定义的一部分，不能任意交换。本文采用
列向量和主动旋转约定，因此右侧旋转先作用：按矩阵作用顺序，零位肢段依次接受
$\boldsymbol R_z(q_y)$、$\boldsymbol R_x(q_r)$ 和 $\boldsymbol R_y(q_p)$。各轴均以本章统一的
矩阵约定解释，而不是在推导中临时更换为另一套欧拉角约定。矩阵乘法不满足交换律，改变顺序会改变同一组
角度对应的空间姿态。本顺序与工程中 G1 左肩三个关节的语义保持一致，而不是声称它是唯一的
人体解剖角定义；解剖学结果的轴序应另按 ISB 推荐报告~\cite{wu2005isb}。

该模型的优点是人体估计量与 G1 五个左臂关节具有一致的语义，避免了先恢复完整人体骨架
再进行第二次逆运动学求解。模型的限制是：仅由两个肢段方向不能在所有姿态下唯一分离肩部
轴向旋转与前臂旋转。第~\ref{sec:ch4_observability_robustness} 节将对这一可观测性问题进行
专门处理。

\subsection{测量信息与五自由度状态的对应关系}
\label{subsec:ch4_measurement_dof_correspondence}

为了说明两个接收器为何能够支持五个运动量的估计，需要区分“数据分量个数”和“独立
自由度个数”。一个三维单位方向属于单位球面
\begin{equation}
  S^2=\left\{\boldsymbol v\in\mathbb R^3\mid
  \boldsymbol v^{\mathrm T}\boldsymbol v=1\right\}.
  \label{eq:ch4_unit_sphere_definition}
\end{equation}
虽然 $\boldsymbol v$ 具有三个坐标分量，但单位长度约束消除了一个独立变量，因此一个方向
只有两个自由度。直观上，这两个自由度类似球面上的经度和纬度；它们能够确定肢段指向，
却不能确定肢段绕自身纵轴旋转了多少。

旋转矩阵属于特殊正交群
\begin{equation}
  SO(3)=\left\{\boldsymbol R\in\mathbb R^{3\times3}\mid
  \boldsymbol R^{\mathrm T}\boldsymbol R=\boldsymbol I,
  \det(\boldsymbol R)=1\right\}.
  \label{eq:ch4_so3_definition}
\end{equation}
旋转矩阵虽然包含九个元素，但正交和单位长度约束使其只有三个独立自由度。也可以将其理解
为“一个轴的空间方向”所需的两个自由度，加上“绕该轴旋转”所需的一个自由度。

“安装标定”与“在线估计”是前后相接的两个阶段。安装标定只在重新佩戴后执行，利用已知
参考姿势求出坐标对齐、局部肢段轴、位置模型和角度修正参数；在线估计则在每一个同步采样
时刻使用这些固定参数，把未知动作转换为五自由度状态，不再重新拟合参数。完成安装标定后，
在线估计实际使用两类几何信息：大臂单位方向
${}^{B}\hat{\boldsymbol u}\in S^2$，以及经过安装对齐的前臂完整姿态
${}^{B}\boldsymbol R_F\in SO(3)$。因此可用独立信息量为
\begin{equation}
  \underbrace{2}_{\text{upper-arm direction}}+
  \underbrace{3}_{\text{forearm attitude}}=5.
  \label{eq:ch4_measurement_dof_count}
\end{equation}
前臂姿态的三个自由度还可以分解为前臂方向 ${}^{B}\hat{\boldsymbol f}$ 的两个自由度和
绕前臂纵轴旋转的一个自由度。由此形成表~\ref{tab:ch4_measurement_to_joint} 所示的逐级
求解关系。

\begin{table}[htbp]
  \centering
  \caption{测量几何量与五自由度状态的对应关系}
  \label{tab:ch4_measurement_to_joint}
  \begin{tabularx}{\linewidth}{@{}lXX@{}}
    \toprule
    待估计量 & 主要几何依据 & 物理含义 \\
    \midrule
    $q_p,q_r$ & 大臂单位方向 ${}^{B}\hat{\boldsymbol u}$ & 大臂指向的两个自由度 \\
    $q_e$ & ${}^{B}\hat{\boldsymbol u}$ 与 ${}^{B}\hat{\boldsymbol f}$ 的夹角
      & 肘部由弯曲向伸直的变化量 \\
    $q_y$ & 前臂方向在大臂法平面内的投影 & 肘平面绕大臂纵轴的朝向 \\
    $q_w$ & 去除肩、肘和固定安装旋转后的前臂残余姿态
      & 绕前臂纵轴的旋前/旋后 \\
    \bottomrule
  \end{tabularx}
\end{table}

上述自由度计数只说明一般姿态下信息量足够，并不保证所有姿态都存在唯一解。当大臂正侧举
时，pitch 方向退化；当肘部完全伸直时，肩部轴向旋转与前臂旋转发生耦合。换言之，该测量
映射在一般姿态下局部可辨识，但在特定几何构型下秩会下降。系统对这些区域的处理将在
第~\ref{sec:ch4_observability_robustness} 节给出。

\section{双接收器同步采集与数据预处理}
\label{sec:ch4_synchronized_acquisition}

\subsection{同步位姿对}
\label{subsec:ch4_synchronized_pose_pairs}

两个 WEMT 接收器分别通过网络返回独立的数据流。如果直接使用两条消息各自最近的样本，
快速运动时可能把不同时刻的大臂和前臂姿态组合在一起，从而产生虚假的肘角和肩部旋转。
因此，采集节点只在以下条件同时满足时发布一对位姿：
\begin{enumerate}[label=\arabic*)]
  \item 两个接收器均处于连接状态且没有设备错误；
  \item 两个接收器都产生了新的硬件帧；
  \item 任一位姿的接收时间不超过 $100\,\mathrm{ms}$；
  \item 两个接收器的到达时间差不超过 $30\,\mathrm{ms}$；
  \item 两个原始四元数均通过有效性检查。
\end{enumerate}

通过检查后，系统给两个 \texttt{geometry\_msgs/PoseStamped} 消息赋予同一个 ROS 时间戳。
该消息由三维位置、四元数姿态、参考坐标系标识和时间戳组成。两条消息分别发布至 ROS 话题
\texttt{/em/upper\_arm\_pose} 和 \texttt{/em/forearm\_pose}。这样，后续标定和在线估计
可按相同时间戳进行严格配对。WEMT 接口内部的平滑窗口被设为 1，即发布未经时间平均的
单帧样本；这里的 1 是窗口中包含的样本数，而不是 1 秒。该设置避免位置和姿态使用不同
滤波造成额外的相位差；统一滤波在完成双流配对后进行。

\subsection{四元数有效性与旋转矩阵}
\label{subsec:ch4_quaternion_validation}

设设备返回四元数为
\begin{equation}
  \boldsymbol q=\begin{bmatrix}q_w&q_x&q_y&q_z\end{bmatrix}^{\mathrm T}.
  \label{eq:ch4_quaternion}
\end{equation}
系统首先检查其模长是否满足
\begin{equation}
  0.95\leq \lVert\boldsymbol q\rVert_2\leq1.05.
  \label{eq:ch4_quaternion_norm_gate}
\end{equation}

未通过检查的数据对被直接丢弃，而不是先强制归一化后继续使用。这是因为持续出现约
0.707 等异常模长时，问题通常来自设备原始数据失效；单纯归一化会掩盖错误并把无效姿态
传入控制链。有效四元数归一化后转换为旋转矩阵：
\begin{equation}
\boldsymbol R(\boldsymbol q)=
\begin{bmatrix}
1-2(q_y^2+q_z^2) & 2(q_xq_y-q_zq_w) & 2(q_xq_z+q_yq_w)\\
2(q_xq_y+q_zq_w) & 1-2(q_x^2+q_z^2) & 2(q_yq_z-q_xq_w)\\
2(q_xq_z-q_yq_w) & 2(q_yq_z+q_xq_w) & 1-2(q_x^2+q_y^2)
\end{bmatrix}.
\label{eq:ch4_quaternion_to_rotation}
\end{equation}

由于当前设备接口只明确四元数分量顺序，而接收器到发射器的主动/被动旋转方向仍需由系统
侧验证，系统在标定时
分别评价原旋转矩阵 $\boldsymbol R$ 和转置矩阵 $\boldsymbol R^{\mathrm T}$。选择标准以前臂
旋转轴在多参考姿势下的方向残差为主，以相对旋转残差作为分数接近时的辅助判据。最终选定
的约定被写入当次标定文件，并在在线估计中保持一致。这一步不是任意改变四元数分量，而是
在互为逆旋转的两种物理解释之间进行一次安装级一致性检验；若后续由设备协议明确方向，仍可
保留该检验作为接线或固件变化的诊断。

\section{安装级自适应标定}
\label{sec:ch4_installation_calibration}

\subsection{标定流程概览}
\label{subsec:ch4_calibration_overview}

完整标定由三个前置质量/动作检查和七个静态参考姿势组成，如图~\ref{fig:ch4_calibration_flow}
所示。
\begin{figure}[htbp]
  \centering
  \fbox{\parbox{0.88\linewidth}{\centering
    图示占位：静态质量检查 $\rightarrow$ 接收器身份检查 $\rightarrow$
    前臂旋转轴标定 $\rightarrow$ 七姿势采集 $\rightarrow$
    参数拟合 $\rightarrow$ 质量门槛 $\rightarrow$ 保存或拒绝。}}
  \caption{安装级自适应标定流程（待替换为正式绘图）}
  \label{fig:ch4_calibration_flow}
\end{figure}

标定开始前，操作者保持自然站立，系统采集 $4\,\mathrm{s}$ 静态数据，用于计算位置和姿态
噪声。随后，操作者在保持大臂基本不动的情况下反复屈伸肘部并转动前臂。若前臂接收器的
姿态变化明显大于大臂接收器，才认为两个接收器的身份与安装位置正确。最后，操作者保持大臂
竖直、肘部弯曲约 $90^\circ$，仅使前臂绕“肘到腕”的纵轴做旋前/旋后，用于辨识前臂接收器
局部坐标系中的旋转轴。

七个静态参考姿势及其理论角度见表~\ref{tab:ch4_reference_poses}。每个姿势离开后重新构建
3 次，每次保持 $2\,\mathrm{s}$。这种“离开---重建---保持”的方式能够发现单次摆姿偶然
正确但重复性不足的问题。

\begin{table}[htbp]
  \centering
  \caption{安装级标定使用的静态参考姿势}
  \label{tab:ch4_reference_poses}
  \begin{tabularx}{\linewidth}{@{}XrrrrX@{}}
    \toprule
    参考姿势 & $q_p$ & $q_r$ & $q_y$ & $q_e$ & 主要作用 \\
    \midrule
    DOWN             & $0^\circ$   & $0^\circ$  & $0^\circ$ & $0^\circ$  & 定义系统零位 \\
    FORWARD\_30      & $-30^\circ$ & $0^\circ$  & $0^\circ$ & $0^\circ$  & 约束大臂位置模型 \\
    FORWARD\_60      & $-60^\circ$ & $0^\circ$  & $0^\circ$ & $0^\circ$  & 约束大臂位置模型 \\
    FORWARD          & $-90^\circ$ & $0^\circ$  & $0^\circ$ & $0^\circ$  & 约束前举方向 \\
    SIDEWAYS         & $0^\circ$   & $90^\circ$ & $0^\circ$ & $90^\circ$ & 约束侧举和伸肘方向 \\
    STRAIGHT\_DOWN   & $0^\circ$   & $0^\circ$  & $0^\circ$ & $90^\circ$ & 约束伸肘姿态 \\
    STRAIGHT\_FORWARD& $-90^\circ$ & $0^\circ$  & $0^\circ$ & $90^\circ$ & 约束主要工作空间 \\
    \bottomrule
  \end{tabularx}
\end{table}

图~\ref{fig:ch4_reference_pose_set} 给出七个参考姿势应呈现的肢段关系。FORWARD\_30 和
FORWARD\_60 的角度不要求操作者用量角器达到绝对精确值；其作用是为前举路径增加两个
中间方向。操作者应优先保证躯干不动、动作平面一致、肘部状态正确以及三次重复可重现，实际
偏差将反映在参考残差和位置模型残差中。

\begin{figure}[htbp]
  \centering
  \fbox{\parbox{0.92\linewidth}{\centering
    示意图内容：依次绘制 DOWN、FORWARD\_30、FORWARD\_60、FORWARD、SIDEWAYS、
    STRAIGHT\_DOWN 和 STRAIGHT\_FORWARD；用实线表示大臂、虚线表示前臂，并标注
    “弯肘约 $90^\circ$”或“伸肘”以及前举/侧举方向。}}
  \caption{七个安装标定参考姿势及其肢段约束（待替换为实拍图或正式绘图）}
  \label{fig:ch4_reference_pose_set}
\end{figure}

FORWARD\_30 和 FORWARD\_60 只用于提高大臂位置模型在前举路径上的可辨识性；前臂方向
拟合使用 DOWN、FORWARD、SIDEWAYS、STRAIGHT\_DOWN 和 STRAIGHT\_FORWARD 五类姿势。

\subsection{静态质量和重复性评价}
\label{subsec:ch4_static_quality_repeatability}

对于一组静态位置样本 $\boldsymbol p_i$，位置均方根波动定义为
\begin{equation}
  e_{p,\mathrm{RMS}}=
  \sqrt{\frac{1}{N}\sum_{i=1}^{N}
  \lVert\boldsymbol p_i-\bar{\boldsymbol p}\rVert_2^2}.
  \label{eq:ch4_static_position_rms}
\end{equation}

姿态中心采用旋转矩阵均值在 $SO(3)$ 上的最近旋转，姿态波动则用每个样本到中心旋转的
测地角计算 RMS。由于旋转不能像普通欧氏向量那样逐元素平均后直接使用，最近旋转投影和
四元数/旋转平均是姿态处理中常用的几何做法~\cite{hartley2013rotation,markley2007averaging}。
当前工程实现中，当任一接收器的位置 RMS 超过 $3\,\mathrm{mm}$，或姿态
RMS 超过 $3^\circ$ 时，系统提示检查绑带、金属物体和跟踪距离。

每个参考姿势的三次重复使用旋转 medoid（样本中心）作为代表，即只在三个实际观测中选择
一个，使它到其余观测的旋转测地距离总和最小。与直接平均相比，medoid 必然是一个真实采集
样本，对单次离群姿态更稳健。若任一重复与代表样本之间的最大角距离超过 $15^\circ$，则只重新采集当前
姿势，已经通过的姿势继续保留。该策略避免一次摆姿错误使整轮标定全部作废，同时保证每个
参考姿势具有最低重复性。

\subsection{前臂局部旋转轴辨识}
\label{subsec:ch4_forearm_twist_axis}

在理想旋前/旋后动作中，前臂在发射器坐标系中的纵轴方向保持近似不变，而接收器绕该轴
转动。实际只有一个前臂接收器，下标 $i$ 表示该接收器在第 $i$ 个采样时刻的观测。设其旋转
矩阵为 ${}^{T}\boldsymbol R_{F,i}$，希望找到接收器外壳坐标系中的局部单位向量
$\boldsymbol a_t$，使
\begin{equation}
  {}^{T}\boldsymbol R_{F,i}\boldsymbol a_t\approx\boldsymbol d_t,
  \qquad i=1,\ldots,N,
  \label{eq:ch4_twist_axis_invariance}
\end{equation}
其中，$\boldsymbol d_t$ 为发射器坐标系中的稳定旋转轴。寻找 $\boldsymbol a_t$ 的目的，是把
任意安装朝向下的“接收器外壳某一轴”辨识为人体前臂纵轴；若轴选择正确，接收器即使绕它
转动，该轴变换到发射器坐标系后的方向仍应近似不变。令
\begin{equation}
  \bar{\boldsymbol R}_F=\frac{1}{N}
  \sum_{i=1}^{N}{}^{T}\boldsymbol R_{F,i},
  \label{eq:ch4_mean_forearm_rotation}
\end{equation}
则 $\boldsymbol a_t$ 可写为
\begin{equation}
  \boldsymbol a_t=
  \arg\max_{\lVert\boldsymbol a\rVert=1}
  \lVert\bar{\boldsymbol R}_F\boldsymbol a\rVert_2,
  \label{eq:ch4_twist_axis_svd}
\end{equation}
即 $\bar{\boldsymbol R}_F$ 最大奇异值对应的右奇异向量。式
~\eqref{eq:ch4_twist_axis_invariance} 给出物理不变性，式
~\eqref{eq:ch4_mean_forearm_rotation} 将全部姿态观测聚合，式
~\eqref{eq:ch4_twist_axis_svd} 则寻找经这些旋转后最不发散的局部方向；最大奇异值对应方向
正是保持合向量长度最大的方向。系统对残差进行基于中位数绝对
偏差的迭代剔除，降低偶发姿态跳变的影响；同时通过样本对之间的相对旋转检查动作范围。
只有当前臂旋转范围至少达到 $55^\circ$、大臂姿态变化不超过 $15^\circ$，且旋转轴拟合
RMS 不超过 $15^\circ$ 时，才继续执行后续标定。

该旋转轴专门服务于腕部旋前/旋后估计。它不与前臂“肘到腕”的方向轴强制共用，因为皮肤
滑动、绑带安装和参考动作误差会使两者的最优拟合结果存在差异。实际实现分别辨识前臂方向
轴与旋转轴，以避免一个轴同时承担两个不同物理任务造成系统性偏差。

\subsection{大臂方向的仿射位置--方向模型}
\label{subsec:ch4_upper_arm_affine_model}

大臂接收器固定在软组织表面，其姿态容易受到皮肤相对肱骨转动的影响。实验开发过程中发现，
直接使用大臂接收器姿态估计肱骨轴向运动会引入较大误差。因此，当前混合模型主要利用大臂
接收器的位置恢复大臂方向。这里的目标不是根据位置直接回归关节角，而是先恢复具有明确
几何意义的单位方向 ${}^{B}\hat{\boldsymbol u}$，再由第~\ref{sec:ch4_online_arm_estimation}
节的逆运动学关系求解关节变量。

若接收器与肩关节构成理想刚体，大臂运动时接收器应在以肩关节为中心的球面上运动。但实际
轨迹受绑带、软组织和发射器安装影响，前举路径往往呈现非球形畸变。为此，本研究采用仿射
位置--方向模型
\begin{equation}
  {}^{T}\boldsymbol p_{U,i}
  =\boldsymbol c+\boldsymbol M\,{}^{B}\boldsymbol u_i^{*}
  +\boldsymbol\varepsilon_i,
  \label{eq:ch4_upper_affine_model}
\end{equation}
其中，$\boldsymbol c\in\mathbb R^3$ 为本次安装下的轨迹中心，
$\boldsymbol M\in\mathbb R^{3\times3}$ 为从人体标准方向到发射器坐标位移的可逆仿射
矩阵，${}^{B}\boldsymbol u_i^{*}$ 是第 $i$ 个参考姿势的理论大臂方向，
$\boldsymbol\varepsilon_i$ 汇集摆姿误差、软组织运动和测量噪声。$\boldsymbol c$ 是用于
位置--方向映射的虚拟中心，不应直接解释为解剖学肩关节中心。

这里称为“混合模型”，是因为大臂方向主要由接收器\emph{位置}恢复，而前臂方向和轴向
旋转由接收器\emph{姿态}恢复。所谓“恢复大臂方向”，是指由测得的位置
${}^{T}\boldsymbol p_U$ 反求单位向量 ${}^{B}\hat{\boldsymbol u}$，而不是估计大臂长度或
肩关节中心。图~\ref{fig:ch4_affine_geometry} 给出式~\eqref{eq:ch4_upper_affine_model} 的
几何含义。

\begin{figure}[htbp]
  \centering
  \fbox{\parbox{0.92\linewidth}{\centering
    示意图内容：以 $\boldsymbol c$ 为椭球中心，单位方向
    ${}^{B}\boldsymbol u_i^{*}$ 经过 $\boldsymbol M$ 映射为从 $\boldsymbol c$ 指向
    ${}^{T}\boldsymbol p_{U,i}$ 的位移；标出实测点、重投影点 $\hat{\boldsymbol p}_U$、
    残差 $\boldsymbol\varepsilon_i$ 与模型误差 $e_M$。左侧可并列绘制单位球，右侧绘制其
    经过 $\boldsymbol M$ 后的椭球像。}}
  \caption{大臂仿射位置--方向模型的几何解释（待替换为正式绘图）}
  \label{fig:ch4_affine_geometry}
\end{figure}

该模型被称为“仿射椭球”，是因为单位方向满足
$({}^{B}\boldsymbol u)^{\mathrm T}{}^{B}\boldsymbol u=1$。忽略噪声并令
$\boldsymbol x={}^{T}\boldsymbol p_U-\boldsymbol c$，由
$\boldsymbol x=\boldsymbol M{}^{B}\boldsymbol u$ 可得
\begin{equation}
  \boldsymbol x^{\mathrm T}\boldsymbol M^{-\mathrm T}
  \boldsymbol M^{-1}\boldsymbol x=1,
  \label{eq:ch4_affine_ellipsoid_implicit}
\end{equation}
这里引入 $\boldsymbol x$ 只是为了简化表达：它是从虚拟中心 $\boldsymbol c$ 指向当前接收器
位置的位移，不是新的测量量。将 ${}^{B}\boldsymbol u=\boldsymbol M^{-1}\boldsymbol x$
代回单位长度约束即可得到式~\eqref{eq:ch4_affine_ellipsoid_implicit}，这正是以
$\boldsymbol c$ 为中心的椭球面二次型。若对矩阵作奇异值分解
\begin{equation}
  \boldsymbol M=\boldsymbol U\boldsymbol\Sigma\boldsymbol V^{\mathrm T},
  \qquad
  \boldsymbol\Sigma=\operatorname{diag}(\sigma_1,\sigma_2,\sigma_3),
  \label{eq:ch4_affine_matrix_svd}
\end{equation}
则 $\boldsymbol U$ 给出椭球主轴方向，$\sigma_1,\sigma_2,\sigma_3$ 给出三个半轴长度。
当三个奇异值相等时，说明三个主方向的尺度相同；若
$\boldsymbol M=r\boldsymbol R$，线性变换只包含统一尺度 $r$ 和旋转 $\boldsymbol R$，模型便
退化为半径为 $r$ 的球面。若奇异值不同，单位球在三个主方向上被拉伸为不同半轴长度，形成
椭球。因此，椭球是单位方向集合经过本次安装相关线性变换后的位置轨迹，而不是对人体肩关节形状的
解剖学假设。实际七个参考姿势只覆盖该椭球面的有限区域，本模型应理解为工作空间内的经验
标定模型，而非完整肩部轨迹的全局物理模型。

所有参考样本通过最小二乘求解
\begin{equation}
  (\hat{\boldsymbol c},\hat{\boldsymbol M})=
  \arg\min_{\boldsymbol c,\boldsymbol M}
  \sum_i\left\|
  {}^{T}\boldsymbol p_{U,i}-\boldsymbol c-
  \boldsymbol M{}^{B}\boldsymbol u_i^{*}
  \right\|_2^2.
  \label{eq:ch4_upper_affine_fit}
\end{equation}
式~\eqref{eq:ch4_upper_affine_fit} 对 $\boldsymbol c$ 和 $\boldsymbol M$ 是线性最小二乘
问题。它求出的 $\hat{\boldsymbol c}$ 和 $\hat{\boldsymbol M}$ 用于在线阶段把新的接收器
位置反算为大臂方向。为了把常量中心项和三维方向项写在同一个线性方程中，将每个理论方向
在末尾附加常数 1，形成增广向量
$[({}^{B}\boldsymbol u_i^{*})^{\mathrm T},1]$。相应设计矩阵必须具有满列秩 4，即四个列方向
都具有独立信息；否则轨迹
中心和全部仿射分量不能被唯一辨识。DOWN、FORWARD 和 SIDEWAYS 三类方向不足以满足这一
条件，因此标定流程额外引入 FORWARD\_30 和 FORWARD\_60，以增加前举路径上的方向变化
并改善模型条件。

在线阶段，根据当前大臂接收器位置计算未归一化方向
\begin{equation}
  \tilde{\boldsymbol u}=
  \hat{\boldsymbol M}^{-1}
  ({}^{T}\boldsymbol p_U-\hat{\boldsymbol c}),
  \label{eq:ch4_upper_direction_raw}
\end{equation}
并归一化得到
\begin{equation}
  {}^{B}\hat{\boldsymbol u}=
  \frac{\tilde{\boldsymbol u}}
  {\lVert\tilde{\boldsymbol u}\rVert_2}.
  \label{eq:ch4_upper_direction_normalized}
\end{equation}
若当前点严格位于标定椭球面上，则
$\lVert\tilde{\boldsymbol u}\rVert_2=1$；实际数据受噪声影响通常偏离曲面，因此仍需归一化。
这一过程可以理解为先用逆仿射变换消除安装相关的方向性畸变，再将结果投影回单位方向空间。

矩阵 $\hat{\boldsymbol M}$ 的奇异值表示沿三个主方向的位移尺度。最小奇异值过小时，说明
某一方向几乎没有被参考姿势激励，求逆时该方向的微小位置噪声会被显著放大。为此定义条件数
\begin{equation}
  \kappa(\hat{\boldsymbol M})=
  \frac{\sigma_{\max}(\hat{\boldsymbol M})}
       {\sigma_{\min}(\hat{\boldsymbol M})}.
  \label{eq:ch4_affine_condition_number}
\end{equation}
条件数越接近 1，三个方向的尺度越均衡；条件数过大则表明参考方向分布不足、摆姿不一致或
拟合接近奇异。当前标定质量门槛要求 $\kappa(\hat{\boldsymbol M})\leq4.0$。

为了检测设备移动或当前动作严重超出标定范围，系统将估计方向重新投影到椭球表面：
\begin{equation}
  \hat{\boldsymbol p}_U=
  \hat{\boldsymbol c}+
  \hat{\boldsymbol M}{}^{B}\hat{\boldsymbol u},
  \label{eq:ch4_upper_position_reprojection}
\end{equation}
并定义位置模型误差
\begin{equation}
  e_M=\left\|{}^{T}\boldsymbol p_U-
  \hat{\boldsymbol p}_U\right\|_2.
  \label{eq:ch4_upper_model_error}
\end{equation}
在线运行时，若 $e_M>50\,\mathrm{mm}$，当前帧被保护性丢弃。该误差不等同于人体大臂
方向的绝对真值误差，而是接收器位置与本次安装模型的一致性指标。小的 $e_M$ 只能说明
当前数据与标定轨迹相容；要评价大臂方向的绝对精度，仍需要独立的光学动作捕捉或其他真值
测量系统。

\subsection{前臂方向与坐标对齐}
\label{subsec:ch4_forearm_direction_alignment}

设 $\boldsymbol a_d$ 为前臂接收器局部坐标系中的“肘到腕”方向，$\boldsymbol A_d$ 为发射器
坐标系到人体标准坐标系的旋转。对于第 $i$ 个参考姿势，预测方向为
\begin{equation}
  {}^{B}\hat{\boldsymbol f}_i=
  \boldsymbol A_d\,{}^{T}\boldsymbol R_{F,i}\boldsymbol a_d.
  \label{eq:ch4_forearm_direction_prediction}
\end{equation}

式~\eqref{eq:ch4_forearm_direction_prediction} 表示两步坐标变换：先用
${}^{T}\boldsymbol R_{F,i}$ 把接收器内部的候选方向 $\boldsymbol a_d$ 旋转到发射器坐标系，
再用 $\boldsymbol A_d$ 将其统一表达在人体坐标系。方向轴和对齐矩阵通过交替优化求解：固定
$\boldsymbol a_d$ 时，$\boldsymbol A_d$ 可归结为
Wahba 旋转对齐问题，即寻找一个旋转，使多组测量方向与对应理论方向的平方误差之和最小
~\cite{wahba1965least}；固定 $\boldsymbol A_d$ 时，可将理论方向反变换到各接收器局部坐标系
并取均值，更新 $\boldsymbol a_d$。算法从六个坐标轴正负方向初始化，保留参考姿势方向 RMS
最小的解。

旋转对齐子问题为
\begin{equation}
  \hat{\boldsymbol A}_d=
  \arg\min_{\boldsymbol A\in SO(3)}
  \sum_i\left\|
  \boldsymbol A\,{}^{T}\boldsymbol R_{F,i}\boldsymbol a_d-
  {}^{B}\boldsymbol f_i^{*}
  \right\|_2^2,
  \label{eq:ch4_forearm_wahba_problem}
\end{equation}
通过互协方差矩阵的奇异值分解求得，并显式约束 $\det(\boldsymbol A)=1$。在线阶段的前臂
单位方向为
\begin{equation}
  {}^{B}\hat{\boldsymbol f}=
  \frac{
    \hat{\boldsymbol A}_d\,{}^{T}\boldsymbol R_F\hat{\boldsymbol a}_d
  }{
    \left\|
      \hat{\boldsymbol A}_d\,{}^{T}\boldsymbol R_F\hat{\boldsymbol a}_d
    \right\|_2
  }.
  \label{eq:ch4_forearm_direction_online}
\end{equation}

除了方向映射，系统还建立独立的姿态对齐矩阵 $\boldsymbol A_t$ 和固定安装旋转
$\boldsymbol R_m$，用于计算 wrist roll。这里的“姿态对齐”是把接收器报告的完整三维旋转
转换到人体模型坐标系；$\boldsymbol R_m$ 则记录接收器外壳相对理论前臂坐标系的固定安装角。
在 DOWN 姿势下，理论前臂姿态作为腕部零位；三次 DOWN 重复的平均旋转确定本次安装的固定
旋转。在线计算时先去除肩、肘造成的理论旋转，再去除 $\boldsymbol R_m$，剩余的绕前臂纵轴
分量才被解释为腕部扭转。由此，外壳安装角和前臂方向误差不会直接被误认为 wrist roll。

\subsection{肘角的安装级仿射修正}
\label{subsec:ch4_elbow_affine_correction}

由两个单位方向可以得到几何肘角
\begin{equation}
  q_{e,g}=\arcsin\left(
  \operatorname{clip}
  ( {}^{B}\hat{\boldsymbol u}^{\mathrm T}
  {}^{B}\hat{\boldsymbol f},-1,1)
  \right).
  \label{eq:ch4_geometric_elbow_angle}
\end{equation}

下标 $g$ 表示 geometric，即由两条肢段方向直接得到的几何角。在本研究的零位约定中，
弯肘约 $90^\circ$ 时两个方向近似正交，因而 $q_{e,g}\approx0$；
手臂伸直时二者近似平行，因而 $q_{e,g}\approx90^\circ$。由于两个方向拟合各自存在固定
残差，几何角可能同时存在零位偏差和幅值缩放误差：固定零偏只能整体平移角度，不能改变从
弯曲端到伸直端的量程；单独改变增益又不能校正零位。因此本研究对每次安装分别拟合一维
仿射修正，即“缩放加平移”模型
\begin{equation}
  q_e=\operatorname{clip}
  (s_e q_{e,g}+b_e,\,0,\,\pi/2),
  \label{eq:ch4_elbow_affine_correction}
\end{equation}
其中增益 $s_e$ 和偏置 $b_e$ 由五类前臂参考姿势的理论肘角和几何肘角进行线性最小二乘
拟合。为拒绝异常标定，当前实现要求 $0.75\leq s_e\leq1.50$，且
$|b_e|\leq20^\circ$。若拟合不可辨识或超出合理范围，系统保留未经修正的几何肘角，同时
输出警告。

\subsection{标定质量门槛与实机资格}
\label{subsec:ch4_calibration_quality_gates}

标定参数并非只要能够求解就可进入真机。系统在保存前检查表~\ref{tab:ch4_quality_gates}
所列工程门槛。
\begin{table}[htbp]
  \centering
  \caption{安装级标定质量门槛}
  \label{tab:ch4_quality_gates}
  \begin{tabularx}{\linewidth}{@{}XllX@{}}
    \toprule
    检查项 & 建议/诊断阈值 & 硬限制 & 不通过时的处理 \\
    \midrule
    参考姿势重复差异 & --- & $15^\circ$ & 重采当前姿势 \\
    大臂方向参考误差 & $10^\circ$ & $15^\circ$ & 警告/拒绝保存 \\
    前臂方向参考误差 & $10^\circ$；真机资格 $20^\circ$ & $25^\circ$ & 警告、仅诊断或拒绝 \\
    前臂腕部零位重复误差 & --- & $15^\circ$ & 拒绝保存 \\
    仿射矩阵条件数 & 越接近 1 越好 & 4.0 & 拒绝保存 \\
    仿射位置模型 RMS/最大误差 & $30/50\,\mathrm{mm}$ & $35/65\,\mathrm{mm}$ & 警告/拒绝保存 \\
    \bottomrule
  \end{tabularx}
\end{table}

表中的“建议/诊断阈值”用于提示边缘质量，“硬限制”决定标定文件能否保存，两者不应简单
理解为“仿真门槛”和“实机门槛”。仿真与 dry-run 仍可使用通过硬限制但质量边缘的标定；
真机资格另由更严格的前臂方向条件控制。若前臂方向误差超过 $20^\circ$ 但仍未超过硬限制
$25^\circ$，标定可以作为诊断资料保存，
但会被标记为 \texttt{dry\_run\_only: true}。该标定只允许用于 dry-run 和离线诊断，
重定向节点拒绝向真机主题发布关节目标。该设计把“算法能够产生数值”和“标定具备实机
资格”明确分开。

这些阈值是本项目的工程安全标准，并非电磁跟踪器或人形机器人领域的通用精度标准。参考
姿势残差描述标定模型与人工构建姿势之间的一致性，不能替代独立视觉系统测得的人体姿态
真值误差。

\section{在线人体手臂姿态估计}
\label{sec:ch4_online_arm_estimation}

安装标定完成后，每一对同步位姿首先被转换为三个中间几何量：大臂单位方向
${}^{B}\hat{\boldsymbol u}$、前臂单位方向 ${}^{B}\hat{\boldsymbol f}$，以及经过腕部姿态
对齐的前臂接收器旋转 ${}^{B}\boldsymbol R_F$。其中，大臂方向由
式~\eqref{eq:ch4_upper_direction_raw}--\eqref{eq:ch4_upper_direction_normalized} 对当前大臂
接收器位置进行逆仿射映射得到；前臂方向由式~\eqref{eq:ch4_forearm_direction_online} 对当前
前臂接收器姿态和已标定局部轴进行变换得到；“腕部姿态对齐”则指用 $\boldsymbol A_t$ 将完整
接收器旋转转换到人体模型坐标系。五个关节量并不是同时进行黑箱优化，而是
按照几何依赖关系逐级求解：先由大臂方向确定 $q_p,q_r$，再由两条肢段方向确定几何肘角
$q_{e,g}$ 和肩部轴向角 $q_y$，最后从前臂残余旋转中提取 $q_w$。几何肘角随后经过本次
安装的增益和偏置修正，形成最终输出 $q_e$。

这种求解顺序保留了每个变量的物理意义，也使不可观测条件能够被显式检测。它与输入信息的
对应关系可概括为
\begin{equation}
  ({}^{T}\boldsymbol p_U,{}^{T}\boldsymbol R_F)
  \longrightarrow
  ({}^{B}\hat{\boldsymbol u},{}^{B}\hat{\boldsymbol f},
   {}^{B}\boldsymbol R_F)
  \longrightarrow
  (q_p,q_r,q_y,q_{e,g},q_w)
  \longrightarrow
  (q_p,q_r,q_y,q_e,q_w).
  \label{eq:ch4_online_estimation_chain}
\end{equation}

\subsection{肩部 pitch 与 roll}
\label{subsec:ch4_shoulder_pitch_roll}

得到大臂单位方向
\begin{equation}
  {}^{B}\hat{\boldsymbol u}=
  \begin{bmatrix}u_x&u_y&u_z\end{bmatrix}^{\mathrm T}
  \label{eq:ch4_upper_direction_components}
\end{equation}
后，由于 $\boldsymbol R_z(q_y)$ 绕零位大臂轴旋转，不改变
$\boldsymbol a_{u,0}=[0,0,-1]^{\mathrm T}$，大臂方向只与 $q_p,q_r$ 有关。展开
式~\eqref{eq:ch4_segment_directions} 可得
\begin{equation}
  {}^{B}\hat{\boldsymbol u}=
  \begin{bmatrix}
    -\sin q_p\cos q_r\\
    \sin q_r\\
    -\cos q_p\cos q_r
  \end{bmatrix}.
  \label{eq:ch4_upper_direction_expanded}
\end{equation}
这里的“零位大臂轴”就是式~\eqref{eq:ch4_zero_segment_axes} 定义的
$\boldsymbol a_{u,0}$。由于该向量位于 $z$ 轴上，绕 $z$ 轴旋转不会改变它，即
$\boldsymbol R_z(q_y)\boldsymbol a_{u,0}=\boldsymbol a_{u,0}$；随后依次左乘
$\boldsymbol R_x(q_r)$ 和 $\boldsymbol R_y(q_p)$，直接进行矩阵乘法即可得到
式~\eqref{eq:ch4_upper_direction_expanded}。

因此，肩部 pitch 和 roll 的主值可写为
\begin{equation}
  q_p=\operatorname{atan2}(-u_x,-u_z),
  \qquad
  q_r=\arcsin(\operatorname{clip}(u_y,-1,1)).
  \label{eq:ch4_pitch_roll_inverse}
\end{equation}

“主值”表示反三角函数在约定区间内返回的代表解，因为同一方向可对应相差 $2\pi$ 的周期
等价角。由式~\eqref{eq:ch4_upper_direction_expanded} 有
$u_y=\sin q_r$，故 $q_r$ 可由反正弦得到；而 $-u_x$ 与 $-u_z$ 分别包含
$\sin q_p\cos q_r$ 和 $\cos q_p\cos q_r$，用双参数反正切可同时保留象限信息并消去共同
因子 $\cos q_r$。$\operatorname{clip}(u_y,-1,1)$ 只用于消除浮点误差造成的轻微越界，
例如数值上出现 $1+10^{-12}$ 时避免 $\arcsin$ 失效，并非裁剪正常人体动作。

由于 \texttt{atan2} 的输出位于 $[-\pi,\pi]$，系统将当前结果展开到最接近上一帧的等价角，
避免跨越 $\pm\pi$ 时产生整周跳变。

当大臂接近正侧方时，$u_x$ 和 $u_z$ 同时接近零，pitch 对测量噪声高度敏感。定义
\begin{equation}
  o_p=\sqrt{u_x^2+u_z^2}=|\cos q_r|
  \label{eq:ch4_pitch_observability}
\end{equation}
为 pitch 可观测度。在线算法采用带滞回的两状态机制：当 $o_p\leq0.50$ 时进入侧举奇异区，
将 pitch 切换到统一的标准解 $q_p=0$；只有当手臂返回水平面以下且 $o_p\geq0.65$ 时才
重新采用测量解。与此同时，roll 使用标准 pitch 坐标系下的 \texttt{atan2} 表达，使其
能够连续通过 $90^\circ$，而不是被 \texttt{asin} 折回到主值范围。该处理保证从 DOWN
直接侧举和从 FORWARD 转向侧举时，同一最终人体姿势映射到同一机器人目标。

由式~\eqref{eq:ch4_upper_direction_expanded} 可见，$u_x$ 和 $u_z$ 都含有因子
$\cos q_r$；当侧举接近 $|q_r|=90^\circ$ 时该因子趋近零，pitch 的分子和分母同时退化。
阈值 0.50 和 0.65 分别约对应 $|q_r|=60^\circ$ 和 $49^\circ$，是当前工程为避免临界点
抖动设置的进入/退出滞回阈值，而不是人体生理常数。进入奇异区后固定 $q_p=0$，相当于从
无穷多个等价 pitch--roll 表示中选择统一标准解；退出阈值更高且要求回到水平面以下，使状态
不会在噪声作用下频繁切换。

\subsection{肘部伸展量}
\label{subsec:ch4_online_elbow_angle}

在运动学模型中，大臂和前臂方向分别为
${}^{B}\hat{\boldsymbol u}={}^B\boldsymbol R_S\boldsymbol a_{u,0}$ 和
${}^{B}\hat{\boldsymbol f}={}^B\boldsymbol R_S
\boldsymbol R_y(q_{e,g})\boldsymbol a_{f,0}$。利用旋转矩阵保持内积的性质，有
\begin{align}
  ({}^{B}\hat{\boldsymbol u})^{\mathrm T}
  {}^{B}\hat{\boldsymbol f}
  &=\boldsymbol a_{u,0}^{\mathrm T}
    \boldsymbol R_y(q_{e,g})\boldsymbol a_{f,0}
  \nonumber\\
  &=\sin q_{e,g}.
  \label{eq:ch4_elbow_dot_product_derivation}
\end{align}
旋转矩阵满足 $\boldsymbol R^{\mathrm T}\boldsymbol R=\boldsymbol I$，因此对任意向量
$\boldsymbol a,\boldsymbol b$ 都有
$(\boldsymbol R\boldsymbol a)^{\mathrm T}(\boldsymbol R\boldsymbol b)
=\boldsymbol a^{\mathrm T}\boldsymbol b$，即共同旋转不改变夹角或内积。在本式中，肩部旋转
${}^{B}\boldsymbol R_S$ 因而被消去；再代入
$\boldsymbol a_{u,0}=[0,0,-1]^{\mathrm T}$、
$\boldsymbol a_{f,0}=[1,0,0]^{\mathrm T}$ 和 $\boldsymbol R_y(q_{e,g})$ 的矩阵形式，内积
恰好等于 $\sin q_{e,g}$。

因此，式~\eqref{eq:ch4_geometric_elbow_angle} 可以直接由两个肢段方向求得，而不依赖肩部
三个角的具体数值。在本研究的约定下，$q_{e,g}=0$ 对应约 $90^\circ$ 屈肘，
$q_{e,g}=\pi/2$ 对应伸直。为避免与临床文献中常用的“屈曲角”混淆，本文将其称为肘部
伸展量。

这种零位来自当前 G1 左臂模型与重定向接口的关节语义：机器人模型零位附近，大臂向下而
前臂向前，二者约成直角；正向 $q_e$ 将前臂逐渐转到与大臂同向。因此本文的 $q_e$ 是
“从约 $90^\circ$ 弯肘零位出发的伸展量”，不等同于临床上通常从完全伸直定义的肘屈曲角。

几何肘角直接反映当前两条估计方向之间的关系。最终发布的肘部变量使用
式~\eqref{eq:ch4_elbow_affine_correction} 进行安装级增益和零偏修正。这里应区分
$q_{e,g}$ 与 $q_e$：前者用于保持几何分解与测量方向一致，后者是面向机器人目标的校正后
输出。
必须保留两者，是因为 wrist 残余姿态分解要求理论前臂方向与当前实测大臂、前臂方向严格
一致，故使用未修正的 $q_{e,g}$；而输出给重定向层的肘角需要补偿安装造成的尺度和零偏，故
使用 $q_e$。若提前用 $q_e$ 构造理论前臂姿态，标定补偿量会被错误地吸收到 wrist twist 中。

\subsection{肩部 yaw}
\label{subsec:ch4_shoulder_yaw}

大臂方向本身不能确定绕大臂轴的旋转。系统使用前臂在垂直于大臂方向平面内的分量估计
肩部 yaw：
\begin{equation}
  \boldsymbol f_{\perp}=
  {}^{B}\hat{\boldsymbol f}-
  \left({}^{B}\hat{\boldsymbol f}^{\mathrm T}
  {}^{B}\hat{\boldsymbol u}\right)
  {}^{B}\hat{\boldsymbol u}.
  \label{eq:ch4_forearm_perpendicular_component}
\end{equation}

若 $\lVert\boldsymbol f_{\perp}\rVert$ 足够大，将其归一化，并与在当前 pitch、roll 下的零
yaw 参考方向比较。几何上，$\boldsymbol f_{\perp}$ 是前臂方向在大臂法平面内的投影，
它确定了“肩--肘--腕”平面绕大臂纵轴的朝向。设零 yaw 参考方向投影并归一化后为
$\boldsymbol f_{0,\perp}$，则
\begin{equation}
  q_y=\operatorname{atan2}
  \left(
  -{}^{B}\hat{\boldsymbol u}^{\mathrm T}
  (\boldsymbol f_{0,\perp}\times\boldsymbol f_{\perp}),
  \boldsymbol f_{0,\perp}^{\mathrm T}\boldsymbol f_{\perp}
  \right).
  \label{eq:ch4_shoulder_yaw_inverse}
\end{equation}

投影的作用是去除前臂沿大臂轴的分量，只保留能够描述肘平面绕大臂轴朝向的分量。归一化
则去除投影长度，只比较方向；但只有投影长度足够大时，归一化才不会放大噪声。零 yaw 参考
方向 $\boldsymbol f_{0,\perp}$ 由当前 $q_p,q_r$ 和 $q_y=0$ 的运动学模型生成，因此比较二者
即可得到实际肘平面相对零 yaw 肘平面的有符号转角。式~\eqref{eq:ch4_shoulder_yaw_inverse}
中，点积给出余弦项，沿大臂轴的叉积投影给出带符号正弦项，$\operatorname{atan2}$ 将二者
合成为保留象限的有符号角。前置负号用于匹配 G1 肩 yaw 的正方向约定。

\begin{figure}[htbp]
  \centering
  \fbox{\parbox{0.92\linewidth}{\centering
    左图：在垂直于大臂方向 ${}^{B}\hat{\boldsymbol u}$ 的平面内，画出前臂投影
    $\boldsymbol f_{\perp}$、零 yaw 参考 $\boldsymbol f_{0,\perp}$ 和有符号角 $q_y$；
    右图：对比弯肘时投影较长、yaw 与 wrist 可分辨，以及直臂时投影趋近零、两种轴向旋转
    耦合的情形。}}
  \caption{肩部 yaw 的平面投影求解及直臂不可观测性（待替换为正式绘图）}
  \label{fig:ch4_yaw_projection_observability}
\end{figure}

与 pitch 相同，yaw 也展开到最接近上一帧的等价角。由于旋转 $q_y$ 与
$q_y+2k\pi$（$k\in\mathbb Z$）表示同一姿态，展开只选择最连续的数值分支，并不改变物理
旋转。

\subsection{前臂旋前/旋后}
\label{subsec:ch4_forearm_pronation_supination}

当前实现使用与两条测量方向（大臂单位方向 ${}^{B}\hat{\boldsymbol u}$ 和前臂单位方向
${}^{B}\hat{\boldsymbol f}$）直接一致的几何肘角 $q_{e,g}$ 构造不含 wrist roll 的理论
前臂旋转
\begin{equation}
  {}^{B}\boldsymbol R_{E,g}=
  \boldsymbol R_y(q_p)\boldsymbol R_x(q_r)
  \boldsymbol R_z(q_y)\boldsymbol R_y(q_{e,g}).
  \label{eq:ch4_forearm_rotation_without_wrist}
\end{equation}
当前接收器在人体标准坐标系中的旋转为
\begin{equation}
  {}^{B}\boldsymbol R_F=
  \boldsymbol A_t\,{}^{T}\boldsymbol R_F.
  \label{eq:ch4_forearm_sensor_body_rotation}
\end{equation}
去除理论前臂旋转和固定安装旋转后，残余旋转为
\begin{equation}
  \boldsymbol R_{res}=
  ({}^{B}\boldsymbol R_{E,g})^{\mathrm T}
  {}^{B}\boldsymbol R_F\boldsymbol R_m^{\mathrm T}.
  \label{eq:ch4_wrist_residual_rotation}
\end{equation}

上述三式按“预测---观测---求残差”的顺序工作。式
~\eqref{eq:ch4_forearm_rotation_without_wrist} 用已估计的肩、肘量预测前臂在未转掌时的姿态；
式~\eqref{eq:ch4_forearm_sensor_body_rotation} 将传感器观测变换到同一人体坐标系；式
~\eqref{eq:ch4_wrist_residual_rotation} 左乘理论姿态的逆、右乘固定安装旋转的逆，从观测中
去除肩部运动、肘部运动和外壳安装角。理想剩余量只包含前臂绕自身纵轴的旋转。本文沿用
机器人关节名将该量称为 wrist roll，但其人体意义是前臂旋前/旋后，并非手腕屈伸。

等价地，前臂接收器姿态模型可写为
\begin{equation}
  {}^{B}\boldsymbol R_F
  \approx{}^{B}\boldsymbol R_{E,g}
  \boldsymbol R_x(q_w)\boldsymbol R_m.
  \label{eq:ch4_forearm_sensor_attitude_model}
\end{equation}
其中，$\boldsymbol R_m$ 表示本次安装中接收器与前臂模型坐标系之间的固定旋转。

理想情况下，$\boldsymbol R_{res}$ 仅包含绕前臂局部 $+x$ 轴的旋转。为抑制小幅残余 swing
误差，系统将 $\boldsymbol R_{res}$ 转为四元数
$\boldsymbol q_{res}=[q_0,q_x,q_y,q_z]^{\mathrm T}$，只保留向量部分在 $+x$ 轴上的投影，
得到
\begin{equation}
  q_w=\operatorname{wrap}_{[-\pi,\pi)}
  \left(2\operatorname{atan2}(q_x,q_0)\right).
  \label{eq:ch4_wrist_twist_angle}
\end{equation}
然后将结果展开到上一帧附近，以维持时间连续性。这一过程等价于围绕前臂纵轴的
swing--twist 分解：swing 改变纵轴指向，twist 则绕该纵轴旋转；只保留 $q_x$ 相当于提取
局部 $+x$ 轴上的 twist 分量。此类分解常用于把三维关节旋转分离为摆动和轴向扭转
~\cite{murthy2019joint}。结果展开的原因与 yaw 相同：四元数给出的主角位于有限区间，相邻
帧若跨过 $\pm\pi$ 会在数值上跳变约 $2\pi$，选择离上一帧最近的等价角可保持连续。
wrist 提取完成后，系统才将 $q_{e,g}$ 修正为最终输出 $q_e$；因此，
肘角的经验增益和零偏不会被误解释为额外的前臂轴向旋转。

\section{不可观测区域与稳健性处理}
\label{sec:ch4_observability_robustness}

\subsection{直臂状态下 yaw 与 wrist 的耦合}
\label{subsec:ch4_yaw_wrist_unobservable}

当肘部接近完全伸直时，大臂和前臂方向近似共线，
$\lVert\boldsymbol f_{\perp}\rVert\rightarrow0$。此时，仅凭两个肢段方向无法区分“肩部
绕大臂轴旋转”和“前臂自身旋转”；同一物理姿态可以对应多组 $(q_y,q_w)$。这是测量模型
的几何不可观测性，而不是数值求解精度不足。

原因是直臂时肩 yaw 轴与前臂 twist 轴近似重合，末端姿态主要取决于二者的组合量。对任意
小角度 $\delta$，参数对
\begin{equation}
  (q_y,q_w)\sim(q_y+\delta,q_w-\delta)
  \label{eq:ch4_yaw_wrist_equivalent_pairs}
\end{equation}
可产生相同或近似相同的总轴向旋转，因而无法由当前两接收器数据唯一判断该旋转应归属于
肩部还是前臂。弯肘后两条旋转轴不再共线，图
~\ref{fig:ch4_yaw_projection_observability} 所示投影重新出现，两个变量才恢复可分辨性。

系统定义
\begin{equation}
  o_y=\lVert\boldsymbol f_{\perp}\rVert_2
  \label{eq:ch4_yaw_observability}
\end{equation}
为 yaw 可观测度。由于 $\boldsymbol u$ 和 $\boldsymbol f$ 均为单位向量，且
$\boldsymbol u^{\mathrm T}\boldsymbol f=\sin q_{e,g}$，可进一步得到
\begin{equation}
  o_y=\sqrt{1-(\boldsymbol u^{\mathrm T}\boldsymbol f)^2}
  =|\cos q_{e,g}|.
  \label{eq:ch4_yaw_observability_elbow_relation}
\end{equation}
因此，这一退化不是经验上偶然出现的数值问题，而是由五自由度模型的几何结构直接决定。
弯肘约 $90^\circ$ 时 $q_{e,g}\approx0$、$o_y\approx1$，肘平面方向清晰；手臂趋于伸直时
$q_{e,g}\rightarrow\pi/2$、$o_y\rightarrow0$，肩 yaw 轴与前臂 twist 轴趋于重合，测量
中只能保留两者的组合旋转而不能唯一分离。

当 $o_y<0.90$ 时，保持最后一次可靠的 $q_y$ 和 $q_w$，避免把微小方向
噪声放大为几十度的轴向旋转。由于本模型中 $o_y$ 随肘部伸直而下降，该阈值会在完全伸直
之前提前冻结轴向角。系统恢复到明显弯肘姿态后，再继续更新 yaw 和 wrist。

因此，直臂前举和直臂侧举可以用于评价大臂方向及肘角，但不适合
把直臂主动转掌作为 wrist 精度的主要证据。wrist 单关节实验应保持肘部弯曲约 $60^\circ$
至 $90^\circ$，使肩部轴向旋转和前臂旋转保持可分辨。

\subsection{周期角展开与时间连续性}
\label{subsec:ch4_temporal_conditioning}

肩部 yaw、wrist roll 和其他周期角首先展开到最接近上一帧的 $2\pi$ 等价分支。反三角函数
通常只返回 $[-\pi,\pi)$ 内的主值，因此真实角度从 $+\pi$ 连续增加时，下一帧数值会变成
接近 $-\pi$；两帧物理上很接近，数值却相差近 $2\pi$。角展开在
$\theta+2k\pi$ 的等价集合中选择最接近上一帧的分支，从而避免该边界产生整周跳变。侧举
奇异区和直臂不可观测区的内部状态也只在满足带滞回的
重新观测条件后更新，防止临界点附近反复切换。

本章输出的是经过可观测性处理的五自由度人体估计，不在此处施加 G1 关节限位、机器人
命令低通滤波或真机速度限制。上述机器人目标调理统一放在第5章说明，从而把“人体状态
估计”和“机器人命令生成”分开。

\subsection{无效数据与看门狗}
\label{subsec:ch4_invalid_data_watchdog}

在线估计过程中，以下情况不会产生新的关节目标：
\begin{itemize}
  \item 接收器数据陈旧或双流时间差过大；
  \item 原始四元数模长异常；
  \item 两个话题的坐标系标识不一致；
  \item 大臂仿射位置模型误差超过 $50\,\mathrm{mm}$；
  \item 标定矩阵无效或当前方向向量接近零。
\end{itemize}

“看门狗”是监测数据或程序是否在规定时间内持续更新的超时机制。若连续
$0.25\,\mathrm{s}$ 没有新的有效同步数据对，重定向节点发出看门狗警告并停止更新
目标。真机桥还具有独立的 UDP 目标超时和控制权释放机制。因此，采集层、重定向层与执行
层分别检查数据新鲜度，避免单点故障持续驱动机器人。

\section{算法流程与实现}
\label{sec:ch4_algorithm_implementation}

一次完整在线估计可归纳为算法 4-1。

\begin{quote}
\textbf{算法 4-1：双电磁接收器人体左臂姿态估计}

\textbf{输入：}同步接收器位姿
$({}^{T}\boldsymbol p_U,{}^{T}\boldsymbol R_U,
{}^{T}\boldsymbol p_F,{}^{T}\boldsymbol R_F)$，当次安装标定参数，上一帧角度。

\textbf{输出：}五自由度人体手臂估计
$\hat{\boldsymbol q}_h=[q_p,q_r,q_y,q_e,q_w]^{\mathrm T}$。

\begin{enumerate}[label=\arabic*.]
  \item 检查数据新鲜度、双流时间差、四元数模长和共同时间戳；
  \item 根据当次标定选择原始或转置旋转约定；
  \item 用大臂接收器位置和仿射椭球模型计算 ${}^{B}\hat{\boldsymbol u}$，并检查模型误差；
  \item 用前臂方向对齐矩阵和局部方向轴计算 ${}^{B}\hat{\boldsymbol f}$；
  \item 根据 ${}^{B}\hat{\boldsymbol u}$ 计算 pitch 和 roll，并处理侧举奇异区；
  \item 根据两条肢段方向的内积计算几何肘角 $q_{e,g}$；
  \item 计算 yaw 可观测度；可观测时求解 yaw，否则保持上一帧 yaw；
  \item 使用 $q_{e,g}$ 去除肩部、肘部和接收器固定安装旋转，提取 wrist twist；yaw
        不可观测时同时保持上一帧 wrist；
  \item 对 $q_{e,g}$ 应用本次安装的增益和零偏，得到最终肘部输出 $q_e$；
  \item 将周期角展开到上一帧附近，输出五自由度人体估计；
  \item 记录原始方向、可观测度、原始角、修正角和质量指标，用于离线评价。
\end{enumerate}
\end{quote}

实现中，双接收器采集、标定和在线估计分别由 \texttt{dual\_em\_tracker\_node}、
\texttt{dual\_calibration\_node} 和 \texttt{dual\_retargeting\_node} 完成。纯几何运算被
独立放置在 \texttt{retargeting\_kinematics.py} 和
\texttt{dual\_retargeting\_kinematics.py} 两个无 ROS 依赖的运动学模块中，便于通过单元
测试验证坐标变换、奇异区分支和角度连续性。每轮运行结束时，系统将全频率诊断数据保存为
NPZ（NumPy compressed archive，NumPy 压缩数组归档）文件。该格式可在一个压缩文件中按
名称保存多个同长度数组，便于保持时间序列字段之间的一一对应。记录内容包括接收器原始位姿、人体
肢段方向、五个原始关节角、肘角修正前后数值、pitch/yaw 可观测度以及最终命令。这些数据
为第6章实验设计和第7章结果分析提供统一的数据来源。

\section{本章小结}
\label{sec:ch4_summary}

本章提出并实现了面向重复佩戴的双电磁接收器人体左臂估计方法。该方法没有假设传感器在
每次佩戴后保持相同朝向，而是通过静态质量检查、接收器身份检查、独立前臂旋转轴动作和
七个参考姿势，重新辨识当前安装的坐标对齐、肢段方向轴、大臂位置模型以及肘角修正参数。

本章进一步说明了测量信息与模型自由度之间的关系：大臂单位方向提供两个独立自由度，前臂
完整姿态提供三个独立自由度，从而在一般姿态下支持五自由度状态的逐级几何求解。这五个量
是与 G1 关节语义一致的模型广义坐标，而非完整人体上肢的解剖学真值。

为降低软组织和绑带运动对大臂姿态测量的影响，系统使用仿射位置--方向模型恢复大臂方向；
该模型在几何上把单位方向球面映射为安装相关的椭球轨迹，但并不假设人体肩关节本身为椭球。
前臂方向与旋前/旋后使用相互独立的轴和对齐参数。在线阶段，系统从两个肢段方向计算肩部
pitch、roll、yaw 和几何肘角，通过残余旋转的 swing--twist 分解得到 wrist roll，再将
安装级增益和零偏作用于最终肘部输出。针对侧举
时 pitch 退化和直臂时 yaw/wrist 耦合，本章分别设计了带滞回的标准分支和最后可靠值保持
策略。最后，标定质量门槛、同步检查、位置模型误差和看门狗共同阻止低质量数据进入后续
控制链。

下一章将在本章输出的五自由度人体动作基础上，介绍 G1 关节目标生成、关节和速度限制、
躯干碰撞保护以及真机 Arm SDK 安全桥。

。

\chapter{自适应标定与人体手臂姿态估计}
\label{chap:adaptive_calibration_arm_estimation}

\section{安装标定与姿态估计相关研究}
\label{sec:ch4_related_calibration}

可穿戴传感器报告的是传感器外壳相对参考坐标系的位姿，而人体运动分析需要的是肢段坐标系或
关节变量。两者之间的传感器--肢段关系会随佩戴位置和朝向改变，是可穿戴上肢测量中的核心
误差来源。国际生物力学学会（ISB）给出了肩、肘、腕和手部坐标系及关节运动报告的推荐
定义，为不同测量系统之间统一轴向和角度语义提供了基础~\cite{wu2005isb}。Picerno 等通过
解剖标志点建立上肢惯性传感器与肢段坐标系之间的关系，并使用光学
系统验证肩、肘运动学~\cite{picerno2019upper}。Seel 等利用运动学约束从功能动作中识别关节轴，
说明不必要求传感器外壳与解剖轴精确平行~\cite{seel2014imu}。Zabat 等进一步表明，相比单一
静态姿势，多姿态信息可以减小安装偏差和软组织运动引起的角度误差~\cite{zabat2019multiple}。
近期研究还指出，待测系统与参考系统若采用不同的传感器--肢段定义，会形成明显的系统性关节角
偏差~\cite{adjel2025influence}。

现有标定方法大体可分为静态、功能性和多姿态方法。静态方法操作简单，但对参考姿势执行误差
敏感；功能性方法通过重复运动拟合关节轴，降低了对安装朝向的要求，但依赖充分且近似单轴的
动作激励；多姿态方法利用若干互补方向联合估计安装关系，并可通过各姿态残差发现错误动作或
绑带移动。对肩部三自由度、肘部屈伸和前臂轴向旋转同时估计时，单一姿态不能提供足够约束，
尤其在大臂与前臂接近共线时，肘角和轴向旋转的可观测性会降低。

Cutti 等使用惯性与磁传感器进行肩、肘动态测量的研究也表明，统一的肢段定义、安装标定和
磁环境检查是获得可解释上肢角度的共同前提~\cite{cutti2008ambulatory}。

电磁设备还存在“设备场标定”与“人体安装标定”两个不同层次。前者从已知空间位姿建立电磁
观测与接收器位姿之间的映射，并处理工作空间内的位置相关误差；相关综述与标准化评估研究
均指出，金属畸变、工作距离和空间位置会共同影响电磁跟踪误差
~\cite{kindratenko2000survey,frantz2003accuracy,hummel2006evaluation,franz2014electromagnetic,
maclachlan2025calibration}。后者在设备位姿已经
可用的前提下，估计接收器外壳与人体肢段之间的关系。本文不修改设备固件或制造商场模型，而是
针对每次佩戴执行多姿态安装级标定，并把姿态残差、位置模型残差和重复性作为进入在线估计前的
质量门控。

\section{问题描述与总体思路}
\label{sec:ch4_problem_overview}

本研究使用一套电磁位姿跟踪系统获取操作者左臂运动。电磁发射器固定于操作者躯干，两个
接收器分别固定于左大臂和左前臂。系统以发射器坐标系为共同参考，以约 $50\,\mathrm{Hz}$
的频率输出两个接收器的三维位置和姿态。与基于人体外观的视觉方法相比，电磁测量不依赖
纹理和光照，也不存在人体自身造成的直接视觉遮挡。然而，接收器在每次佩戴后的安装朝向
并不相同，绑带和软组织还会使接收器运动偏离理想的刚体关节模型。因此，原始接收器位姿
不能直接作为机器人关节角使用，必须先完成针对本次安装的标定。这里的“原始位姿”专指
设备直接报告的接收器外壳相对发射器的位置和旋转；它尚未去除外壳安装角，也尚未转换为
人体大臂、前臂方向或关节角。

本章所研究的问题可表述为：给定大臂接收器和前臂接收器在电磁发射器坐标系中的同步位姿，
估计人体左臂的五个运动量
\begin{equation}
  \boldsymbol q_h =
  \begin{bmatrix}
    q_p & q_r & q_y & q_e & q_w
  \end{bmatrix}^{\mathrm T},
  \label{eq:ch4_human_arm_state}
\end{equation}
其中，$q_p$、$q_r$ 和 $q_y$ 分别表示肩部前举/后摆、外展/内收和轴向旋转，$q_e$
表示肘部由约 $90^\circ$ 屈曲姿态向伸直方向的变化量，$q_w$ 表示前臂旋前/旋后。
上述变量沿用 Unitree G1 左臂的关节语义，但其数值首先描述人体动作，随后才在重定向和
安全层中被限制为机器人可执行目标。

需要强调的是，式~\eqref{eq:ch4_human_arm_state} 中的五个量是本研究所选运动学模型的
广义坐标，而不是由传感器直接测得的五个解剖学关节角。人体肩胛骨、锁骨、腕部屈伸以及
关节中心的平移均未被单独建模。本文的目标是构造一组与 G1 左臂关节语义一致、能够由当前
两接收器配置稳定估计的低维运动描述，而不是恢复完整的人体上肢骨骼模型。

系统采用“安装级自适应标定”思路。这里的“自适应”并非在线学习或自适应控制，而是指每次
重新佩戴设备后，系统根据本次安装采集的参考动作重新辨识：
\begin{enumerate}[label=\arabic*)]
  \item 电磁发射器坐标系到人体标准坐标系的映射；
  \item 接收器局部坐标系中的肢段纵轴；
  \item 大臂接收器位置与人体大臂方向之间的映射；
  \item 前臂旋前/旋后的局部旋转轴和零位；
  \item 肘角几何估计的增益和偏置。
\end{enumerate}

在线阶段只使用本次安装生成的参数。任一接收器或发射器被拆卸、重新绑扎或在支架中移动
后，原标定不再有效，必须重新执行标定流程。
换言之，“安装级”表示参数在一次佩戴完成后被辨识，并在该次运行期间保持固定；它不表示
参数随每一帧在线更新。这样既允许系统适应重复佩戴造成的安装差异，又避免在线更新参数时
把操作者的真实动作错误地吸收到标定模型中。

\subsection{坐标系与符号}
\label{subsec:ch4_frames_notation}

本章使用表~\ref{tab:ch4_coordinate_frames} 所列坐标系。
\begin{table}[htbp]
  \centering
  \caption{人体手臂估计使用的坐标系}
  \label{tab:ch4_coordinate_frames}
  \begin{tabularx}{\linewidth}{@{}llX@{}}
    \toprule
    符号 & 英文含义 & 物理含义 \\
    \midrule
    $\{T\}$ & Transmitter & 电磁发射器坐标系 \\
    $\{B\}$ & Body & 人体标准坐标系，与机器人约定一致：$+x$ 向前、$+y$ 向左、$+z$ 向上 \\
    $\{U\}$ & Upper arm & 大臂接收器局部坐标系 \\
    $\{F\}$ & Forearm & 前臂接收器局部坐标系 \\
    \bottomrule
  \end{tabularx}
\end{table}

记 ${}^{T}\boldsymbol p_U$ 和 ${}^{T}\boldsymbol p_F$ 为两个接收器在 $\{T\}$ 中的位置，
${}^{T}\boldsymbol R_U$ 和 ${}^{T}\boldsymbol R_F$ 为将接收器局部向量旋转到 $\{T\}$ 的
旋转矩阵。上标表示量所在的表达坐标系，下标表示被描述的对象；例如
${}^{T}\boldsymbol p_U$ 表示“大臂接收器的位置，用发射器坐标系表达”。
${}^{B}\boldsymbol u$ 和 ${}^{B}\boldsymbol f$ 分别表示人体大臂和前臂在
$\{B\}$ 中的单位方向。本文所称“肢段”是由相邻关节界定、在当前模型中近似视为刚体的
身体部分：大臂肢段连接肩与肘，前臂肢段连接肘与腕。所谓单位方向，是指只保留肢段的朝向
而忽略其长度，即
\begin{equation}
  \lVert{}^{B}\boldsymbol u\rVert_2=1,
  \qquad
  \lVert{}^{B}\boldsymbol f\rVert_2=1.
  \label{eq:ch4_unit_segment_directions}
\end{equation}
将单位向量平移到 $\{B\}$ 的数学原点后，其终点位于单位球面上。向量是自由向量，平移
起点不会改变其方向；因此这里的单位球面只是统一表示所有三维方向的数学空间，并不要求
人体肩关节位于 $\{B\}$ 原点，也不表示接收器在物理上沿半径为 1 的球面运动。

为与 G1 关节模型保持一致，本研究将零位大臂方向和零位前臂方向定义为
\begin{equation}
  \boldsymbol a_{u,0}=\begin{bmatrix}0&0&-1\end{bmatrix}^{\mathrm T},
  \qquad
  \boldsymbol a_{f,0}=\begin{bmatrix}1&0&0\end{bmatrix}^{\mathrm T}.
  \label{eq:ch4_zero_segment_axes}
\end{equation}
因此，零位对应“大臂自然下垂、肘部弯曲约 $90^\circ$、前臂水平向前”的姿态。
这里的“零位方向”是运动学模型在所有广义坐标均为零时规定的参考方向，并非接收器局部
坐标轴，也不等同于医学上的
标准解剖姿势。当所有模型关节变量为零时，大臂和前臂方向分别等于
$\boldsymbol a_{u,0}$ 和 $\boldsymbol a_{f,0}$；其他姿态均由关节旋转作用于这两个参考
方向得到。图~\ref{fig:ch4_frames_zero_pose} 汇总了坐标轴、接收器和零位方向之间的关系。

\begin{figure}[htbp]
  \centering
  \fbox{\parbox{0.90\linewidth}{\centering
    示意图内容：人体左臂处于 DOWN 零位；标出 $\{B\}$ 的 $+x$ 前、$+y$ 左、$+z$ 上，
    发射器坐标系 $\{T\}$，大臂与前臂接收器坐标系 $\{U\}$、$\{F\}$，以及
    $\boldsymbol a_{u,0}=[0,0,-1]^{\mathrm T}$ 和
    $\boldsymbol a_{f,0}=[1,0,0]^{\mathrm T}$。}}
  \caption{双接收器安装、坐标系与运动学零位示意（待替换为正式绘图）}
  \label{fig:ch4_frames_zero_pose}
\end{figure}

\subsection{五自由度运动学模型}
\label{subsec:ch4_five_dof_model}

本研究用肩部三个旋转、肘部一个旋转以及前臂一个轴向旋转描述左臂运动，即
\begin{equation}
  3\ \text{(shoulder)}+1\ \text{(elbow)}+1\ \text{(forearm twist)}
  =5\ \text{DOF}.
  \label{eq:ch4_five_dof_count}
\end{equation}
肩部三个变量 $q_p$（pitch，前举/后摆）、$q_r$（roll，外展/内收）和 $q_y$（yaw，绕
大臂纵轴旋转）共同确定大臂的空间方向及肘平面朝向；$q_e$ 描述肘部由弯曲向伸直的变化，
$q_w$ 描述前臂旋前/旋后。本模型没有包含 G1 腕部的 pitch 和 yaw，也没有单独描述肩胛带
运动，因此应理解为面向动作重定向的五自由度表示，而非完整的人体上肢解剖模型。

肩部和肘部的正向运动学采用
\begin{align}
  {}^{B}\boldsymbol R_S
  &=\boldsymbol R_y(q_p)\boldsymbol R_x(q_r)\boldsymbol R_z(q_y),
  \label{eq:ch4_shoulder_rotation}\\
  {}^{B}\boldsymbol R_E
  &={}^B\boldsymbol R_S\boldsymbol R_y(q_e),
  \label{eq:ch4_forearm_rotation}
\end{align}
其中，${}^{B}\boldsymbol R_S$ 表示肩部旋转，${}^{B}\boldsymbol R_E$ 表示不包含腕部
旋转的前臂坐标系。相应的肢段方向为
\begin{equation}
  {}^{B}\boldsymbol u={}^B\boldsymbol R_S\boldsymbol a_{u,0},
  \qquad
  {}^{B}\boldsymbol f={}^B\boldsymbol R_E\boldsymbol a_{f,0}.
  \label{eq:ch4_segment_directions}
\end{equation}

式~\eqref{eq:ch4_shoulder_rotation} 的乘法顺序是模型定义的一部分，不能任意交换。本文采用
列向量和主动旋转约定，因此右侧旋转先作用：按矩阵作用顺序，零位肢段依次接受
$\boldsymbol R_z(q_y)$、$\boldsymbol R_x(q_r)$ 和 $\boldsymbol R_y(q_p)$。各轴均以本章统一的
矩阵约定解释，而不是在推导中临时更换为另一套欧拉角约定。矩阵乘法不满足交换律，改变顺序会改变同一组
角度对应的空间姿态。本顺序与工程中 G1 左肩三个关节的语义保持一致，而不是声称它是唯一的
人体解剖角定义；解剖学结果的轴序应另按 ISB 推荐报告~\cite{wu2005isb}。

该模型的优点是人体估计量与 G1 五个左臂关节具有一致的语义，避免了先恢复完整人体骨架
再进行第二次逆运动学求解。模型的限制是：仅由两个肢段方向不能在所有姿态下唯一分离肩部
轴向旋转与前臂旋转。第~\ref{sec:ch4_observability_robustness} 节将对这一可观测性问题进行
专门处理。

\subsection{测量信息与五自由度状态的对应关系}
\label{subsec:ch4_measurement_dof_correspondence}

为了说明两个接收器为何能够支持五个运动量的估计，需要区分“数据分量个数”和“独立
自由度个数”。一个三维单位方向属于单位球面
\begin{equation}
  S^2=\left\{\boldsymbol v\in\mathbb R^3\mid
  \boldsymbol v^{\mathrm T}\boldsymbol v=1\right\}.
  \label{eq:ch4_unit_sphere_definition}
\end{equation}
虽然 $\boldsymbol v$ 具有三个坐标分量，但单位长度约束消除了一个独立变量，因此一个方向
只有两个自由度。直观上，这两个自由度类似球面上的经度和纬度；它们能够确定肢段指向，
却不能确定肢段绕自身纵轴旋转了多少。

旋转矩阵属于特殊正交群
\begin{equation}
  SO(3)=\left\{\boldsymbol R\in\mathbb R^{3\times3}\mid
  \boldsymbol R^{\mathrm T}\boldsymbol R=\boldsymbol I,
  \det(\boldsymbol R)=1\right\}.
  \label{eq:ch4_so3_definition}
\end{equation}
旋转矩阵虽然包含九个元素，但正交和单位长度约束使其只有三个独立自由度。也可以将其理解
为“一个轴的空间方向”所需的两个自由度，加上“绕该轴旋转”所需的一个自由度。

“安装标定”与“在线估计”是前后相接的两个阶段。安装标定只在重新佩戴后执行，利用已知
参考姿势求出坐标对齐、局部肢段轴、位置模型和角度修正参数；在线估计则在每一个同步采样
时刻使用这些固定参数，把未知动作转换为五自由度状态，不再重新拟合参数。完成安装标定后，
在线估计实际使用两类几何信息：大臂单位方向
${}^{B}\hat{\boldsymbol u}\in S^2$，以及经过安装对齐的前臂完整姿态
${}^{B}\boldsymbol R_F\in SO(3)$。因此可用独立信息量为
\begin{equation}
  \underbrace{2}_{\text{upper-arm direction}}+
  \underbrace{3}_{\text{forearm attitude}}=5.
  \label{eq:ch4_measurement_dof_count}
\end{equation}
前臂姿态的三个自由度还可以分解为前臂方向 ${}^{B}\hat{\boldsymbol f}$ 的两个自由度和
绕前臂纵轴旋转的一个自由度。由此形成表~\ref{tab:ch4_measurement_to_joint} 所示的逐级
求解关系。

\begin{table}[htbp]
  \centering
  \caption{测量几何量与五自由度状态的对应关系}
  \label{tab:ch4_measurement_to_joint}
  \begin{tabularx}{\linewidth}{@{}lXX@{}}
    \toprule
    待估计量 & 主要几何依据 & 物理含义 \\
    \midrule
    $q_p,q_r$ & 大臂单位方向 ${}^{B}\hat{\boldsymbol u}$ & 大臂指向的两个自由度 \\
    $q_e$ & ${}^{B}\hat{\boldsymbol u}$ 与 ${}^{B}\hat{\boldsymbol f}$ 的夹角
      & 肘部由弯曲向伸直的变化量 \\
    $q_y$ & 前臂方向在大臂法平面内的投影 & 肘平面绕大臂纵轴的朝向 \\
    $q_w$ & 去除肩、肘和固定安装旋转后的前臂残余姿态
      & 绕前臂纵轴的旋前/旋后 \\
    \bottomrule
  \end{tabularx}
\end{table}

上述自由度计数只说明一般姿态下信息量足够，并不保证所有姿态都存在唯一解。当大臂正侧举
时，pitch 方向退化；当肘部完全伸直时，肩部轴向旋转与前臂旋转发生耦合。换言之，该测量
映射在一般姿态下局部可辨识，但在特定几何构型下秩会下降。系统对这些区域的处理将在
第~\ref{sec:ch4_observability_robustness} 节给出。

\section{双接收器同步采集与数据预处理}
\label{sec:ch4_synchronized_acquisition}

\subsection{同步位姿对}
\label{subsec:ch4_synchronized_pose_pairs}

两个 WEMT 接收器分别通过网络返回独立的数据流。如果直接使用两条消息各自最近的样本，
快速运动时可能把不同时刻的大臂和前臂姿态组合在一起，从而产生虚假的肘角和肩部旋转。
因此，采集节点只在以下条件同时满足时发布一对位姿：
\begin{enumerate}[label=\arabic*)]
  \item 两个接收器均处于连接状态且没有设备错误；
  \item 两个接收器都产生了新的硬件帧；
  \item 任一位姿的接收时间不超过 $100\,\mathrm{ms}$；
  \item 两个接收器的到达时间差不超过 $30\,\mathrm{ms}$；
  \item 两个原始四元数均通过有效性检查。
\end{enumerate}

通过检查后，系统给两个 \texttt{geometry\_msgs/PoseStamped} 消息赋予同一个 ROS 时间戳。
该消息由三维位置、四元数姿态、参考坐标系标识和时间戳组成。两条消息分别发布至 ROS 话题
\texttt{/em/upper\_arm\_pose} 和 \texttt{/em/forearm\_pose}。这样，后续标定和在线估计
可按相同时间戳进行严格配对。WEMT 接口内部的平滑窗口被设为 1，即发布未经时间平均的
单帧样本；这里的 1 是窗口中包含的样本数，而不是 1 秒。该设置避免位置和姿态使用不同
滤波造成额外的相位差；统一滤波在完成双流配对后进行。

\subsection{四元数有效性与旋转矩阵}
\label{subsec:ch4_quaternion_validation}

设设备返回四元数为
\begin{equation}
  \boldsymbol q=\begin{bmatrix}q_w&q_x&q_y&q_z\end{bmatrix}^{\mathrm T}.
  \label{eq:ch4_quaternion}
\end{equation}
系统首先检查其模长是否满足
\begin{equation}
  0.95\leq \lVert\boldsymbol q\rVert_2\leq1.05.
  \label{eq:ch4_quaternion_norm_gate}
\end{equation}

未通过检查的数据对被直接丢弃，而不是先强制归一化后继续使用。这是因为持续出现约
0.707 等异常模长时，问题通常来自设备原始数据失效；单纯归一化会掩盖错误并把无效姿态
传入控制链。有效四元数归一化后转换为旋转矩阵：
\begin{equation}
\boldsymbol R(\boldsymbol q)=
\begin{bmatrix}
1-2(q_y^2+q_z^2) & 2(q_xq_y-q_zq_w) & 2(q_xq_z+q_yq_w)\\
2(q_xq_y+q_zq_w) & 1-2(q_x^2+q_z^2) & 2(q_yq_z-q_xq_w)\\
2(q_xq_z-q_yq_w) & 2(q_yq_z+q_xq_w) & 1-2(q_x^2+q_y^2)
\end{bmatrix}.
\label{eq:ch4_quaternion_to_rotation}
\end{equation}

由于当前设备接口只明确四元数分量顺序，而接收器到发射器的主动/被动旋转方向仍需由系统
侧验证，系统在标定时
分别评价原旋转矩阵 $\boldsymbol R$ 和转置矩阵 $\boldsymbol R^{\mathrm T}$。选择标准以前臂
旋转轴在多参考姿势下的方向残差为主，以相对旋转残差作为分数接近时的辅助判据。最终选定
的约定被写入当次标定文件，并在在线估计中保持一致。这一步不是任意改变四元数分量，而是
在互为逆旋转的两种物理解释之间进行一次安装级一致性检验；若后续由设备协议明确方向，仍可
保留该检验作为接线或固件变化的诊断。

\section{安装级自适应标定}
\label{sec:ch4_installation_calibration}

\subsection{标定流程概览}
\label{subsec:ch4_calibration_overview}

完整标定由三个前置质量/动作检查和七个静态参考姿势组成，如图~\ref{fig:ch4_calibration_flow}
所示。
\begin{figure}[htbp]
  \centering
  \fbox{\parbox{0.88\linewidth}{\centering
    图示占位：静态质量检查 $\rightarrow$ 接收器身份检查 $\rightarrow$
    前臂旋转轴标定 $\rightarrow$ 七姿势采集 $\rightarrow$
    参数拟合 $\rightarrow$ 质量门槛 $\rightarrow$ 保存或拒绝。}}
  \caption{安装级自适应标定流程（待替换为正式绘图）}
  \label{fig:ch4_calibration_flow}
\end{figure}

标定开始前，操作者保持自然站立，系统采集 $4\,\mathrm{s}$ 静态数据，用于计算位置和姿态
噪声。随后，操作者在保持大臂基本不动的情况下反复屈伸肘部并转动前臂。若前臂接收器的
姿态变化明显大于大臂接收器，才认为两个接收器的身份与安装位置正确。最后，操作者保持大臂
竖直、肘部弯曲约 $90^\circ$，仅使前臂绕“肘到腕”的纵轴做旋前/旋后，用于辨识前臂接收器
局部坐标系中的旋转轴。

七个静态参考姿势及其理论角度见表~\ref{tab:ch4_reference_poses}。每个姿势离开后重新构建
3 次，每次保持 $2\,\mathrm{s}$。这种“离开---重建---保持”的方式能够发现单次摆姿偶然
正确但重复性不足的问题。

\begin{table}[htbp]
  \centering
  \caption{安装级标定使用的静态参考姿势}
  \label{tab:ch4_reference_poses}
  \begin{tabularx}{\linewidth}{@{}XrrrrX@{}}
    \toprule
    参考姿势 & $q_p$ & $q_r$ & $q_y$ & $q_e$ & 主要作用 \\
    \midrule
    DOWN             & $0^\circ$   & $0^\circ$  & $0^\circ$ & $0^\circ$  & 定义系统零位 \\
    FORWARD\_30      & $-30^\circ$ & $0^\circ$  & $0^\circ$ & $0^\circ$  & 约束大臂位置模型 \\
    FORWARD\_60      & $-60^\circ$ & $0^\circ$  & $0^\circ$ & $0^\circ$  & 约束大臂位置模型 \\
    FORWARD          & $-90^\circ$ & $0^\circ$  & $0^\circ$ & $0^\circ$  & 约束前举方向 \\
    SIDEWAYS         & $0^\circ$   & $90^\circ$ & $0^\circ$ & $90^\circ$ & 约束侧举和伸肘方向 \\
    STRAIGHT\_DOWN   & $0^\circ$   & $0^\circ$  & $0^\circ$ & $90^\circ$ & 约束伸肘姿态 \\
    STRAIGHT\_FORWARD& $-90^\circ$ & $0^\circ$  & $0^\circ$ & $90^\circ$ & 约束主要工作空间 \\
    \bottomrule
  \end{tabularx}
\end{table}

图~\ref{fig:ch4_reference_pose_set} 给出七个参考姿势应呈现的肢段关系。FORWARD\_30 和
FORWARD\_60 的角度不要求操作者用量角器达到绝对精确值；其作用是为前举路径增加两个
中间方向。操作者应优先保证躯干不动、动作平面一致、肘部状态正确以及三次重复可重现，实际
偏差将反映在参考残差和位置模型残差中。

\begin{figure}[htbp]
  \centering
  \fbox{\parbox{0.92\linewidth}{\centering
    示意图内容：依次绘制 DOWN、FORWARD\_30、FORWARD\_60、FORWARD、SIDEWAYS、
    STRAIGHT\_DOWN 和 STRAIGHT\_FORWARD；用实线表示大臂、虚线表示前臂，并标注
    “弯肘约 $90^\circ$”或“伸肘”以及前举/侧举方向。}}
  \caption{七个安装标定参考姿势及其肢段约束（待替换为实拍图或正式绘图）}
  \label{fig:ch4_reference_pose_set}
\end{figure}

FORWARD\_30 和 FORWARD\_60 只用于提高大臂位置模型在前举路径上的可辨识性；前臂方向
拟合使用 DOWN、FORWARD、SIDEWAYS、STRAIGHT\_DOWN 和 STRAIGHT\_FORWARD 五类姿势。

\subsection{静态质量和重复性评价}
\label{subsec:ch4_static_quality_repeatability}

对于一组静态位置样本 $\boldsymbol p_i$，位置均方根波动定义为
\begin{equation}
  e_{p,\mathrm{RMS}}=
  \sqrt{\frac{1}{N}\sum_{i=1}^{N}
  \lVert\boldsymbol p_i-\bar{\boldsymbol p}\rVert_2^2}.
  \label{eq:ch4_static_position_rms}
\end{equation}

姿态中心采用旋转矩阵均值在 $SO(3)$ 上的最近旋转，姿态波动则用每个样本到中心旋转的
测地角计算 RMS。由于旋转不能像普通欧氏向量那样逐元素平均后直接使用，最近旋转投影和
四元数/旋转平均是姿态处理中常用的几何做法~\cite{hartley2013rotation,markley2007averaging}。
当前工程实现中，当任一接收器的位置 RMS 超过 $3\,\mathrm{mm}$，或姿态
RMS 超过 $3^\circ$ 时，系统提示检查绑带、金属物体和跟踪距离。

每个参考姿势的三次重复使用旋转 medoid（样本中心）作为代表，即只在三个实际观测中选择
一个，使它到其余观测的旋转测地距离总和最小。与直接平均相比，medoid 必然是一个真实采集
样本，对单次离群姿态更稳健。若任一重复与代表样本之间的最大角距离超过 $15^\circ$，则只重新采集当前
姿势，已经通过的姿势继续保留。该策略避免一次摆姿错误使整轮标定全部作废，同时保证每个
参考姿势具有最低重复性。

\subsection{前臂局部旋转轴辨识}
\label{subsec:ch4_forearm_twist_axis}

在理想旋前/旋后动作中，前臂在发射器坐标系中的纵轴方向保持近似不变，而接收器绕该轴
转动。实际只有一个前臂接收器，下标 $i$ 表示该接收器在第 $i$ 个采样时刻的观测。设其旋转
矩阵为 ${}^{T}\boldsymbol R_{F,i}$，希望找到接收器外壳坐标系中的局部单位向量
$\boldsymbol a_t$，使
\begin{equation}
  {}^{T}\boldsymbol R_{F,i}\boldsymbol a_t\approx\boldsymbol d_t,
  \qquad i=1,\ldots,N,
  \label{eq:ch4_twist_axis_invariance}
\end{equation}
其中，$\boldsymbol d_t$ 为发射器坐标系中的稳定旋转轴。寻找 $\boldsymbol a_t$ 的目的，是把
任意安装朝向下的“接收器外壳某一轴”辨识为人体前臂纵轴；若轴选择正确，接收器即使绕它
转动，该轴变换到发射器坐标系后的方向仍应近似不变。令
\begin{equation}
  \bar{\boldsymbol R}_F=\frac{1}{N}
  \sum_{i=1}^{N}{}^{T}\boldsymbol R_{F,i},
  \label{eq:ch4_mean_forearm_rotation}
\end{equation}
则 $\boldsymbol a_t$ 可写为
\begin{equation}
  \boldsymbol a_t=
  \arg\max_{\lVert\boldsymbol a\rVert=1}
  \lVert\bar{\boldsymbol R}_F\boldsymbol a\rVert_2,
  \label{eq:ch4_twist_axis_svd}
\end{equation}
即 $\bar{\boldsymbol R}_F$ 最大奇异值对应的右奇异向量。式
~\eqref{eq:ch4_twist_axis_invariance} 给出物理不变性，式
~\eqref{eq:ch4_mean_forearm_rotation} 将全部姿态观测聚合，式
~\eqref{eq:ch4_twist_axis_svd} 则寻找经这些旋转后最不发散的局部方向；最大奇异值对应方向
正是保持合向量长度最大的方向。系统对残差进行基于中位数绝对
偏差的迭代剔除，降低偶发姿态跳变的影响；同时通过样本对之间的相对旋转检查动作范围。
只有当前臂旋转范围至少达到 $55^\circ$、大臂姿态变化不超过 $15^\circ$，且旋转轴拟合
RMS 不超过 $15^\circ$ 时，才继续执行后续标定。

该旋转轴专门服务于腕部旋前/旋后估计。它不与前臂“肘到腕”的方向轴强制共用，因为皮肤
滑动、绑带安装和参考动作误差会使两者的最优拟合结果存在差异。实际实现分别辨识前臂方向
轴与旋转轴，以避免一个轴同时承担两个不同物理任务造成系统性偏差。

\subsection{大臂方向的仿射位置--方向模型}
\label{subsec:ch4_upper_arm_affine_model}

大臂接收器固定在软组织表面，其姿态容易受到皮肤相对肱骨转动的影响。实验开发过程中发现，
直接使用大臂接收器姿态估计肱骨轴向运动会引入较大误差。因此，当前混合模型主要利用大臂
接收器的位置恢复大臂方向。这里的目标不是根据位置直接回归关节角，而是先恢复具有明确
几何意义的单位方向 ${}^{B}\hat{\boldsymbol u}$，再由第~\ref{sec:ch4_online_arm_estimation}
节的逆运动学关系求解关节变量。

若接收器与肩关节构成理想刚体，大臂运动时接收器应在以肩关节为中心的球面上运动。但实际
轨迹受绑带、软组织和发射器安装影响，前举路径往往呈现非球形畸变。为此，本研究采用仿射
位置--方向模型
\begin{equation}
  {}^{T}\boldsymbol p_{U,i}
  =\boldsymbol c+\boldsymbol M\,{}^{B}\boldsymbol u_i^{*}
  +\boldsymbol\varepsilon_i,
  \label{eq:ch4_upper_affine_model}
\end{equation}
其中，$\boldsymbol c\in\mathbb R^3$ 为本次安装下的轨迹中心，
$\boldsymbol M\in\mathbb R^{3\times3}$ 为从人体标准方向到发射器坐标位移的可逆仿射
矩阵，${}^{B}\boldsymbol u_i^{*}$ 是第 $i$ 个参考姿势的理论大臂方向，
$\boldsymbol\varepsilon_i$ 汇集摆姿误差、软组织运动和测量噪声。$\boldsymbol c$ 是用于
位置--方向映射的虚拟中心，不应直接解释为解剖学肩关节中心。

这里称为“混合模型”，是因为大臂方向主要由接收器\emph{位置}恢复，而前臂方向和轴向
旋转由接收器\emph{姿态}恢复。所谓“恢复大臂方向”，是指由测得的位置
${}^{T}\boldsymbol p_U$ 反求单位向量 ${}^{B}\hat{\boldsymbol u}$，而不是估计大臂长度或
肩关节中心。图~\ref{fig:ch4_affine_geometry} 给出式~\eqref{eq:ch4_upper_affine_model} 的
几何含义。

\begin{figure}[htbp]
  \centering
  \fbox{\parbox{0.92\linewidth}{\centering
    示意图内容：以 $\boldsymbol c$ 为椭球中心，单位方向
    ${}^{B}\boldsymbol u_i^{*}$ 经过 $\boldsymbol M$ 映射为从 $\boldsymbol c$ 指向
    ${}^{T}\boldsymbol p_{U,i}$ 的位移；标出实测点、重投影点 $\hat{\boldsymbol p}_U$、
    残差 $\boldsymbol\varepsilon_i$ 与模型误差 $e_M$。左侧可并列绘制单位球，右侧绘制其
    经过 $\boldsymbol M$ 后的椭球像。}}
  \caption{大臂仿射位置--方向模型的几何解释（待替换为正式绘图）}
  \label{fig:ch4_affine_geometry}
\end{figure}

该模型被称为“仿射椭球”，是因为单位方向满足
$({}^{B}\boldsymbol u)^{\mathrm T}{}^{B}\boldsymbol u=1$。忽略噪声并令
$\boldsymbol x={}^{T}\boldsymbol p_U-\boldsymbol c$，由
$\boldsymbol x=\boldsymbol M{}^{B}\boldsymbol u$ 可得
\begin{equation}
  \boldsymbol x^{\mathrm T}\boldsymbol M^{-\mathrm T}
  \boldsymbol M^{-1}\boldsymbol x=1,
  \label{eq:ch4_affine_ellipsoid_implicit}
\end{equation}
这里引入 $\boldsymbol x$ 只是为了简化表达：它是从虚拟中心 $\boldsymbol c$ 指向当前接收器
位置的位移，不是新的测量量。将 ${}^{B}\boldsymbol u=\boldsymbol M^{-1}\boldsymbol x$
代回单位长度约束即可得到式~\eqref{eq:ch4_affine_ellipsoid_implicit}，这正是以
$\boldsymbol c$ 为中心的椭球面二次型。若对矩阵作奇异值分解
\begin{equation}
  \boldsymbol M=\boldsymbol U\boldsymbol\Sigma\boldsymbol V^{\mathrm T},
  \qquad
  \boldsymbol\Sigma=\operatorname{diag}(\sigma_1,\sigma_2,\sigma_3),
  \label{eq:ch4_affine_matrix_svd}
\end{equation}
则 $\boldsymbol U$ 给出椭球主轴方向，$\sigma_1,\sigma_2,\sigma_3$ 给出三个半轴长度。
当三个奇异值相等时，说明三个主方向的尺度相同；若
$\boldsymbol M=r\boldsymbol R$，线性变换只包含统一尺度 $r$ 和旋转 $\boldsymbol R$，模型便
退化为半径为 $r$ 的球面。若奇异值不同，单位球在三个主方向上被拉伸为不同半轴长度，形成
椭球。因此，椭球是单位方向集合经过本次安装相关线性变换后的位置轨迹，而不是对人体肩关节形状的
解剖学假设。实际七个参考姿势只覆盖该椭球面的有限区域，本模型应理解为工作空间内的经验
标定模型，而非完整肩部轨迹的全局物理模型。

所有参考样本通过最小二乘求解
\begin{equation}
  (\hat{\boldsymbol c},\hat{\boldsymbol M})=
  \arg\min_{\boldsymbol c,\boldsymbol M}
  \sum_i\left\|
  {}^{T}\boldsymbol p_{U,i}-\boldsymbol c-
  \boldsymbol M{}^{B}\boldsymbol u_i^{*}
  \right\|_2^2.
  \label{eq:ch4_upper_affine_fit}
\end{equation}
式~\eqref{eq:ch4_upper_affine_fit} 对 $\boldsymbol c$ 和 $\boldsymbol M$ 是线性最小二乘
问题。它求出的 $\hat{\boldsymbol c}$ 和 $\hat{\boldsymbol M}$ 用于在线阶段把新的接收器
位置反算为大臂方向。为了把常量中心项和三维方向项写在同一个线性方程中，将每个理论方向
在末尾附加常数 1，形成增广向量
$[({}^{B}\boldsymbol u_i^{*})^{\mathrm T},1]$。相应设计矩阵必须具有满列秩 4，即四个列方向
都具有独立信息；否则轨迹
中心和全部仿射分量不能被唯一辨识。DOWN、FORWARD 和 SIDEWAYS 三类方向不足以满足这一
条件，因此标定流程额外引入 FORWARD\_30 和 FORWARD\_60，以增加前举路径上的方向变化
并改善模型条件。

在线阶段，根据当前大臂接收器位置计算未归一化方向
\begin{equation}
  \tilde{\boldsymbol u}=
  \hat{\boldsymbol M}^{-1}
  ({}^{T}\boldsymbol p_U-\hat{\boldsymbol c}),
  \label{eq:ch4_upper_direction_raw}
\end{equation}
并归一化得到
\begin{equation}
  {}^{B}\hat{\boldsymbol u}=
  \frac{\tilde{\boldsymbol u}}
  {\lVert\tilde{\boldsymbol u}\rVert_2}.
  \label{eq:ch4_upper_direction_normalized}
\end{equation}
若当前点严格位于标定椭球面上，则
$\lVert\tilde{\boldsymbol u}\rVert_2=1$；实际数据受噪声影响通常偏离曲面，因此仍需归一化。
这一过程可以理解为先用逆仿射变换消除安装相关的方向性畸变，再将结果投影回单位方向空间。

矩阵 $\hat{\boldsymbol M}$ 的奇异值表示沿三个主方向的位移尺度。最小奇异值过小时，说明
某一方向几乎没有被参考姿势激励，求逆时该方向的微小位置噪声会被显著放大。为此定义条件数
\begin{equation}
  \kappa(\hat{\boldsymbol M})=
  \frac{\sigma_{\max}(\hat{\boldsymbol M})}
       {\sigma_{\min}(\hat{\boldsymbol M})}.
  \label{eq:ch4_affine_condition_number}
\end{equation}
条件数越接近 1，三个方向的尺度越均衡；条件数过大则表明参考方向分布不足、摆姿不一致或
拟合接近奇异。当前标定质量门槛要求 $\kappa(\hat{\boldsymbol M})\leq4.0$。

为了检测设备移动或当前动作严重超出标定范围，系统将估计方向重新投影到椭球表面：
\begin{equation}
  \hat{\boldsymbol p}_U=
  \hat{\boldsymbol c}+
  \hat{\boldsymbol M}{}^{B}\hat{\boldsymbol u},
  \label{eq:ch4_upper_position_reprojection}
\end{equation}
并定义位置模型误差
\begin{equation}
  e_M=\left\|{}^{T}\boldsymbol p_U-
  \hat{\boldsymbol p}_U\right\|_2.
  \label{eq:ch4_upper_model_error}
\end{equation}
在线运行时，若 $e_M>50\,\mathrm{mm}$，当前帧被保护性丢弃。该误差不等同于人体大臂
方向的绝对真值误差，而是接收器位置与本次安装模型的一致性指标。小的 $e_M$ 只能说明
当前数据与标定轨迹相容；要评价大臂方向的绝对精度，仍需要独立的光学动作捕捉或其他真值
测量系统。

\subsection{前臂方向与坐标对齐}
\label{subsec:ch4_forearm_direction_alignment}

设 $\boldsymbol a_d$ 为前臂接收器局部坐标系中的“肘到腕”方向，$\boldsymbol A_d$ 为发射器
坐标系到人体标准坐标系的旋转。对于第 $i$ 个参考姿势，预测方向为
\begin{equation}
  {}^{B}\hat{\boldsymbol f}_i=
  \boldsymbol A_d\,{}^{T}\boldsymbol R_{F,i}\boldsymbol a_d.
  \label{eq:ch4_forearm_direction_prediction}
\end{equation}

式~\eqref{eq:ch4_forearm_direction_prediction} 表示两步坐标变换：先用
${}^{T}\boldsymbol R_{F,i}$ 把接收器内部的候选方向 $\boldsymbol a_d$ 旋转到发射器坐标系，
再用 $\boldsymbol A_d$ 将其统一表达在人体坐标系。方向轴和对齐矩阵通过交替优化求解：固定
$\boldsymbol a_d$ 时，$\boldsymbol A_d$ 可归结为
Wahba 旋转对齐问题，即寻找一个旋转，使多组测量方向与对应理论方向的平方误差之和最小
~\cite{wahba1965least}；固定 $\boldsymbol A_d$ 时，可将理论方向反变换到各接收器局部坐标系
并取均值，更新 $\boldsymbol a_d$。算法从六个坐标轴正负方向初始化，保留参考姿势方向 RMS
最小的解。

旋转对齐子问题为
\begin{equation}
  \hat{\boldsymbol A}_d=
  \arg\min_{\boldsymbol A\in SO(3)}
  \sum_i\left\|
  \boldsymbol A\,{}^{T}\boldsymbol R_{F,i}\boldsymbol a_d-
  {}^{B}\boldsymbol f_i^{*}
  \right\|_2^2,
  \label{eq:ch4_forearm_wahba_problem}
\end{equation}
通过互协方差矩阵的奇异值分解求得，并显式约束 $\det(\boldsymbol A)=1$。在线阶段的前臂
单位方向为
\begin{equation}
  {}^{B}\hat{\boldsymbol f}=
  \frac{
    \hat{\boldsymbol A}_d\,{}^{T}\boldsymbol R_F\hat{\boldsymbol a}_d
  }{
    \left\|
      \hat{\boldsymbol A}_d\,{}^{T}\boldsymbol R_F\hat{\boldsymbol a}_d
    \right\|_2
  }.
  \label{eq:ch4_forearm_direction_online}
\end{equation}

除了方向映射，系统还建立独立的姿态对齐矩阵 $\boldsymbol A_t$ 和固定安装旋转
$\boldsymbol R_m$，用于计算 wrist roll。这里的“姿态对齐”是把接收器报告的完整三维旋转
转换到人体模型坐标系；$\boldsymbol R_m$ 则记录接收器外壳相对理论前臂坐标系的固定安装角。
在 DOWN 姿势下，理论前臂姿态作为腕部零位；三次 DOWN 重复的平均旋转确定本次安装的固定
旋转。在线计算时先去除肩、肘造成的理论旋转，再去除 $\boldsymbol R_m$，剩余的绕前臂纵轴
分量才被解释为腕部扭转。由此，外壳安装角和前臂方向误差不会直接被误认为 wrist roll。

\subsection{肘角的安装级仿射修正}
\label{subsec:ch4_elbow_affine_correction}

由两个单位方向可以得到几何肘角
\begin{equation}
  q_{e,g}=\arcsin\left(
  \operatorname{clip}
  ( {}^{B}\hat{\boldsymbol u}^{\mathrm T}
  {}^{B}\hat{\boldsymbol f},-1,1)
  \right).
  \label{eq:ch4_geometric_elbow_angle}
\end{equation}

下标 $g$ 表示 geometric，即由两条肢段方向直接得到的几何角。在本研究的零位约定中，
弯肘约 $90^\circ$ 时两个方向近似正交，因而 $q_{e,g}\approx0$；
手臂伸直时二者近似平行，因而 $q_{e,g}\approx90^\circ$。由于两个方向拟合各自存在固定
残差，几何角可能同时存在零位偏差和幅值缩放误差：固定零偏只能整体平移角度，不能改变从
弯曲端到伸直端的量程；单独改变增益又不能校正零位。因此本研究对每次安装分别拟合一维
仿射修正，即“缩放加平移”模型
\begin{equation}
  q_e=\operatorname{clip}
  (s_e q_{e,g}+b_e,\,0,\,\pi/2),
  \label{eq:ch4_elbow_affine_correction}
\end{equation}
其中增益 $s_e$ 和偏置 $b_e$ 由五类前臂参考姿势的理论肘角和几何肘角进行线性最小二乘
拟合。为拒绝异常标定，当前实现要求 $0.75\leq s_e\leq1.50$，且
$|b_e|\leq20^\circ$。若拟合不可辨识或超出合理范围，系统保留未经修正的几何肘角，同时
输出警告。

\subsection{标定质量门槛与实机资格}
\label{subsec:ch4_calibration_quality_gates}

标定参数并非只要能够求解就可进入真机。系统在保存前检查表~\ref{tab:ch4_quality_gates}
所列工程门槛。
\begin{table}[htbp]
  \centering
  \caption{安装级标定质量门槛}
  \label{tab:ch4_quality_gates}
  \begin{tabularx}{\linewidth}{@{}XllX@{}}
    \toprule
    检查项 & 建议/诊断阈值 & 硬限制 & 不通过时的处理 \\
    \midrule
    参考姿势重复差异 & --- & $15^\circ$ & 重采当前姿势 \\
    大臂方向参考误差 & $10^\circ$ & $15^\circ$ & 警告/拒绝保存 \\
    前臂方向参考误差 & $10^\circ$；真机资格 $20^\circ$ & $25^\circ$ & 警告、仅诊断或拒绝 \\
    前臂腕部零位重复误差 & --- & $15^\circ$ & 拒绝保存 \\
    仿射矩阵条件数 & 越接近 1 越好 & 4.0 & 拒绝保存 \\
    仿射位置模型 RMS/最大误差 & $30/50\,\mathrm{mm}$ & $35/65\,\mathrm{mm}$ & 警告/拒绝保存 \\
    \bottomrule
  \end{tabularx}
\end{table}

表中的“建议/诊断阈值”用于提示边缘质量，“硬限制”决定标定文件能否保存，两者不应简单
理解为“仿真门槛”和“实机门槛”。仿真与 dry-run 仍可使用通过硬限制但质量边缘的标定；
真机资格另由更严格的前臂方向条件控制。若前臂方向误差超过 $20^\circ$ 但仍未超过硬限制
$25^\circ$，标定可以作为诊断资料保存，
但会被标记为 \texttt{dry\_run\_only: true}。该标定只允许用于 dry-run 和离线诊断，
重定向节点拒绝向真机主题发布关节目标。该设计把“算法能够产生数值”和“标定具备实机
资格”明确分开。

这些阈值是本项目的工程安全标准，并非电磁跟踪器或人形机器人领域的通用精度标准。参考
姿势残差描述标定模型与人工构建姿势之间的一致性，不能替代独立视觉系统测得的人体姿态
真值误差。

\section{在线人体手臂姿态估计}
\label{sec:ch4_online_arm_estimation}

安装标定完成后，每一对同步位姿首先被转换为三个中间几何量：大臂单位方向
${}^{B}\hat{\boldsymbol u}$、前臂单位方向 ${}^{B}\hat{\boldsymbol f}$，以及经过腕部姿态
对齐的前臂接收器旋转 ${}^{B}\boldsymbol R_F$。其中，大臂方向由
式~\eqref{eq:ch4_upper_direction_raw}--\eqref{eq:ch4_upper_direction_normalized} 对当前大臂
接收器位置进行逆仿射映射得到；前臂方向由式~\eqref{eq:ch4_forearm_direction_online} 对当前
前臂接收器姿态和已标定局部轴进行变换得到；“腕部姿态对齐”则指用 $\boldsymbol A_t$ 将完整
接收器旋转转换到人体模型坐标系。五个关节量并不是同时进行黑箱优化，而是
按照几何依赖关系逐级求解：先由大臂方向确定 $q_p,q_r$，再由两条肢段方向确定几何肘角
$q_{e,g}$ 和肩部轴向角 $q_y$，最后从前臂残余旋转中提取 $q_w$。几何肘角随后经过本次
安装的增益和偏置修正，形成最终输出 $q_e$。

这种求解顺序保留了每个变量的物理意义，也使不可观测条件能够被显式检测。它与输入信息的
对应关系可概括为
\begin{equation}
  ({}^{T}\boldsymbol p_U,{}^{T}\boldsymbol R_F)
  \longrightarrow
  ({}^{B}\hat{\boldsymbol u},{}^{B}\hat{\boldsymbol f},
   {}^{B}\boldsymbol R_F)
  \longrightarrow
  (q_p,q_r,q_y,q_{e,g},q_w)
  \longrightarrow
  (q_p,q_r,q_y,q_e,q_w).
  \label{eq:ch4_online_estimation_chain}
\end{equation}

\subsection{肩部 pitch 与 roll}
\label{subsec:ch4_shoulder_pitch_roll}

得到大臂单位方向
\begin{equation}
  {}^{B}\hat{\boldsymbol u}=
  \begin{bmatrix}u_x&u_y&u_z\end{bmatrix}^{\mathrm T}
  \label{eq:ch4_upper_direction_components}
\end{equation}
后，由于 $\boldsymbol R_z(q_y)$ 绕零位大臂轴旋转，不改变
$\boldsymbol a_{u,0}=[0,0,-1]^{\mathrm T}$，大臂方向只与 $q_p,q_r$ 有关。展开
式~\eqref{eq:ch4_segment_directions} 可得
\begin{equation}
  {}^{B}\hat{\boldsymbol u}=
  \begin{bmatrix}
    -\sin q_p\cos q_r\\
    \sin q_r\\
    -\cos q_p\cos q_r
  \end{bmatrix}.
  \label{eq:ch4_upper_direction_expanded}
\end{equation}
这里的“零位大臂轴”就是式~\eqref{eq:ch4_zero_segment_axes} 定义的
$\boldsymbol a_{u,0}$。由于该向量位于 $z$ 轴上，绕 $z$ 轴旋转不会改变它，即
$\boldsymbol R_z(q_y)\boldsymbol a_{u,0}=\boldsymbol a_{u,0}$；随后依次左乘
$\boldsymbol R_x(q_r)$ 和 $\boldsymbol R_y(q_p)$，直接进行矩阵乘法即可得到
式~\eqref{eq:ch4_upper_direction_expanded}。

因此，肩部 pitch 和 roll 的主值可写为
\begin{equation}
  q_p=\operatorname{atan2}(-u_x,-u_z),
  \qquad
  q_r=\arcsin(\operatorname{clip}(u_y,-1,1)).
  \label{eq:ch4_pitch_roll_inverse}
\end{equation}

“主值”表示反三角函数在约定区间内返回的代表解，因为同一方向可对应相差 $2\pi$ 的周期
等价角。由式~\eqref{eq:ch4_upper_direction_expanded} 有
$u_y=\sin q_r$，故 $q_r$ 可由反正弦得到；而 $-u_x$ 与 $-u_z$ 分别包含
$\sin q_p\cos q_r$ 和 $\cos q_p\cos q_r$，用双参数反正切可同时保留象限信息并消去共同
因子 $\cos q_r$。$\operatorname{clip}(u_y,-1,1)$ 只用于消除浮点误差造成的轻微越界，
例如数值上出现 $1+10^{-12}$ 时避免 $\arcsin$ 失效，并非裁剪正常人体动作。

由于 \texttt{atan2} 的输出位于 $[-\pi,\pi]$，系统将当前结果展开到最接近上一帧的等价角，
避免跨越 $\pm\pi$ 时产生整周跳变。

当大臂接近正侧方时，$u_x$ 和 $u_z$ 同时接近零，pitch 对测量噪声高度敏感。定义
\begin{equation}
  o_p=\sqrt{u_x^2+u_z^2}=|\cos q_r|
  \label{eq:ch4_pitch_observability}
\end{equation}
为 pitch 可观测度。在线算法采用带滞回的两状态机制：当 $o_p\leq0.50$ 时进入侧举奇异区，
将 pitch 切换到统一的标准解 $q_p=0$；只有当手臂返回水平面以下且 $o_p\geq0.65$ 时才
重新采用测量解。与此同时，roll 使用标准 pitch 坐标系下的 \texttt{atan2} 表达，使其
能够连续通过 $90^\circ$，而不是被 \texttt{asin} 折回到主值范围。该处理保证从 DOWN
直接侧举和从 FORWARD 转向侧举时，同一最终人体姿势映射到同一机器人目标。

由式~\eqref{eq:ch4_upper_direction_expanded} 可见，$u_x$ 和 $u_z$ 都含有因子
$\cos q_r$；当侧举接近 $|q_r|=90^\circ$ 时该因子趋近零，pitch 的分子和分母同时退化。
阈值 0.50 和 0.65 分别约对应 $|q_r|=60^\circ$ 和 $49^\circ$，是当前工程为避免临界点
抖动设置的进入/退出滞回阈值，而不是人体生理常数。进入奇异区后固定 $q_p=0$，相当于从
无穷多个等价 pitch--roll 表示中选择统一标准解；退出阈值更高且要求回到水平面以下，使状态
不会在噪声作用下频繁切换。

\subsection{肘部伸展量}
\label{subsec:ch4_online_elbow_angle}

在运动学模型中，大臂和前臂方向分别为
${}^{B}\hat{\boldsymbol u}={}^B\boldsymbol R_S\boldsymbol a_{u,0}$ 和
${}^{B}\hat{\boldsymbol f}={}^B\boldsymbol R_S
\boldsymbol R_y(q_{e,g})\boldsymbol a_{f,0}$。利用旋转矩阵保持内积的性质，有
\begin{align}
  ({}^{B}\hat{\boldsymbol u})^{\mathrm T}
  {}^{B}\hat{\boldsymbol f}
  &=\boldsymbol a_{u,0}^{\mathrm T}
    \boldsymbol R_y(q_{e,g})\boldsymbol a_{f,0}
  \nonumber\\
  &=\sin q_{e,g}.
  \label{eq:ch4_elbow_dot_product_derivation}
\end{align}
旋转矩阵满足 $\boldsymbol R^{\mathrm T}\boldsymbol R=\boldsymbol I$，因此对任意向量
$\boldsymbol a,\boldsymbol b$ 都有
$(\boldsymbol R\boldsymbol a)^{\mathrm T}(\boldsymbol R\boldsymbol b)
=\boldsymbol a^{\mathrm T}\boldsymbol b$，即共同旋转不改变夹角或内积。在本式中，肩部旋转
${}^{B}\boldsymbol R_S$ 因而被消去；再代入
$\boldsymbol a_{u,0}=[0,0,-1]^{\mathrm T}$、
$\boldsymbol a_{f,0}=[1,0,0]^{\mathrm T}$ 和 $\boldsymbol R_y(q_{e,g})$ 的矩阵形式，内积
恰好等于 $\sin q_{e,g}$。

因此，式~\eqref{eq:ch4_geometric_elbow_angle} 可以直接由两个肢段方向求得，而不依赖肩部
三个角的具体数值。在本研究的约定下，$q_{e,g}=0$ 对应约 $90^\circ$ 屈肘，
$q_{e,g}=\pi/2$ 对应伸直。为避免与临床文献中常用的“屈曲角”混淆，本文将其称为肘部
伸展量。

这种零位来自当前 G1 左臂模型与重定向接口的关节语义：机器人模型零位附近，大臂向下而
前臂向前，二者约成直角；正向 $q_e$ 将前臂逐渐转到与大臂同向。因此本文的 $q_e$ 是
“从约 $90^\circ$ 弯肘零位出发的伸展量”，不等同于临床上通常从完全伸直定义的肘屈曲角。

几何肘角直接反映当前两条估计方向之间的关系。最终发布的肘部变量使用
式~\eqref{eq:ch4_elbow_affine_correction} 进行安装级增益和零偏修正。这里应区分
$q_{e,g}$ 与 $q_e$：前者用于保持几何分解与测量方向一致，后者是面向机器人目标的校正后
输出。
必须保留两者，是因为 wrist 残余姿态分解要求理论前臂方向与当前实测大臂、前臂方向严格
一致，故使用未修正的 $q_{e,g}$；而输出给重定向层的肘角需要补偿安装造成的尺度和零偏，故
使用 $q_e$。若提前用 $q_e$ 构造理论前臂姿态，标定补偿量会被错误地吸收到 wrist twist 中。

\subsection{肩部 yaw}
\label{subsec:ch4_shoulder_yaw}

大臂方向本身不能确定绕大臂轴的旋转。系统使用前臂在垂直于大臂方向平面内的分量估计
肩部 yaw：
\begin{equation}
  \boldsymbol f_{\perp}=
  {}^{B}\hat{\boldsymbol f}-
  \left({}^{B}\hat{\boldsymbol f}^{\mathrm T}
  {}^{B}\hat{\boldsymbol u}\right)
  {}^{B}\hat{\boldsymbol u}.
  \label{eq:ch4_forearm_perpendicular_component}
\end{equation}

若 $\lVert\boldsymbol f_{\perp}\rVert$ 足够大，将其归一化，并与在当前 pitch、roll 下的零
yaw 参考方向比较。几何上，$\boldsymbol f_{\perp}$ 是前臂方向在大臂法平面内的投影，
它确定了“肩--肘--腕”平面绕大臂纵轴的朝向。设零 yaw 参考方向投影并归一化后为
$\boldsymbol f_{0,\perp}$，则
\begin{equation}
  q_y=\operatorname{atan2}
  \left(
  -{}^{B}\hat{\boldsymbol u}^{\mathrm T}
  (\boldsymbol f_{0,\perp}\times\boldsymbol f_{\perp}),
  \boldsymbol f_{0,\perp}^{\mathrm T}\boldsymbol f_{\perp}
  \right).
  \label{eq:ch4_shoulder_yaw_inverse}
\end{equation}

投影的作用是去除前臂沿大臂轴的分量，只保留能够描述肘平面绕大臂轴朝向的分量。归一化
则去除投影长度，只比较方向；但只有投影长度足够大时，归一化才不会放大噪声。零 yaw 参考
方向 $\boldsymbol f_{0,\perp}$ 由当前 $q_p,q_r$ 和 $q_y=0$ 的运动学模型生成，因此比较二者
即可得到实际肘平面相对零 yaw 肘平面的有符号转角。式~\eqref{eq:ch4_shoulder_yaw_inverse}
中，点积给出余弦项，沿大臂轴的叉积投影给出带符号正弦项，$\operatorname{atan2}$ 将二者
合成为保留象限的有符号角。前置负号用于匹配 G1 肩 yaw 的正方向约定。

\begin{figure}[htbp]
  \centering
  \fbox{\parbox{0.92\linewidth}{\centering
    左图：在垂直于大臂方向 ${}^{B}\hat{\boldsymbol u}$ 的平面内，画出前臂投影
    $\boldsymbol f_{\perp}$、零 yaw 参考 $\boldsymbol f_{0,\perp}$ 和有符号角 $q_y$；
    右图：对比弯肘时投影较长、yaw 与 wrist 可分辨，以及直臂时投影趋近零、两种轴向旋转
    耦合的情形。}}
  \caption{肩部 yaw 的平面投影求解及直臂不可观测性（待替换为正式绘图）}
  \label{fig:ch4_yaw_projection_observability}
\end{figure}

与 pitch 相同，yaw 也展开到最接近上一帧的等价角。由于旋转 $q_y$ 与
$q_y+2k\pi$（$k\in\mathbb Z$）表示同一姿态，展开只选择最连续的数值分支，并不改变物理
旋转。

\subsection{前臂旋前/旋后}
\label{subsec:ch4_forearm_pronation_supination}

当前实现使用与两条测量方向（大臂单位方向 ${}^{B}\hat{\boldsymbol u}$ 和前臂单位方向
${}^{B}\hat{\boldsymbol f}$）直接一致的几何肘角 $q_{e,g}$ 构造不含 wrist roll 的理论
前臂旋转
\begin{equation}
  {}^{B}\boldsymbol R_{E,g}=
  \boldsymbol R_y(q_p)\boldsymbol R_x(q_r)
  \boldsymbol R_z(q_y)\boldsymbol R_y(q_{e,g}).
  \label{eq:ch4_forearm_rotation_without_wrist}
\end{equation}
当前接收器在人体标准坐标系中的旋转为
\begin{equation}
  {}^{B}\boldsymbol R_F=
  \boldsymbol A_t\,{}^{T}\boldsymbol R_F.
  \label{eq:ch4_forearm_sensor_body_rotation}
\end{equation}
去除理论前臂旋转和固定安装旋转后，残余旋转为
\begin{equation}
  \boldsymbol R_{res}=
  ({}^{B}\boldsymbol R_{E,g})^{\mathrm T}
  {}^{B}\boldsymbol R_F\boldsymbol R_m^{\mathrm T}.
  \label{eq:ch4_wrist_residual_rotation}
\end{equation}

上述三式按“预测---观测---求残差”的顺序工作。式
~\eqref{eq:ch4_forearm_rotation_without_wrist} 用已估计的肩、肘量预测前臂在未转掌时的姿态；
式~\eqref{eq:ch4_forearm_sensor_body_rotation} 将传感器观测变换到同一人体坐标系；式
~\eqref{eq:ch4_wrist_residual_rotation} 左乘理论姿态的逆、右乘固定安装旋转的逆，从观测中
去除肩部运动、肘部运动和外壳安装角。理想剩余量只包含前臂绕自身纵轴的旋转。本文沿用
机器人关节名将该量称为 wrist roll，但其人体意义是前臂旋前/旋后，并非手腕屈伸。

等价地，前臂接收器姿态模型可写为
\begin{equation}
  {}^{B}\boldsymbol R_F
  \approx{}^{B}\boldsymbol R_{E,g}
  \boldsymbol R_x(q_w)\boldsymbol R_m.
  \label{eq:ch4_forearm_sensor_attitude_model}
\end{equation}
其中，$\boldsymbol R_m$ 表示本次安装中接收器与前臂模型坐标系之间的固定旋转。

理想情况下，$\boldsymbol R_{res}$ 仅包含绕前臂局部 $+x$ 轴的旋转。为抑制小幅残余 swing
误差，系统将 $\boldsymbol R_{res}$ 转为四元数
$\boldsymbol q_{res}=[q_0,q_x,q_y,q_z]^{\mathrm T}$，只保留向量部分在 $+x$ 轴上的投影，
得到
\begin{equation}
  q_w=\operatorname{wrap}_{[-\pi,\pi)}
  \left(2\operatorname{atan2}(q_x,q_0)\right).
  \label{eq:ch4_wrist_twist_angle}
\end{equation}
然后将结果展开到上一帧附近，以维持时间连续性。这一过程等价于围绕前臂纵轴的
swing--twist 分解：swing 改变纵轴指向，twist 则绕该纵轴旋转；只保留 $q_x$ 相当于提取
局部 $+x$ 轴上的 twist 分量。此类分解常用于把三维关节旋转分离为摆动和轴向扭转
~\cite{murthy2019joint}。结果展开的原因与 yaw 相同：四元数给出的主角位于有限区间，相邻
帧若跨过 $\pm\pi$ 会在数值上跳变约 $2\pi$，选择离上一帧最近的等价角可保持连续。
wrist 提取完成后，系统才将 $q_{e,g}$ 修正为最终输出 $q_e$；因此，
肘角的经验增益和零偏不会被误解释为额外的前臂轴向旋转。

\section{不可观测区域与稳健性处理}
\label{sec:ch4_observability_robustness}

\subsection{直臂状态下 yaw 与 wrist 的耦合}
\label{subsec:ch4_yaw_wrist_unobservable}

当肘部接近完全伸直时，大臂和前臂方向近似共线，
$\lVert\boldsymbol f_{\perp}\rVert\rightarrow0$。此时，仅凭两个肢段方向无法区分“肩部
绕大臂轴旋转”和“前臂自身旋转”；同一物理姿态可以对应多组 $(q_y,q_w)$。这是测量模型
的几何不可观测性，而不是数值求解精度不足。

原因是直臂时肩 yaw 轴与前臂 twist 轴近似重合，末端姿态主要取决于二者的组合量。对任意
小角度 $\delta$，参数对
\begin{equation}
  (q_y,q_w)\sim(q_y+\delta,q_w-\delta)
  \label{eq:ch4_yaw_wrist_equivalent_pairs}
\end{equation}
可产生相同或近似相同的总轴向旋转，因而无法由当前两接收器数据唯一判断该旋转应归属于
肩部还是前臂。弯肘后两条旋转轴不再共线，图
~\ref{fig:ch4_yaw_projection_observability} 所示投影重新出现，两个变量才恢复可分辨性。

系统定义
\begin{equation}
  o_y=\lVert\boldsymbol f_{\perp}\rVert_2
  \label{eq:ch4_yaw_observability}
\end{equation}
为 yaw 可观测度。由于 $\boldsymbol u$ 和 $\boldsymbol f$ 均为单位向量，且
$\boldsymbol u^{\mathrm T}\boldsymbol f=\sin q_{e,g}$，可进一步得到
\begin{equation}
  o_y=\sqrt{1-(\boldsymbol u^{\mathrm T}\boldsymbol f)^2}
  =|\cos q_{e,g}|.
  \label{eq:ch4_yaw_observability_elbow_relation}
\end{equation}
因此，这一退化不是经验上偶然出现的数值问题，而是由五自由度模型的几何结构直接决定。
弯肘约 $90^\circ$ 时 $q_{e,g}\approx0$、$o_y\approx1$，肘平面方向清晰；手臂趋于伸直时
$q_{e,g}\rightarrow\pi/2$、$o_y\rightarrow0$，肩 yaw 轴与前臂 twist 轴趋于重合，测量
中只能保留两者的组合旋转而不能唯一分离。

当 $o_y<0.90$ 时，保持最后一次可靠的 $q_y$ 和 $q_w$，避免把微小方向
噪声放大为几十度的轴向旋转。由于本模型中 $o_y$ 随肘部伸直而下降，该阈值会在完全伸直
之前提前冻结轴向角。系统恢复到明显弯肘姿态后，再继续更新 yaw 和 wrist。

因此，直臂前举和直臂侧举可以用于评价大臂方向及肘角，但不适合
把直臂主动转掌作为 wrist 精度的主要证据。wrist 单关节实验应保持肘部弯曲约 $60^\circ$
至 $90^\circ$，使肩部轴向旋转和前臂旋转保持可分辨。

\subsection{周期角展开与时间连续性}
\label{subsec:ch4_temporal_conditioning}

肩部 yaw、wrist roll 和其他周期角首先展开到最接近上一帧的 $2\pi$ 等价分支。反三角函数
通常只返回 $[-\pi,\pi)$ 内的主值，因此真实角度从 $+\pi$ 连续增加时，下一帧数值会变成
接近 $-\pi$；两帧物理上很接近，数值却相差近 $2\pi$。角展开在
$\theta+2k\pi$ 的等价集合中选择最接近上一帧的分支，从而避免该边界产生整周跳变。侧举
奇异区和直臂不可观测区的内部状态也只在满足带滞回的
重新观测条件后更新，防止临界点附近反复切换。

本章输出的是经过可观测性处理的五自由度人体估计，不在此处施加 G1 关节限位、机器人
命令低通滤波或真机速度限制。上述机器人目标调理统一放在第5章说明，从而把“人体状态
估计”和“机器人命令生成”分开。

\subsection{无效数据与看门狗}
\label{subsec:ch4_invalid_data_watchdog}

在线估计过程中，以下情况不会产生新的关节目标：
\begin{itemize}
  \item 接收器数据陈旧或双流时间差过大；
  \item 原始四元数模长异常；
  \item 两个话题的坐标系标识不一致；
  \item 大臂仿射位置模型误差超过 $50\,\mathrm{mm}$；
  \item 标定矩阵无效或当前方向向量接近零。
\end{itemize}

“看门狗”是监测数据或程序是否在规定时间内持续更新的超时机制。若连续
$0.25\,\mathrm{s}$ 没有新的有效同步数据对，重定向节点发出看门狗警告并停止更新
目标。真机桥还具有独立的 UDP 目标超时和控制权释放机制。因此，采集层、重定向层与执行
层分别检查数据新鲜度，避免单点故障持续驱动机器人。

\section{算法流程与实现}
\label{sec:ch4_algorithm_implementation}

一次完整在线估计可归纳为算法 4-1。

\begin{quote}
\textbf{算法 4-1：双电磁接收器人体左臂姿态估计}

\textbf{输入：}同步接收器位姿
$({}^{T}\boldsymbol p_U,{}^{T}\boldsymbol R_U,
{}^{T}\boldsymbol p_F,{}^{T}\boldsymbol R_F)$，当次安装标定参数，上一帧角度。

\textbf{输出：}五自由度人体手臂估计
$\hat{\boldsymbol q}_h=[q_p,q_r,q_y,q_e,q_w]^{\mathrm T}$。

\begin{enumerate}[label=\arabic*.]
  \item 检查数据新鲜度、双流时间差、四元数模长和共同时间戳；
  \item 根据当次标定选择原始或转置旋转约定；
  \item 用大臂接收器位置和仿射椭球模型计算 ${}^{B}\hat{\boldsymbol u}$，并检查模型误差；
  \item 用前臂方向对齐矩阵和局部方向轴计算 ${}^{B}\hat{\boldsymbol f}$；
  \item 根据 ${}^{B}\hat{\boldsymbol u}$ 计算 pitch 和 roll，并处理侧举奇异区；
  \item 根据两条肢段方向的内积计算几何肘角 $q_{e,g}$；
  \item 计算 yaw 可观测度；可观测时求解 yaw，否则保持上一帧 yaw；
  \item 使用 $q_{e,g}$ 去除肩部、肘部和接收器固定安装旋转，提取 wrist twist；yaw
        不可观测时同时保持上一帧 wrist；
  \item 对 $q_{e,g}$ 应用本次安装的增益和零偏，得到最终肘部输出 $q_e$；
  \item 将周期角展开到上一帧附近，输出五自由度人体估计；
  \item 记录原始方向、可观测度、原始角、修正角和质量指标，用于离线评价。
\end{enumerate}
\end{quote}

实现中，双接收器采集、标定和在线估计分别由 \texttt{dual\_em\_tracker\_node}、
\texttt{dual\_calibration\_node} 和 \texttt{dual\_retargeting\_node} 完成。纯几何运算被
独立放置在 \texttt{retargeting\_kinematics.py} 和
\texttt{dual\_retargeting\_kinematics.py} 两个无 ROS 依赖的运动学模块中，便于通过单元
测试验证坐标变换、奇异区分支和角度连续性。每轮运行结束时，系统将全频率诊断数据保存为
NPZ（NumPy compressed archive，NumPy 压缩数组归档）文件。该格式可在一个压缩文件中按
名称保存多个同长度数组，便于保持时间序列字段之间的一一对应。记录内容包括接收器原始位姿、人体
肢段方向、五个原始关节角、肘角修正前后数值、pitch/yaw 可观测度以及最终命令。这些数据
为第6章实验设计和第7章结果分析提供统一的数据来源。

\section{本章小结}
\label{sec:ch4_summary}

本章提出并实现了面向重复佩戴的双电磁接收器人体左臂估计方法。该方法没有假设传感器在
每次佩戴后保持相同朝向，而是通过静态质量检查、接收器身份检查、独立前臂旋转轴动作和
七个参考姿势，重新辨识当前安装的坐标对齐、肢段方向轴、大臂位置模型以及肘角修正参数。

本章进一步说明了测量信息与模型自由度之间的关系：大臂单位方向提供两个独立自由度，前臂
完整姿态提供三个独立自由度，从而在一般姿态下支持五自由度状态的逐级几何求解。这五个量
是与 G1 关节语义一致的模型广义坐标，而非完整人体上肢的解剖学真值。

为降低软组织和绑带运动对大臂姿态测量的影响，系统使用仿射位置--方向模型恢复大臂方向；
该模型在几何上把单位方向球面映射为安装相关的椭球轨迹，但并不假设人体肩关节本身为椭球。
前臂方向与旋前/旋后使用相互独立的轴和对齐参数。在线阶段，系统从两个肢段方向计算肩部
pitch、roll、yaw 和几何肘角，通过残余旋转的 swing--twist 分解得到 wrist roll，再将
安装级增益和零偏作用于最终肘部输出。针对侧举
时 pitch 退化和直臂时 yaw/wrist 耦合，本章分别设计了带滞回的标准分支和最后可靠值保持
策略。最后，标定质量门槛、同步检查、位置模型误差和看门狗共同阻止低质量数据进入后续
控制链。

下一章将在本章输出的五自由度人体动作基础上，介绍 G1 关节目标生成、关节和速度限制、
躯干碰撞保护以及真机 Arm SDK 安全桥。

。

\chapter{自适应标定与人体手臂姿态估计}
\label{chap:adaptive_calibration_arm_estimation}

\section{安装标定与姿态估计相关研究}
\label{sec:ch4_related_calibration}

可穿戴传感器报告的是传感器外壳相对参考坐标系的位姿，而人体运动分析需要的是肢段坐标系或
关节变量。两者之间的传感器--肢段关系会随佩戴位置和朝向改变，是可穿戴上肢测量中的核心
误差来源。国际生物力学学会（ISB）给出了肩、肘、腕和手部坐标系及关节运动报告的推荐
定义，为不同测量系统之间统一轴向和角度语义提供了基础~\cite{wu2005isb}。Picerno 等通过
解剖标志点建立上肢惯性传感器与肢段坐标系之间的关系，并使用光学
系统验证肩、肘运动学~\cite{picerno2019upper}。Seel 等利用运动学约束从功能动作中识别关节轴，
说明不必要求传感器外壳与解剖轴精确平行~\cite{seel2014imu}。Zabat 等进一步表明，相比单一
静态姿势，多姿态信息可以减小安装偏差和软组织运动引起的角度误差~\cite{zabat2019multiple}。
近期研究还指出，待测系统与参考系统若采用不同的传感器--肢段定义，会形成明显的系统性关节角
偏差~\cite{adjel2025influence}。

现有标定方法大体可分为静态、功能性和多姿态方法。静态方法操作简单，但对参考姿势执行误差
敏感；功能性方法通过重复运动拟合关节轴，降低了对安装朝向的要求，但依赖充分且近似单轴的
动作激励；多姿态方法利用若干互补方向联合估计安装关系，并可通过各姿态残差发现错误动作或
绑带移动。对肩部三自由度、肘部屈伸和前臂轴向旋转同时估计时，单一姿态不能提供足够约束，
尤其在大臂与前臂接近共线时，肘角和轴向旋转的可观测性会降低。

Cutti 等使用惯性与磁传感器进行肩、肘动态测量的研究也表明，统一的肢段定义、安装标定和
磁环境检查是获得可解释上肢角度的共同前提~\cite{cutti2008ambulatory}。

电磁设备还存在“设备场标定”与“人体安装标定”两个不同层次。前者从已知空间位姿建立电磁
观测与接收器位姿之间的映射，并处理工作空间内的位置相关误差；相关综述与标准化评估研究
均指出，金属畸变、工作距离和空间位置会共同影响电磁跟踪误差
~\cite{kindratenko2000survey,frantz2003accuracy,hummel2006evaluation,franz2014electromagnetic,
maclachlan2025calibration}。后者在设备位姿已经
可用的前提下，估计接收器外壳与人体肢段之间的关系。本文不修改设备固件或制造商场模型，而是
针对每次佩戴执行多姿态安装级标定，并把姿态残差、位置模型残差和重复性作为进入在线估计前的
质量门控。

\section{问题描述与总体思路}
\label{sec:ch4_problem_overview}

本研究使用一套电磁位姿跟踪系统获取操作者左臂运动。电磁发射器固定于操作者躯干，两个
接收器分别固定于左大臂和左前臂。系统以发射器坐标系为共同参考，以约 $50\,\mathrm{Hz}$
的频率输出两个接收器的三维位置和姿态。与基于人体外观的视觉方法相比，电磁测量不依赖
纹理和光照，也不存在人体自身造成的直接视觉遮挡。然而，接收器在每次佩戴后的安装朝向
并不相同，绑带和软组织还会使接收器运动偏离理想的刚体关节模型。因此，原始接收器位姿
不能直接作为机器人关节角使用，必须先完成针对本次安装的标定。这里的“原始位姿”专指
设备直接报告的接收器外壳相对发射器的位置和旋转；它尚未去除外壳安装角，也尚未转换为
人体大臂、前臂方向或关节角。

本章所研究的问题可表述为：给定大臂接收器和前臂接收器在电磁发射器坐标系中的同步位姿，
估计人体左臂的五个运动量
\begin{equation}
  \boldsymbol q_h =
  \begin{bmatrix}
    q_p & q_r & q_y & q_e & q_w
  \end{bmatrix}^{\mathrm T},
  \label{eq:ch4_human_arm_state}
\end{equation}
其中，$q_p$、$q_r$ 和 $q_y$ 分别表示肩部前举/后摆、外展/内收和轴向旋转，$q_e$
表示肘部由约 $90^\circ$ 屈曲姿态向伸直方向的变化量，$q_w$ 表示前臂旋前/旋后。
上述变量沿用 Unitree G1 左臂的关节语义，但其数值首先描述人体动作，随后才在重定向和
安全层中被限制为机器人可执行目标。

需要强调的是，式~\eqref{eq:ch4_human_arm_state} 中的五个量是本研究所选运动学模型的
广义坐标，而不是由传感器直接测得的五个解剖学关节角。人体肩胛骨、锁骨、腕部屈伸以及
关节中心的平移均未被单独建模。本文的目标是构造一组与 G1 左臂关节语义一致、能够由当前
两接收器配置稳定估计的低维运动描述，而不是恢复完整的人体上肢骨骼模型。

系统采用“安装级自适应标定”思路。这里的“自适应”并非在线学习或自适应控制，而是指每次
重新佩戴设备后，系统根据本次安装采集的参考动作重新辨识：
\begin{enumerate}[label=\arabic*)]
  \item 电磁发射器坐标系到人体标准坐标系的映射；
  \item 接收器局部坐标系中的肢段纵轴；
  \item 大臂接收器位置与人体大臂方向之间的映射；
  \item 前臂旋前/旋后的局部旋转轴和零位；
  \item 肘角几何估计的增益和偏置。
\end{enumerate}

在线阶段只使用本次安装生成的参数。任一接收器或发射器被拆卸、重新绑扎或在支架中移动
后，原标定不再有效，必须重新执行标定流程。
换言之，“安装级”表示参数在一次佩戴完成后被辨识，并在该次运行期间保持固定；它不表示
参数随每一帧在线更新。这样既允许系统适应重复佩戴造成的安装差异，又避免在线更新参数时
把操作者的真实动作错误地吸收到标定模型中。

\subsection{坐标系与符号}
\label{subsec:ch4_frames_notation}

本章使用表~\ref{tab:ch4_coordinate_frames} 所列坐标系。
\begin{table}[htbp]
  \centering
  \caption{人体手臂估计使用的坐标系}
  \label{tab:ch4_coordinate_frames}
  \begin{tabularx}{\linewidth}{@{}llX@{}}
    \toprule
    符号 & 英文含义 & 物理含义 \\
    \midrule
    $\{T\}$ & Transmitter & 电磁发射器坐标系 \\
    $\{B\}$ & Body & 人体标准坐标系，与机器人约定一致：$+x$ 向前、$+y$ 向左、$+z$ 向上 \\
    $\{U\}$ & Upper arm & 大臂接收器局部坐标系 \\
    $\{F\}$ & Forearm & 前臂接收器局部坐标系 \\
    \bottomrule
  \end{tabularx}
\end{table}

记 ${}^{T}\boldsymbol p_U$ 和 ${}^{T}\boldsymbol p_F$ 为两个接收器在 $\{T\}$ 中的位置，
${}^{T}\boldsymbol R_U$ 和 ${}^{T}\boldsymbol R_F$ 为将接收器局部向量旋转到 $\{T\}$ 的
旋转矩阵。上标表示量所在的表达坐标系，下标表示被描述的对象；例如
${}^{T}\boldsymbol p_U$ 表示“大臂接收器的位置，用发射器坐标系表达”。
${}^{B}\boldsymbol u$ 和 ${}^{B}\boldsymbol f$ 分别表示人体大臂和前臂在
$\{B\}$ 中的单位方向。本文所称“肢段”是由相邻关节界定、在当前模型中近似视为刚体的
身体部分：大臂肢段连接肩与肘，前臂肢段连接肘与腕。所谓单位方向，是指只保留肢段的朝向
而忽略其长度，即
\begin{equation}
  \lVert{}^{B}\boldsymbol u\rVert_2=1,
  \qquad
  \lVert{}^{B}\boldsymbol f\rVert_2=1.
  \label{eq:ch4_unit_segment_directions}
\end{equation}
将单位向量平移到 $\{B\}$ 的数学原点后，其终点位于单位球面上。向量是自由向量，平移
起点不会改变其方向；因此这里的单位球面只是统一表示所有三维方向的数学空间，并不要求
人体肩关节位于 $\{B\}$ 原点，也不表示接收器在物理上沿半径为 1 的球面运动。

为与 G1 关节模型保持一致，本研究将零位大臂方向和零位前臂方向定义为
\begin{equation}
  \boldsymbol a_{u,0}=\begin{bmatrix}0&0&-1\end{bmatrix}^{\mathrm T},
  \qquad
  \boldsymbol a_{f,0}=\begin{bmatrix}1&0&0\end{bmatrix}^{\mathrm T}.
  \label{eq:ch4_zero_segment_axes}
\end{equation}
因此，零位对应“大臂自然下垂、肘部弯曲约 $90^\circ$、前臂水平向前”的姿态。
这里的“零位方向”是运动学模型在所有广义坐标均为零时规定的参考方向，并非接收器局部
坐标轴，也不等同于医学上的
标准解剖姿势。当所有模型关节变量为零时，大臂和前臂方向分别等于
$\boldsymbol a_{u,0}$ 和 $\boldsymbol a_{f,0}$；其他姿态均由关节旋转作用于这两个参考
方向得到。图~\ref{fig:ch4_frames_zero_pose} 汇总了坐标轴、接收器和零位方向之间的关系。

\begin{figure}[htbp]
  \centering
  \fbox{\parbox{0.90\linewidth}{\centering
    示意图内容：人体左臂处于 DOWN 零位；标出 $\{B\}$ 的 $+x$ 前、$+y$ 左、$+z$ 上，
    发射器坐标系 $\{T\}$，大臂与前臂接收器坐标系 $\{U\}$、$\{F\}$，以及
    $\boldsymbol a_{u,0}=[0,0,-1]^{\mathrm T}$ 和
    $\boldsymbol a_{f,0}=[1,0,0]^{\mathrm T}$。}}
  \caption{双接收器安装、坐标系与运动学零位示意（待替换为正式绘图）}
  \label{fig:ch4_frames_zero_pose}
\end{figure}

\subsection{五自由度运动学模型}
\label{subsec:ch4_five_dof_model}

本研究用肩部三个旋转、肘部一个旋转以及前臂一个轴向旋转描述左臂运动，即
\begin{equation}
  3\ \text{(shoulder)}+1\ \text{(elbow)}+1\ \text{(forearm twist)}
  =5\ \text{DOF}.
  \label{eq:ch4_five_dof_count}
\end{equation}
肩部三个变量 $q_p$（pitch，前举/后摆）、$q_r$（roll，外展/内收）和 $q_y$（yaw，绕
大臂纵轴旋转）共同确定大臂的空间方向及肘平面朝向；$q_e$ 描述肘部由弯曲向伸直的变化，
$q_w$ 描述前臂旋前/旋后。本模型没有包含 G1 腕部的 pitch 和 yaw，也没有单独描述肩胛带
运动，因此应理解为面向动作重定向的五自由度表示，而非完整的人体上肢解剖模型。

肩部和肘部的正向运动学采用
\begin{align}
  {}^{B}\boldsymbol R_S
  &=\boldsymbol R_y(q_p)\boldsymbol R_x(q_r)\boldsymbol R_z(q_y),
  \label{eq:ch4_shoulder_rotation}\\
  {}^{B}\boldsymbol R_E
  &={}^B\boldsymbol R_S\boldsymbol R_y(q_e),
  \label{eq:ch4_forearm_rotation}
\end{align}
其中，${}^{B}\boldsymbol R_S$ 表示肩部旋转，${}^{B}\boldsymbol R_E$ 表示不包含腕部
旋转的前臂坐标系。相应的肢段方向为
\begin{equation}
  {}^{B}\boldsymbol u={}^B\boldsymbol R_S\boldsymbol a_{u,0},
  \qquad
  {}^{B}\boldsymbol f={}^B\boldsymbol R_E\boldsymbol a_{f,0}.
  \label{eq:ch4_segment_directions}
\end{equation}

式~\eqref{eq:ch4_shoulder_rotation} 的乘法顺序是模型定义的一部分，不能任意交换。本文采用
列向量和主动旋转约定，因此右侧旋转先作用：按矩阵作用顺序，零位肢段依次接受
$\boldsymbol R_z(q_y)$、$\boldsymbol R_x(q_r)$ 和 $\boldsymbol R_y(q_p)$。各轴均以本章统一的
矩阵约定解释，而不是在推导中临时更换为另一套欧拉角约定。矩阵乘法不满足交换律，改变顺序会改变同一组
角度对应的空间姿态。本顺序与工程中 G1 左肩三个关节的语义保持一致，而不是声称它是唯一的
人体解剖角定义；解剖学结果的轴序应另按 ISB 推荐报告~\cite{wu2005isb}。

该模型的优点是人体估计量与 G1 五个左臂关节具有一致的语义，避免了先恢复完整人体骨架
再进行第二次逆运动学求解。模型的限制是：仅由两个肢段方向不能在所有姿态下唯一分离肩部
轴向旋转与前臂旋转。第~\ref{sec:ch4_observability_robustness} 节将对这一可观测性问题进行
专门处理。

\subsection{测量信息与五自由度状态的对应关系}
\label{subsec:ch4_measurement_dof_correspondence}

为了说明两个接收器为何能够支持五个运动量的估计，需要区分“数据分量个数”和“独立
自由度个数”。一个三维单位方向属于单位球面
\begin{equation}
  S^2=\left\{\boldsymbol v\in\mathbb R^3\mid
  \boldsymbol v^{\mathrm T}\boldsymbol v=1\right\}.
  \label{eq:ch4_unit_sphere_definition}
\end{equation}
虽然 $\boldsymbol v$ 具有三个坐标分量，但单位长度约束消除了一个独立变量，因此一个方向
只有两个自由度。直观上，这两个自由度类似球面上的经度和纬度；它们能够确定肢段指向，
却不能确定肢段绕自身纵轴旋转了多少。

旋转矩阵属于特殊正交群
\begin{equation}
  SO(3)=\left\{\boldsymbol R\in\mathbb R^{3\times3}\mid
  \boldsymbol R^{\mathrm T}\boldsymbol R=\boldsymbol I,
  \det(\boldsymbol R)=1\right\}.
  \label{eq:ch4_so3_definition}
\end{equation}
旋转矩阵虽然包含九个元素，但正交和单位长度约束使其只有三个独立自由度。也可以将其理解
为“一个轴的空间方向”所需的两个自由度，加上“绕该轴旋转”所需的一个自由度。

“安装标定”与“在线估计”是前后相接的两个阶段。安装标定只在重新佩戴后执行，利用已知
参考姿势求出坐标对齐、局部肢段轴、位置模型和角度修正参数；在线估计则在每一个同步采样
时刻使用这些固定参数，把未知动作转换为五自由度状态，不再重新拟合参数。完成安装标定后，
在线估计实际使用两类几何信息：大臂单位方向
${}^{B}\hat{\boldsymbol u}\in S^2$，以及经过安装对齐的前臂完整姿态
${}^{B}\boldsymbol R_F\in SO(3)$。因此可用独立信息量为
\begin{equation}
  \underbrace{2}_{\text{upper-arm direction}}+
  \underbrace{3}_{\text{forearm attitude}}=5.
  \label{eq:ch4_measurement_dof_count}
\end{equation}
前臂姿态的三个自由度还可以分解为前臂方向 ${}^{B}\hat{\boldsymbol f}$ 的两个自由度和
绕前臂纵轴旋转的一个自由度。由此形成表~\ref{tab:ch4_measurement_to_joint} 所示的逐级
求解关系。

\begin{table}[htbp]
  \centering
  \caption{测量几何量与五自由度状态的对应关系}
  \label{tab:ch4_measurement_to_joint}
  \begin{tabularx}{\linewidth}{@{}lXX@{}}
    \toprule
    待估计量 & 主要几何依据 & 物理含义 \\
    \midrule
    $q_p,q_r$ & 大臂单位方向 ${}^{B}\hat{\boldsymbol u}$ & 大臂指向的两个自由度 \\
    $q_e$ & ${}^{B}\hat{\boldsymbol u}$ 与 ${}^{B}\hat{\boldsymbol f}$ 的夹角
      & 肘部由弯曲向伸直的变化量 \\
    $q_y$ & 前臂方向在大臂法平面内的投影 & 肘平面绕大臂纵轴的朝向 \\
    $q_w$ & 去除肩、肘和固定安装旋转后的前臂残余姿态
      & 绕前臂纵轴的旋前/旋后 \\
    \bottomrule
  \end{tabularx}
\end{table}

上述自由度计数只说明一般姿态下信息量足够，并不保证所有姿态都存在唯一解。当大臂正侧举
时，pitch 方向退化；当肘部完全伸直时，肩部轴向旋转与前臂旋转发生耦合。换言之，该测量
映射在一般姿态下局部可辨识，但在特定几何构型下秩会下降。系统对这些区域的处理将在
第~\ref{sec:ch4_observability_robustness} 节给出。

\section{双接收器同步采集与数据预处理}
\label{sec:ch4_synchronized_acquisition}

\subsection{同步位姿对}
\label{subsec:ch4_synchronized_pose_pairs}

两个 WEMT 接收器分别通过网络返回独立的数据流。如果直接使用两条消息各自最近的样本，
快速运动时可能把不同时刻的大臂和前臂姿态组合在一起，从而产生虚假的肘角和肩部旋转。
因此，采集节点只在以下条件同时满足时发布一对位姿：
\begin{enumerate}[label=\arabic*)]
  \item 两个接收器均处于连接状态且没有设备错误；
  \item 两个接收器都产生了新的硬件帧；
  \item 任一位姿的接收时间不超过 $100\,\mathrm{ms}$；
  \item 两个接收器的到达时间差不超过 $30\,\mathrm{ms}$；
  \item 两个原始四元数均通过有效性检查。
\end{enumerate}

通过检查后，系统给两个 \texttt{geometry\_msgs/PoseStamped} 消息赋予同一个 ROS 时间戳。
该消息由三维位置、四元数姿态、参考坐标系标识和时间戳组成。两条消息分别发布至 ROS 话题
\texttt{/em/upper\_arm\_pose} 和 \texttt{/em/forearm\_pose}。这样，后续标定和在线估计
可按相同时间戳进行严格配对。WEMT 接口内部的平滑窗口被设为 1，即发布未经时间平均的
单帧样本；这里的 1 是窗口中包含的样本数，而不是 1 秒。该设置避免位置和姿态使用不同
滤波造成额外的相位差；统一滤波在完成双流配对后进行。

\subsection{四元数有效性与旋转矩阵}
\label{subsec:ch4_quaternion_validation}

设设备返回四元数为
\begin{equation}
  \boldsymbol q=\begin{bmatrix}q_w&q_x&q_y&q_z\end{bmatrix}^{\mathrm T}.
  \label{eq:ch4_quaternion}
\end{equation}
系统首先检查其模长是否满足
\begin{equation}
  0.95\leq \lVert\boldsymbol q\rVert_2\leq1.05.
  \label{eq:ch4_quaternion_norm_gate}
\end{equation}

未通过检查的数据对被直接丢弃，而不是先强制归一化后继续使用。这是因为持续出现约
0.707 等异常模长时，问题通常来自设备原始数据失效；单纯归一化会掩盖错误并把无效姿态
传入控制链。有效四元数归一化后转换为旋转矩阵：
\begin{equation}
\boldsymbol R(\boldsymbol q)=
\begin{bmatrix}
1-2(q_y^2+q_z^2) & 2(q_xq_y-q_zq_w) & 2(q_xq_z+q_yq_w)\\
2(q_xq_y+q_zq_w) & 1-2(q_x^2+q_z^2) & 2(q_yq_z-q_xq_w)\\
2(q_xq_z-q_yq_w) & 2(q_yq_z+q_xq_w) & 1-2(q_x^2+q_y^2)
\end{bmatrix}.
\label{eq:ch4_quaternion_to_rotation}
\end{equation}

由于当前设备接口只明确四元数分量顺序，而接收器到发射器的主动/被动旋转方向仍需由系统
侧验证，系统在标定时
分别评价原旋转矩阵 $\boldsymbol R$ 和转置矩阵 $\boldsymbol R^{\mathrm T}$。选择标准以前臂
旋转轴在多参考姿势下的方向残差为主，以相对旋转残差作为分数接近时的辅助判据。最终选定
的约定被写入当次标定文件，并在在线估计中保持一致。这一步不是任意改变四元数分量，而是
在互为逆旋转的两种物理解释之间进行一次安装级一致性检验；若后续由设备协议明确方向，仍可
保留该检验作为接线或固件变化的诊断。

\section{安装级自适应标定}
\label{sec:ch4_installation_calibration}

\subsection{标定流程概览}
\label{subsec:ch4_calibration_overview}

完整标定由三个前置质量/动作检查和七个静态参考姿势组成，如图~\ref{fig:ch4_calibration_flow}
所示。
\begin{figure}[htbp]
  \centering
  \fbox{\parbox{0.88\linewidth}{\centering
    图示占位：静态质量检查 $\rightarrow$ 接收器身份检查 $\rightarrow$
    前臂旋转轴标定 $\rightarrow$ 七姿势采集 $\rightarrow$
    参数拟合 $\rightarrow$ 质量门槛 $\rightarrow$ 保存或拒绝。}}
  \caption{安装级自适应标定流程（待替换为正式绘图）}
  \label{fig:ch4_calibration_flow}
\end{figure}

标定开始前，操作者保持自然站立，系统采集 $4\,\mathrm{s}$ 静态数据，用于计算位置和姿态
噪声。随后，操作者在保持大臂基本不动的情况下反复屈伸肘部并转动前臂。若前臂接收器的
姿态变化明显大于大臂接收器，才认为两个接收器的身份与安装位置正确。最后，操作者保持大臂
竖直、肘部弯曲约 $90^\circ$，仅使前臂绕“肘到腕”的纵轴做旋前/旋后，用于辨识前臂接收器
局部坐标系中的旋转轴。

七个静态参考姿势及其理论角度见表~\ref{tab:ch4_reference_poses}。每个姿势离开后重新构建
3 次，每次保持 $2\,\mathrm{s}$。这种“离开---重建---保持”的方式能够发现单次摆姿偶然
正确但重复性不足的问题。

\begin{table}[htbp]
  \centering
  \caption{安装级标定使用的静态参考姿势}
  \label{tab:ch4_reference_poses}
  \begin{tabularx}{\linewidth}{@{}XrrrrX@{}}
    \toprule
    参考姿势 & $q_p$ & $q_r$ & $q_y$ & $q_e$ & 主要作用 \\
    \midrule
    DOWN             & $0^\circ$   & $0^\circ$  & $0^\circ$ & $0^\circ$  & 定义系统零位 \\
    FORWARD\_30      & $-30^\circ$ & $0^\circ$  & $0^\circ$ & $0^\circ$  & 约束大臂位置模型 \\
    FORWARD\_60      & $-60^\circ$ & $0^\circ$  & $0^\circ$ & $0^\circ$  & 约束大臂位置模型 \\
    FORWARD          & $-90^\circ$ & $0^\circ$  & $0^\circ$ & $0^\circ$  & 约束前举方向 \\
    SIDEWAYS         & $0^\circ$   & $90^\circ$ & $0^\circ$ & $90^\circ$ & 约束侧举和伸肘方向 \\
    STRAIGHT\_DOWN   & $0^\circ$   & $0^\circ$  & $0^\circ$ & $90^\circ$ & 约束伸肘姿态 \\
    STRAIGHT\_FORWARD& $-90^\circ$ & $0^\circ$  & $0^\circ$ & $90^\circ$ & 约束主要工作空间 \\
    \bottomrule
  \end{tabularx}
\end{table}

图~\ref{fig:ch4_reference_pose_set} 给出七个参考姿势应呈现的肢段关系。FORWARD\_30 和
FORWARD\_60 的角度不要求操作者用量角器达到绝对精确值；其作用是为前举路径增加两个
中间方向。操作者应优先保证躯干不动、动作平面一致、肘部状态正确以及三次重复可重现，实际
偏差将反映在参考残差和位置模型残差中。

\begin{figure}[htbp]
  \centering
  \fbox{\parbox{0.92\linewidth}{\centering
    示意图内容：依次绘制 DOWN、FORWARD\_30、FORWARD\_60、FORWARD、SIDEWAYS、
    STRAIGHT\_DOWN 和 STRAIGHT\_FORWARD；用实线表示大臂、虚线表示前臂，并标注
    “弯肘约 $90^\circ$”或“伸肘”以及前举/侧举方向。}}
  \caption{七个安装标定参考姿势及其肢段约束（待替换为实拍图或正式绘图）}
  \label{fig:ch4_reference_pose_set}
\end{figure}

FORWARD\_30 和 FORWARD\_60 只用于提高大臂位置模型在前举路径上的可辨识性；前臂方向
拟合使用 DOWN、FORWARD、SIDEWAYS、STRAIGHT\_DOWN 和 STRAIGHT\_FORWARD 五类姿势。

\subsection{静态质量和重复性评价}
\label{subsec:ch4_static_quality_repeatability}

对于一组静态位置样本 $\boldsymbol p_i$，位置均方根波动定义为
\begin{equation}
  e_{p,\mathrm{RMS}}=
  \sqrt{\frac{1}{N}\sum_{i=1}^{N}
  \lVert\boldsymbol p_i-\bar{\boldsymbol p}\rVert_2^2}.
  \label{eq:ch4_static_position_rms}
\end{equation}

姿态中心采用旋转矩阵均值在 $SO(3)$ 上的最近旋转，姿态波动则用每个样本到中心旋转的
测地角计算 RMS。由于旋转不能像普通欧氏向量那样逐元素平均后直接使用，最近旋转投影和
四元数/旋转平均是姿态处理中常用的几何做法~\cite{hartley2013rotation,markley2007averaging}。
当前工程实现中，当任一接收器的位置 RMS 超过 $3\,\mathrm{mm}$，或姿态
RMS 超过 $3^\circ$ 时，系统提示检查绑带、金属物体和跟踪距离。

每个参考姿势的三次重复使用旋转 medoid（样本中心）作为代表，即只在三个实际观测中选择
一个，使它到其余观测的旋转测地距离总和最小。与直接平均相比，medoid 必然是一个真实采集
样本，对单次离群姿态更稳健。若任一重复与代表样本之间的最大角距离超过 $15^\circ$，则只重新采集当前
姿势，已经通过的姿势继续保留。该策略避免一次摆姿错误使整轮标定全部作废，同时保证每个
参考姿势具有最低重复性。

\subsection{前臂局部旋转轴辨识}
\label{subsec:ch4_forearm_twist_axis}

在理想旋前/旋后动作中，前臂在发射器坐标系中的纵轴方向保持近似不变，而接收器绕该轴
转动。实际只有一个前臂接收器，下标 $i$ 表示该接收器在第 $i$ 个采样时刻的观测。设其旋转
矩阵为 ${}^{T}\boldsymbol R_{F,i}$，希望找到接收器外壳坐标系中的局部单位向量
$\boldsymbol a_t$，使
\begin{equation}
  {}^{T}\boldsymbol R_{F,i}\boldsymbol a_t\approx\boldsymbol d_t,
  \qquad i=1,\ldots,N,
  \label{eq:ch4_twist_axis_invariance}
\end{equation}
其中，$\boldsymbol d_t$ 为发射器坐标系中的稳定旋转轴。寻找 $\boldsymbol a_t$ 的目的，是把
任意安装朝向下的“接收器外壳某一轴”辨识为人体前臂纵轴；若轴选择正确，接收器即使绕它
转动，该轴变换到发射器坐标系后的方向仍应近似不变。令
\begin{equation}
  \bar{\boldsymbol R}_F=\frac{1}{N}
  \sum_{i=1}^{N}{}^{T}\boldsymbol R_{F,i},
  \label{eq:ch4_mean_forearm_rotation}
\end{equation}
则 $\boldsymbol a_t$ 可写为
\begin{equation}
  \boldsymbol a_t=
  \arg\max_{\lVert\boldsymbol a\rVert=1}
  \lVert\bar{\boldsymbol R}_F\boldsymbol a\rVert_2,
  \label{eq:ch4_twist_axis_svd}
\end{equation}
即 $\bar{\boldsymbol R}_F$ 最大奇异值对应的右奇异向量。式
~\eqref{eq:ch4_twist_axis_invariance} 给出物理不变性，式
~\eqref{eq:ch4_mean_forearm_rotation} 将全部姿态观测聚合，式
~\eqref{eq:ch4_twist_axis_svd} 则寻找经这些旋转后最不发散的局部方向；最大奇异值对应方向
正是保持合向量长度最大的方向。系统对残差进行基于中位数绝对
偏差的迭代剔除，降低偶发姿态跳变的影响；同时通过样本对之间的相对旋转检查动作范围。
只有当前臂旋转范围至少达到 $55^\circ$、大臂姿态变化不超过 $15^\circ$，且旋转轴拟合
RMS 不超过 $15^\circ$ 时，才继续执行后续标定。

该旋转轴专门服务于腕部旋前/旋后估计。它不与前臂“肘到腕”的方向轴强制共用，因为皮肤
滑动、绑带安装和参考动作误差会使两者的最优拟合结果存在差异。实际实现分别辨识前臂方向
轴与旋转轴，以避免一个轴同时承担两个不同物理任务造成系统性偏差。

\subsection{大臂方向的仿射位置--方向模型}
\label{subsec:ch4_upper_arm_affine_model}

大臂接收器固定在软组织表面，其姿态容易受到皮肤相对肱骨转动的影响。实验开发过程中发现，
直接使用大臂接收器姿态估计肱骨轴向运动会引入较大误差。因此，当前混合模型主要利用大臂
接收器的位置恢复大臂方向。这里的目标不是根据位置直接回归关节角，而是先恢复具有明确
几何意义的单位方向 ${}^{B}\hat{\boldsymbol u}$，再由第~\ref{sec:ch4_online_arm_estimation}
节的逆运动学关系求解关节变量。

若接收器与肩关节构成理想刚体，大臂运动时接收器应在以肩关节为中心的球面上运动。但实际
轨迹受绑带、软组织和发射器安装影响，前举路径往往呈现非球形畸变。为此，本研究采用仿射
位置--方向模型
\begin{equation}
  {}^{T}\boldsymbol p_{U,i}
  =\boldsymbol c+\boldsymbol M\,{}^{B}\boldsymbol u_i^{*}
  +\boldsymbol\varepsilon_i,
  \label{eq:ch4_upper_affine_model}
\end{equation}
其中，$\boldsymbol c\in\mathbb R^3$ 为本次安装下的轨迹中心，
$\boldsymbol M\in\mathbb R^{3\times3}$ 为从人体标准方向到发射器坐标位移的可逆仿射
矩阵，${}^{B}\boldsymbol u_i^{*}$ 是第 $i$ 个参考姿势的理论大臂方向，
$\boldsymbol\varepsilon_i$ 汇集摆姿误差、软组织运动和测量噪声。$\boldsymbol c$ 是用于
位置--方向映射的虚拟中心，不应直接解释为解剖学肩关节中心。

这里称为“混合模型”，是因为大臂方向主要由接收器\emph{位置}恢复，而前臂方向和轴向
旋转由接收器\emph{姿态}恢复。所谓“恢复大臂方向”，是指由测得的位置
${}^{T}\boldsymbol p_U$ 反求单位向量 ${}^{B}\hat{\boldsymbol u}$，而不是估计大臂长度或
肩关节中心。图~\ref{fig:ch4_affine_geometry} 给出式~\eqref{eq:ch4_upper_affine_model} 的
几何含义。

\begin{figure}[htbp]
  \centering
  \fbox{\parbox{0.92\linewidth}{\centering
    示意图内容：以 $\boldsymbol c$ 为椭球中心，单位方向
    ${}^{B}\boldsymbol u_i^{*}$ 经过 $\boldsymbol M$ 映射为从 $\boldsymbol c$ 指向
    ${}^{T}\boldsymbol p_{U,i}$ 的位移；标出实测点、重投影点 $\hat{\boldsymbol p}_U$、
    残差 $\boldsymbol\varepsilon_i$ 与模型误差 $e_M$。左侧可并列绘制单位球，右侧绘制其
    经过 $\boldsymbol M$ 后的椭球像。}}
  \caption{大臂仿射位置--方向模型的几何解释（待替换为正式绘图）}
  \label{fig:ch4_affine_geometry}
\end{figure}

该模型被称为“仿射椭球”，是因为单位方向满足
$({}^{B}\boldsymbol u)^{\mathrm T}{}^{B}\boldsymbol u=1$。忽略噪声并令
$\boldsymbol x={}^{T}\boldsymbol p_U-\boldsymbol c$，由
$\boldsymbol x=\boldsymbol M{}^{B}\boldsymbol u$ 可得
\begin{equation}
  \boldsymbol x^{\mathrm T}\boldsymbol M^{-\mathrm T}
  \boldsymbol M^{-1}\boldsymbol x=1,
  \label{eq:ch4_affine_ellipsoid_implicit}
\end{equation}
这里引入 $\boldsymbol x$ 只是为了简化表达：它是从虚拟中心 $\boldsymbol c$ 指向当前接收器
位置的位移，不是新的测量量。将 ${}^{B}\boldsymbol u=\boldsymbol M^{-1}\boldsymbol x$
代回单位长度约束即可得到式~\eqref{eq:ch4_affine_ellipsoid_implicit}，这正是以
$\boldsymbol c$ 为中心的椭球面二次型。若对矩阵作奇异值分解
\begin{equation}
  \boldsymbol M=\boldsymbol U\boldsymbol\Sigma\boldsymbol V^{\mathrm T},
  \qquad
  \boldsymbol\Sigma=\operatorname{diag}(\sigma_1,\sigma_2,\sigma_3),
  \label{eq:ch4_affine_matrix_svd}
\end{equation}
则 $\boldsymbol U$ 给出椭球主轴方向，$\sigma_1,\sigma_2,\sigma_3$ 给出三个半轴长度。
当三个奇异值相等时，说明三个主方向的尺度相同；若
$\boldsymbol M=r\boldsymbol R$，线性变换只包含统一尺度 $r$ 和旋转 $\boldsymbol R$，模型便
退化为半径为 $r$ 的球面。若奇异值不同，单位球在三个主方向上被拉伸为不同半轴长度，形成
椭球。因此，椭球是单位方向集合经过本次安装相关线性变换后的位置轨迹，而不是对人体肩关节形状的
解剖学假设。实际七个参考姿势只覆盖该椭球面的有限区域，本模型应理解为工作空间内的经验
标定模型，而非完整肩部轨迹的全局物理模型。

所有参考样本通过最小二乘求解
\begin{equation}
  (\hat{\boldsymbol c},\hat{\boldsymbol M})=
  \arg\min_{\boldsymbol c,\boldsymbol M}
  \sum_i\left\|
  {}^{T}\boldsymbol p_{U,i}-\boldsymbol c-
  \boldsymbol M{}^{B}\boldsymbol u_i^{*}
  \right\|_2^2.
  \label{eq:ch4_upper_affine_fit}
\end{equation}
式~\eqref{eq:ch4_upper_affine_fit} 对 $\boldsymbol c$ 和 $\boldsymbol M$ 是线性最小二乘
问题。它求出的 $\hat{\boldsymbol c}$ 和 $\hat{\boldsymbol M}$ 用于在线阶段把新的接收器
位置反算为大臂方向。为了把常量中心项和三维方向项写在同一个线性方程中，将每个理论方向
在末尾附加常数 1，形成增广向量
$[({}^{B}\boldsymbol u_i^{*})^{\mathrm T},1]$。相应设计矩阵必须具有满列秩 4，即四个列方向
都具有独立信息；否则轨迹
中心和全部仿射分量不能被唯一辨识。DOWN、FORWARD 和 SIDEWAYS 三类方向不足以满足这一
条件，因此标定流程额外引入 FORWARD\_30 和 FORWARD\_60，以增加前举路径上的方向变化
并改善模型条件。

在线阶段，根据当前大臂接收器位置计算未归一化方向
\begin{equation}
  \tilde{\boldsymbol u}=
  \hat{\boldsymbol M}^{-1}
  ({}^{T}\boldsymbol p_U-\hat{\boldsymbol c}),
  \label{eq:ch4_upper_direction_raw}
\end{equation}
并归一化得到
\begin{equation}
  {}^{B}\hat{\boldsymbol u}=
  \frac{\tilde{\boldsymbol u}}
  {\lVert\tilde{\boldsymbol u}\rVert_2}.
  \label{eq:ch4_upper_direction_normalized}
\end{equation}
若当前点严格位于标定椭球面上，则
$\lVert\tilde{\boldsymbol u}\rVert_2=1$；实际数据受噪声影响通常偏离曲面，因此仍需归一化。
这一过程可以理解为先用逆仿射变换消除安装相关的方向性畸变，再将结果投影回单位方向空间。

矩阵 $\hat{\boldsymbol M}$ 的奇异值表示沿三个主方向的位移尺度。最小奇异值过小时，说明
某一方向几乎没有被参考姿势激励，求逆时该方向的微小位置噪声会被显著放大。为此定义条件数
\begin{equation}
  \kappa(\hat{\boldsymbol M})=
  \frac{\sigma_{\max}(\hat{\boldsymbol M})}
       {\sigma_{\min}(\hat{\boldsymbol M})}.
  \label{eq:ch4_affine_condition_number}
\end{equation}
条件数越接近 1，三个方向的尺度越均衡；条件数过大则表明参考方向分布不足、摆姿不一致或
拟合接近奇异。当前标定质量门槛要求 $\kappa(\hat{\boldsymbol M})\leq4.0$。

为了检测设备移动或当前动作严重超出标定范围，系统将估计方向重新投影到椭球表面：
\begin{equation}
  \hat{\boldsymbol p}_U=
  \hat{\boldsymbol c}+
  \hat{\boldsymbol M}{}^{B}\hat{\boldsymbol u},
  \label{eq:ch4_upper_position_reprojection}
\end{equation}
并定义位置模型误差
\begin{equation}
  e_M=\left\|{}^{T}\boldsymbol p_U-
  \hat{\boldsymbol p}_U\right\|_2.
  \label{eq:ch4_upper_model_error}
\end{equation}
在线运行时，若 $e_M>50\,\mathrm{mm}$，当前帧被保护性丢弃。该误差不等同于人体大臂
方向的绝对真值误差，而是接收器位置与本次安装模型的一致性指标。小的 $e_M$ 只能说明
当前数据与标定轨迹相容；要评价大臂方向的绝对精度，仍需要独立的光学动作捕捉或其他真值
测量系统。

\subsection{前臂方向与坐标对齐}
\label{subsec:ch4_forearm_direction_alignment}

设 $\boldsymbol a_d$ 为前臂接收器局部坐标系中的“肘到腕”方向，$\boldsymbol A_d$ 为发射器
坐标系到人体标准坐标系的旋转。对于第 $i$ 个参考姿势，预测方向为
\begin{equation}
  {}^{B}\hat{\boldsymbol f}_i=
  \boldsymbol A_d\,{}^{T}\boldsymbol R_{F,i}\boldsymbol a_d.
  \label{eq:ch4_forearm_direction_prediction}
\end{equation}

式~\eqref{eq:ch4_forearm_direction_prediction} 表示两步坐标变换：先用
${}^{T}\boldsymbol R_{F,i}$ 把接收器内部的候选方向 $\boldsymbol a_d$ 旋转到发射器坐标系，
再用 $\boldsymbol A_d$ 将其统一表达在人体坐标系。方向轴和对齐矩阵通过交替优化求解：固定
$\boldsymbol a_d$ 时，$\boldsymbol A_d$ 可归结为
Wahba 旋转对齐问题，即寻找一个旋转，使多组测量方向与对应理论方向的平方误差之和最小
~\cite{wahba1965least}；固定 $\boldsymbol A_d$ 时，可将理论方向反变换到各接收器局部坐标系
并取均值，更新 $\boldsymbol a_d$。算法从六个坐标轴正负方向初始化，保留参考姿势方向 RMS
最小的解。

旋转对齐子问题为
\begin{equation}
  \hat{\boldsymbol A}_d=
  \arg\min_{\boldsymbol A\in SO(3)}
  \sum_i\left\|
  \boldsymbol A\,{}^{T}\boldsymbol R_{F,i}\boldsymbol a_d-
  {}^{B}\boldsymbol f_i^{*}
  \right\|_2^2,
  \label{eq:ch4_forearm_wahba_problem}
\end{equation}
通过互协方差矩阵的奇异值分解求得，并显式约束 $\det(\boldsymbol A)=1$。在线阶段的前臂
单位方向为
\begin{equation}
  {}^{B}\hat{\boldsymbol f}=
  \frac{
    \hat{\boldsymbol A}_d\,{}^{T}\boldsymbol R_F\hat{\boldsymbol a}_d
  }{
    \left\|
      \hat{\boldsymbol A}_d\,{}^{T}\boldsymbol R_F\hat{\boldsymbol a}_d
    \right\|_2
  }.
  \label{eq:ch4_forearm_direction_online}
\end{equation}

除了方向映射，系统还建立独立的姿态对齐矩阵 $\boldsymbol A_t$ 和固定安装旋转
$\boldsymbol R_m$，用于计算 wrist roll。这里的“姿态对齐”是把接收器报告的完整三维旋转
转换到人体模型坐标系；$\boldsymbol R_m$ 则记录接收器外壳相对理论前臂坐标系的固定安装角。
在 DOWN 姿势下，理论前臂姿态作为腕部零位；三次 DOWN 重复的平均旋转确定本次安装的固定
旋转。在线计算时先去除肩、肘造成的理论旋转，再去除 $\boldsymbol R_m$，剩余的绕前臂纵轴
分量才被解释为腕部扭转。由此，外壳安装角和前臂方向误差不会直接被误认为 wrist roll。

\subsection{肘角的安装级仿射修正}
\label{subsec:ch4_elbow_affine_correction}

由两个单位方向可以得到几何肘角
\begin{equation}
  q_{e,g}=\arcsin\left(
  \operatorname{clip}
  ( {}^{B}\hat{\boldsymbol u}^{\mathrm T}
  {}^{B}\hat{\boldsymbol f},-1,1)
  \right).
  \label{eq:ch4_geometric_elbow_angle}
\end{equation}

下标 $g$ 表示 geometric，即由两条肢段方向直接得到的几何角。在本研究的零位约定中，
弯肘约 $90^\circ$ 时两个方向近似正交，因而 $q_{e,g}\approx0$；
手臂伸直时二者近似平行，因而 $q_{e,g}\approx90^\circ$。由于两个方向拟合各自存在固定
残差，几何角可能同时存在零位偏差和幅值缩放误差：固定零偏只能整体平移角度，不能改变从
弯曲端到伸直端的量程；单独改变增益又不能校正零位。因此本研究对每次安装分别拟合一维
仿射修正，即“缩放加平移”模型
\begin{equation}
  q_e=\operatorname{clip}
  (s_e q_{e,g}+b_e,\,0,\,\pi/2),
  \label{eq:ch4_elbow_affine_correction}
\end{equation}
其中增益 $s_e$ 和偏置 $b_e$ 由五类前臂参考姿势的理论肘角和几何肘角进行线性最小二乘
拟合。为拒绝异常标定，当前实现要求 $0.75\leq s_e\leq1.50$，且
$|b_e|\leq20^\circ$。若拟合不可辨识或超出合理范围，系统保留未经修正的几何肘角，同时
输出警告。

\subsection{标定质量门槛与实机资格}
\label{subsec:ch4_calibration_quality_gates}

标定参数并非只要能够求解就可进入真机。系统在保存前检查表~\ref{tab:ch4_quality_gates}
所列工程门槛。
\begin{table}[htbp]
  \centering
  \caption{安装级标定质量门槛}
  \label{tab:ch4_quality_gates}
  \begin{tabularx}{\linewidth}{@{}XllX@{}}
    \toprule
    检查项 & 建议/诊断阈值 & 硬限制 & 不通过时的处理 \\
    \midrule
    参考姿势重复差异 & --- & $15^\circ$ & 重采当前姿势 \\
    大臂方向参考误差 & $10^\circ$ & $15^\circ$ & 警告/拒绝保存 \\
    前臂方向参考误差 & $10^\circ$；真机资格 $20^\circ$ & $25^\circ$ & 警告、仅诊断或拒绝 \\
    前臂腕部零位重复误差 & --- & $15^\circ$ & 拒绝保存 \\
    仿射矩阵条件数 & 越接近 1 越好 & 4.0 & 拒绝保存 \\
    仿射位置模型 RMS/最大误差 & $30/50\,\mathrm{mm}$ & $35/65\,\mathrm{mm}$ & 警告/拒绝保存 \\
    \bottomrule
  \end{tabularx}
\end{table}

表中的“建议/诊断阈值”用于提示边缘质量，“硬限制”决定标定文件能否保存，两者不应简单
理解为“仿真门槛”和“实机门槛”。仿真与 dry-run 仍可使用通过硬限制但质量边缘的标定；
真机资格另由更严格的前臂方向条件控制。若前臂方向误差超过 $20^\circ$ 但仍未超过硬限制
$25^\circ$，标定可以作为诊断资料保存，
但会被标记为 \texttt{dry\_run\_only: true}。该标定只允许用于 dry-run 和离线诊断，
重定向节点拒绝向真机主题发布关节目标。该设计把“算法能够产生数值”和“标定具备实机
资格”明确分开。

这些阈值是本项目的工程安全标准，并非电磁跟踪器或人形机器人领域的通用精度标准。参考
姿势残差描述标定模型与人工构建姿势之间的一致性，不能替代独立视觉系统测得的人体姿态
真值误差。

\section{在线人体手臂姿态估计}
\label{sec:ch4_online_arm_estimation}

安装标定完成后，每一对同步位姿首先被转换为三个中间几何量：大臂单位方向
${}^{B}\hat{\boldsymbol u}$、前臂单位方向 ${}^{B}\hat{\boldsymbol f}$，以及经过腕部姿态
对齐的前臂接收器旋转 ${}^{B}\boldsymbol R_F$。其中，大臂方向由
式~\eqref{eq:ch4_upper_direction_raw}--\eqref{eq:ch4_upper_direction_normalized} 对当前大臂
接收器位置进行逆仿射映射得到；前臂方向由式~\eqref{eq:ch4_forearm_direction_online} 对当前
前臂接收器姿态和已标定局部轴进行变换得到；“腕部姿态对齐”则指用 $\boldsymbol A_t$ 将完整
接收器旋转转换到人体模型坐标系。五个关节量并不是同时进行黑箱优化，而是
按照几何依赖关系逐级求解：先由大臂方向确定 $q_p,q_r$，再由两条肢段方向确定几何肘角
$q_{e,g}$ 和肩部轴向角 $q_y$，最后从前臂残余旋转中提取 $q_w$。几何肘角随后经过本次
安装的增益和偏置修正，形成最终输出 $q_e$。

这种求解顺序保留了每个变量的物理意义，也使不可观测条件能够被显式检测。它与输入信息的
对应关系可概括为
\begin{equation}
  ({}^{T}\boldsymbol p_U,{}^{T}\boldsymbol R_F)
  \longrightarrow
  ({}^{B}\hat{\boldsymbol u},{}^{B}\hat{\boldsymbol f},
   {}^{B}\boldsymbol R_F)
  \longrightarrow
  (q_p,q_r,q_y,q_{e,g},q_w)
  \longrightarrow
  (q_p,q_r,q_y,q_e,q_w).
  \label{eq:ch4_online_estimation_chain}
\end{equation}

\subsection{肩部 pitch 与 roll}
\label{subsec:ch4_shoulder_pitch_roll}

得到大臂单位方向
\begin{equation}
  {}^{B}\hat{\boldsymbol u}=
  \begin{bmatrix}u_x&u_y&u_z\end{bmatrix}^{\mathrm T}
  \label{eq:ch4_upper_direction_components}
\end{equation}
后，由于 $\boldsymbol R_z(q_y)$ 绕零位大臂轴旋转，不改变
$\boldsymbol a_{u,0}=[0,0,-1]^{\mathrm T}$，大臂方向只与 $q_p,q_r$ 有关。展开
式~\eqref{eq:ch4_segment_directions} 可得
\begin{equation}
  {}^{B}\hat{\boldsymbol u}=
  \begin{bmatrix}
    -\sin q_p\cos q_r\\
    \sin q_r\\
    -\cos q_p\cos q_r
  \end{bmatrix}.
  \label{eq:ch4_upper_direction_expanded}
\end{equation}
这里的“零位大臂轴”就是式~\eqref{eq:ch4_zero_segment_axes} 定义的
$\boldsymbol a_{u,0}$。由于该向量位于 $z$ 轴上，绕 $z$ 轴旋转不会改变它，即
$\boldsymbol R_z(q_y)\boldsymbol a_{u,0}=\boldsymbol a_{u,0}$；随后依次左乘
$\boldsymbol R_x(q_r)$ 和 $\boldsymbol R_y(q_p)$，直接进行矩阵乘法即可得到
式~\eqref{eq:ch4_upper_direction_expanded}。

因此，肩部 pitch 和 roll 的主值可写为
\begin{equation}
  q_p=\operatorname{atan2}(-u_x,-u_z),
  \qquad
  q_r=\arcsin(\operatorname{clip}(u_y,-1,1)).
  \label{eq:ch4_pitch_roll_inverse}
\end{equation}

“主值”表示反三角函数在约定区间内返回的代表解，因为同一方向可对应相差 $2\pi$ 的周期
等价角。由式~\eqref{eq:ch4_upper_direction_expanded} 有
$u_y=\sin q_r$，故 $q_r$ 可由反正弦得到；而 $-u_x$ 与 $-u_z$ 分别包含
$\sin q_p\cos q_r$ 和 $\cos q_p\cos q_r$，用双参数反正切可同时保留象限信息并消去共同
因子 $\cos q_r$。$\operatorname{clip}(u_y,-1,1)$ 只用于消除浮点误差造成的轻微越界，
例如数值上出现 $1+10^{-12}$ 时避免 $\arcsin$ 失效，并非裁剪正常人体动作。

由于 \texttt{atan2} 的输出位于 $[-\pi,\pi]$，系统将当前结果展开到最接近上一帧的等价角，
避免跨越 $\pm\pi$ 时产生整周跳变。

当大臂接近正侧方时，$u_x$ 和 $u_z$ 同时接近零，pitch 对测量噪声高度敏感。定义
\begin{equation}
  o_p=\sqrt{u_x^2+u_z^2}=|\cos q_r|
  \label{eq:ch4_pitch_observability}
\end{equation}
为 pitch 可观测度。在线算法采用带滞回的两状态机制：当 $o_p\leq0.50$ 时进入侧举奇异区，
将 pitch 切换到统一的标准解 $q_p=0$；只有当手臂返回水平面以下且 $o_p\geq0.65$ 时才
重新采用测量解。与此同时，roll 使用标准 pitch 坐标系下的 \texttt{atan2} 表达，使其
能够连续通过 $90^\circ$，而不是被 \texttt{asin} 折回到主值范围。该处理保证从 DOWN
直接侧举和从 FORWARD 转向侧举时，同一最终人体姿势映射到同一机器人目标。

由式~\eqref{eq:ch4_upper_direction_expanded} 可见，$u_x$ 和 $u_z$ 都含有因子
$\cos q_r$；当侧举接近 $|q_r|=90^\circ$ 时该因子趋近零，pitch 的分子和分母同时退化。
阈值 0.50 和 0.65 分别约对应 $|q_r|=60^\circ$ 和 $49^\circ$，是当前工程为避免临界点
抖动设置的进入/退出滞回阈值，而不是人体生理常数。进入奇异区后固定 $q_p=0$，相当于从
无穷多个等价 pitch--roll 表示中选择统一标准解；退出阈值更高且要求回到水平面以下，使状态
不会在噪声作用下频繁切换。

\subsection{肘部伸展量}
\label{subsec:ch4_online_elbow_angle}

在运动学模型中，大臂和前臂方向分别为
${}^{B}\hat{\boldsymbol u}={}^B\boldsymbol R_S\boldsymbol a_{u,0}$ 和
${}^{B}\hat{\boldsymbol f}={}^B\boldsymbol R_S
\boldsymbol R_y(q_{e,g})\boldsymbol a_{f,0}$。利用旋转矩阵保持内积的性质，有
\begin{align}
  ({}^{B}\hat{\boldsymbol u})^{\mathrm T}
  {}^{B}\hat{\boldsymbol f}
  &=\boldsymbol a_{u,0}^{\mathrm T}
    \boldsymbol R_y(q_{e,g})\boldsymbol a_{f,0}
  \nonumber\\
  &=\sin q_{e,g}.
  \label{eq:ch4_elbow_dot_product_derivation}
\end{align}
旋转矩阵满足 $\boldsymbol R^{\mathrm T}\boldsymbol R=\boldsymbol I$，因此对任意向量
$\boldsymbol a,\boldsymbol b$ 都有
$(\boldsymbol R\boldsymbol a)^{\mathrm T}(\boldsymbol R\boldsymbol b)
=\boldsymbol a^{\mathrm T}\boldsymbol b$，即共同旋转不改变夹角或内积。在本式中，肩部旋转
${}^{B}\boldsymbol R_S$ 因而被消去；再代入
$\boldsymbol a_{u,0}=[0,0,-1]^{\mathrm T}$、
$\boldsymbol a_{f,0}=[1,0,0]^{\mathrm T}$ 和 $\boldsymbol R_y(q_{e,g})$ 的矩阵形式，内积
恰好等于 $\sin q_{e,g}$。

因此，式~\eqref{eq:ch4_geometric_elbow_angle} 可以直接由两个肢段方向求得，而不依赖肩部
三个角的具体数值。在本研究的约定下，$q_{e,g}=0$ 对应约 $90^\circ$ 屈肘，
$q_{e,g}=\pi/2$ 对应伸直。为避免与临床文献中常用的“屈曲角”混淆，本文将其称为肘部
伸展量。

这种零位来自当前 G1 左臂模型与重定向接口的关节语义：机器人模型零位附近，大臂向下而
前臂向前，二者约成直角；正向 $q_e$ 将前臂逐渐转到与大臂同向。因此本文的 $q_e$ 是
“从约 $90^\circ$ 弯肘零位出发的伸展量”，不等同于临床上通常从完全伸直定义的肘屈曲角。

几何肘角直接反映当前两条估计方向之间的关系。最终发布的肘部变量使用
式~\eqref{eq:ch4_elbow_affine_correction} 进行安装级增益和零偏修正。这里应区分
$q_{e,g}$ 与 $q_e$：前者用于保持几何分解与测量方向一致，后者是面向机器人目标的校正后
输出。
必须保留两者，是因为 wrist 残余姿态分解要求理论前臂方向与当前实测大臂、前臂方向严格
一致，故使用未修正的 $q_{e,g}$；而输出给重定向层的肘角需要补偿安装造成的尺度和零偏，故
使用 $q_e$。若提前用 $q_e$ 构造理论前臂姿态，标定补偿量会被错误地吸收到 wrist twist 中。

\subsection{肩部 yaw}
\label{subsec:ch4_shoulder_yaw}

大臂方向本身不能确定绕大臂轴的旋转。系统使用前臂在垂直于大臂方向平面内的分量估计
肩部 yaw：
\begin{equation}
  \boldsymbol f_{\perp}=
  {}^{B}\hat{\boldsymbol f}-
  \left({}^{B}\hat{\boldsymbol f}^{\mathrm T}
  {}^{B}\hat{\boldsymbol u}\right)
  {}^{B}\hat{\boldsymbol u}.
  \label{eq:ch4_forearm_perpendicular_component}
\end{equation}

若 $\lVert\boldsymbol f_{\perp}\rVert$ 足够大，将其归一化，并与在当前 pitch、roll 下的零
yaw 参考方向比较。几何上，$\boldsymbol f_{\perp}$ 是前臂方向在大臂法平面内的投影，
它确定了“肩--肘--腕”平面绕大臂纵轴的朝向。设零 yaw 参考方向投影并归一化后为
$\boldsymbol f_{0,\perp}$，则
\begin{equation}
  q_y=\operatorname{atan2}
  \left(
  -{}^{B}\hat{\boldsymbol u}^{\mathrm T}
  (\boldsymbol f_{0,\perp}\times\boldsymbol f_{\perp}),
  \boldsymbol f_{0,\perp}^{\mathrm T}\boldsymbol f_{\perp}
  \right).
  \label{eq:ch4_shoulder_yaw_inverse}
\end{equation}

投影的作用是去除前臂沿大臂轴的分量，只保留能够描述肘平面绕大臂轴朝向的分量。归一化
则去除投影长度，只比较方向；但只有投影长度足够大时，归一化才不会放大噪声。零 yaw 参考
方向 $\boldsymbol f_{0,\perp}$ 由当前 $q_p,q_r$ 和 $q_y=0$ 的运动学模型生成，因此比较二者
即可得到实际肘平面相对零 yaw 肘平面的有符号转角。式~\eqref{eq:ch4_shoulder_yaw_inverse}
中，点积给出余弦项，沿大臂轴的叉积投影给出带符号正弦项，$\operatorname{atan2}$ 将二者
合成为保留象限的有符号角。前置负号用于匹配 G1 肩 yaw 的正方向约定。

\begin{figure}[htbp]
  \centering
  \fbox{\parbox{0.92\linewidth}{\centering
    左图：在垂直于大臂方向 ${}^{B}\hat{\boldsymbol u}$ 的平面内，画出前臂投影
    $\boldsymbol f_{\perp}$、零 yaw 参考 $\boldsymbol f_{0,\perp}$ 和有符号角 $q_y$；
    右图：对比弯肘时投影较长、yaw 与 wrist 可分辨，以及直臂时投影趋近零、两种轴向旋转
    耦合的情形。}}
  \caption{肩部 yaw 的平面投影求解及直臂不可观测性（待替换为正式绘图）}
  \label{fig:ch4_yaw_projection_observability}
\end{figure}

与 pitch 相同，yaw 也展开到最接近上一帧的等价角。由于旋转 $q_y$ 与
$q_y+2k\pi$（$k\in\mathbb Z$）表示同一姿态，展开只选择最连续的数值分支，并不改变物理
旋转。

\subsection{前臂旋前/旋后}
\label{subsec:ch4_forearm_pronation_supination}

当前实现使用与两条测量方向（大臂单位方向 ${}^{B}\hat{\boldsymbol u}$ 和前臂单位方向
${}^{B}\hat{\boldsymbol f}$）直接一致的几何肘角 $q_{e,g}$ 构造不含 wrist roll 的理论
前臂旋转
\begin{equation}
  {}^{B}\boldsymbol R_{E,g}=
  \boldsymbol R_y(q_p)\boldsymbol R_x(q_r)
  \boldsymbol R_z(q_y)\boldsymbol R_y(q_{e,g}).
  \label{eq:ch4_forearm_rotation_without_wrist}
\end{equation}
当前接收器在人体标准坐标系中的旋转为
\begin{equation}
  {}^{B}\boldsymbol R_F=
  \boldsymbol A_t\,{}^{T}\boldsymbol R_F.
  \label{eq:ch4_forearm_sensor_body_rotation}
\end{equation}
去除理论前臂旋转和固定安装旋转后，残余旋转为
\begin{equation}
  \boldsymbol R_{res}=
  ({}^{B}\boldsymbol R_{E,g})^{\mathrm T}
  {}^{B}\boldsymbol R_F\boldsymbol R_m^{\mathrm T}.
  \label{eq:ch4_wrist_residual_rotation}
\end{equation}

上述三式按“预测---观测---求残差”的顺序工作。式
~\eqref{eq:ch4_forearm_rotation_without_wrist} 用已估计的肩、肘量预测前臂在未转掌时的姿态；
式~\eqref{eq:ch4_forearm_sensor_body_rotation} 将传感器观测变换到同一人体坐标系；式
~\eqref{eq:ch4_wrist_residual_rotation} 左乘理论姿态的逆、右乘固定安装旋转的逆，从观测中
去除肩部运动、肘部运动和外壳安装角。理想剩余量只包含前臂绕自身纵轴的旋转。本文沿用
机器人关节名将该量称为 wrist roll，但其人体意义是前臂旋前/旋后，并非手腕屈伸。

等价地，前臂接收器姿态模型可写为
\begin{equation}
  {}^{B}\boldsymbol R_F
  \approx{}^{B}\boldsymbol R_{E,g}
  \boldsymbol R_x(q_w)\boldsymbol R_m.
  \label{eq:ch4_forearm_sensor_attitude_model}
\end{equation}
其中，$\boldsymbol R_m$ 表示本次安装中接收器与前臂模型坐标系之间的固定旋转。

理想情况下，$\boldsymbol R_{res}$ 仅包含绕前臂局部 $+x$ 轴的旋转。为抑制小幅残余 swing
误差，系统将 $\boldsymbol R_{res}$ 转为四元数
$\boldsymbol q_{res}=[q_0,q_x,q_y,q_z]^{\mathrm T}$，只保留向量部分在 $+x$ 轴上的投影，
得到
\begin{equation}
  q_w=\operatorname{wrap}_{[-\pi,\pi)}
  \left(2\operatorname{atan2}(q_x,q_0)\right).
  \label{eq:ch4_wrist_twist_angle}
\end{equation}
然后将结果展开到上一帧附近，以维持时间连续性。这一过程等价于围绕前臂纵轴的
swing--twist 分解：swing 改变纵轴指向，twist 则绕该纵轴旋转；只保留 $q_x$ 相当于提取
局部 $+x$ 轴上的 twist 分量。此类分解常用于把三维关节旋转分离为摆动和轴向扭转
~\cite{murthy2019joint}。结果展开的原因与 yaw 相同：四元数给出的主角位于有限区间，相邻
帧若跨过 $\pm\pi$ 会在数值上跳变约 $2\pi$，选择离上一帧最近的等价角可保持连续。
wrist 提取完成后，系统才将 $q_{e,g}$ 修正为最终输出 $q_e$；因此，
肘角的经验增益和零偏不会被误解释为额外的前臂轴向旋转。

\section{不可观测区域与稳健性处理}
\label{sec:ch4_observability_robustness}

\subsection{直臂状态下 yaw 与 wrist 的耦合}
\label{subsec:ch4_yaw_wrist_unobservable}

当肘部接近完全伸直时，大臂和前臂方向近似共线，
$\lVert\boldsymbol f_{\perp}\rVert\rightarrow0$。此时，仅凭两个肢段方向无法区分“肩部
绕大臂轴旋转”和“前臂自身旋转”；同一物理姿态可以对应多组 $(q_y,q_w)$。这是测量模型
的几何不可观测性，而不是数值求解精度不足。

原因是直臂时肩 yaw 轴与前臂 twist 轴近似重合，末端姿态主要取决于二者的组合量。对任意
小角度 $\delta$，参数对
\begin{equation}
  (q_y,q_w)\sim(q_y+\delta,q_w-\delta)
  \label{eq:ch4_yaw_wrist_equivalent_pairs}
\end{equation}
可产生相同或近似相同的总轴向旋转，因而无法由当前两接收器数据唯一判断该旋转应归属于
肩部还是前臂。弯肘后两条旋转轴不再共线，图
~\ref{fig:ch4_yaw_projection_observability} 所示投影重新出现，两个变量才恢复可分辨性。

系统定义
\begin{equation}
  o_y=\lVert\boldsymbol f_{\perp}\rVert_2
  \label{eq:ch4_yaw_observability}
\end{equation}
为 yaw 可观测度。由于 $\boldsymbol u$ 和 $\boldsymbol f$ 均为单位向量，且
$\boldsymbol u^{\mathrm T}\boldsymbol f=\sin q_{e,g}$，可进一步得到
\begin{equation}
  o_y=\sqrt{1-(\boldsymbol u^{\mathrm T}\boldsymbol f)^2}
  =|\cos q_{e,g}|.
  \label{eq:ch4_yaw_observability_elbow_relation}
\end{equation}
因此，这一退化不是经验上偶然出现的数值问题，而是由五自由度模型的几何结构直接决定。
弯肘约 $90^\circ$ 时 $q_{e,g}\approx0$、$o_y\approx1$，肘平面方向清晰；手臂趋于伸直时
$q_{e,g}\rightarrow\pi/2$、$o_y\rightarrow0$，肩 yaw 轴与前臂 twist 轴趋于重合，测量
中只能保留两者的组合旋转而不能唯一分离。

当 $o_y<0.90$ 时，保持最后一次可靠的 $q_y$ 和 $q_w$，避免把微小方向
噪声放大为几十度的轴向旋转。由于本模型中 $o_y$ 随肘部伸直而下降，该阈值会在完全伸直
之前提前冻结轴向角。系统恢复到明显弯肘姿态后，再继续更新 yaw 和 wrist。

因此，直臂前举和直臂侧举可以用于评价大臂方向及肘角，但不适合
把直臂主动转掌作为 wrist 精度的主要证据。wrist 单关节实验应保持肘部弯曲约 $60^\circ$
至 $90^\circ$，使肩部轴向旋转和前臂旋转保持可分辨。

\subsection{周期角展开与时间连续性}
\label{subsec:ch4_temporal_conditioning}

肩部 yaw、wrist roll 和其他周期角首先展开到最接近上一帧的 $2\pi$ 等价分支。反三角函数
通常只返回 $[-\pi,\pi)$ 内的主值，因此真实角度从 $+\pi$ 连续增加时，下一帧数值会变成
接近 $-\pi$；两帧物理上很接近，数值却相差近 $2\pi$。角展开在
$\theta+2k\pi$ 的等价集合中选择最接近上一帧的分支，从而避免该边界产生整周跳变。侧举
奇异区和直臂不可观测区的内部状态也只在满足带滞回的
重新观测条件后更新，防止临界点附近反复切换。

本章输出的是经过可观测性处理的五自由度人体估计，不在此处施加 G1 关节限位、机器人
命令低通滤波或真机速度限制。上述机器人目标调理统一放在第5章说明，从而把“人体状态
估计”和“机器人命令生成”分开。

\subsection{无效数据与看门狗}
\label{subsec:ch4_invalid_data_watchdog}

在线估计过程中，以下情况不会产生新的关节目标：
\begin{itemize}
  \item 接收器数据陈旧或双流时间差过大；
  \item 原始四元数模长异常；
  \item 两个话题的坐标系标识不一致；
  \item 大臂仿射位置模型误差超过 $50\,\mathrm{mm}$；
  \item 标定矩阵无效或当前方向向量接近零。
\end{itemize}

“看门狗”是监测数据或程序是否在规定时间内持续更新的超时机制。若连续
$0.25\,\mathrm{s}$ 没有新的有效同步数据对，重定向节点发出看门狗警告并停止更新
目标。真机桥还具有独立的 UDP 目标超时和控制权释放机制。因此，采集层、重定向层与执行
层分别检查数据新鲜度，避免单点故障持续驱动机器人。

\section{算法流程与实现}
\label{sec:ch4_algorithm_implementation}

一次完整在线估计可归纳为算法 4-1。

\begin{quote}
\textbf{算法 4-1：双电磁接收器人体左臂姿态估计}

\textbf{输入：}同步接收器位姿
$({}^{T}\boldsymbol p_U,{}^{T}\boldsymbol R_U,
{}^{T}\boldsymbol p_F,{}^{T}\boldsymbol R_F)$，当次安装标定参数，上一帧角度。

\textbf{输出：}五自由度人体手臂估计
$\hat{\boldsymbol q}_h=[q_p,q_r,q_y,q_e,q_w]^{\mathrm T}$。

\begin{enumerate}[label=\arabic*.]
  \item 检查数据新鲜度、双流时间差、四元数模长和共同时间戳；
  \item 根据当次标定选择原始或转置旋转约定；
  \item 用大臂接收器位置和仿射椭球模型计算 ${}^{B}\hat{\boldsymbol u}$，并检查模型误差；
  \item 用前臂方向对齐矩阵和局部方向轴计算 ${}^{B}\hat{\boldsymbol f}$；
  \item 根据 ${}^{B}\hat{\boldsymbol u}$ 计算 pitch 和 roll，并处理侧举奇异区；
  \item 根据两条肢段方向的内积计算几何肘角 $q_{e,g}$；
  \item 计算 yaw 可观测度；可观测时求解 yaw，否则保持上一帧 yaw；
  \item 使用 $q_{e,g}$ 去除肩部、肘部和接收器固定安装旋转，提取 wrist twist；yaw
        不可观测时同时保持上一帧 wrist；
  \item 对 $q_{e,g}$ 应用本次安装的增益和零偏，得到最终肘部输出 $q_e$；
  \item 将周期角展开到上一帧附近，输出五自由度人体估计；
  \item 记录原始方向、可观测度、原始角、修正角和质量指标，用于离线评价。
\end{enumerate}
\end{quote}

实现中，双接收器采集、标定和在线估计分别由 \texttt{dual\_em\_tracker\_node}、
\texttt{dual\_calibration\_node} 和 \texttt{dual\_retargeting\_node} 完成。纯几何运算被
独立放置在 \texttt{retargeting\_kinematics.py} 和
\texttt{dual\_retargeting\_kinematics.py} 两个无 ROS 依赖的运动学模块中，便于通过单元
测试验证坐标变换、奇异区分支和角度连续性。每轮运行结束时，系统将全频率诊断数据保存为
NPZ（NumPy compressed archive，NumPy 压缩数组归档）文件。该格式可在一个压缩文件中按
名称保存多个同长度数组，便于保持时间序列字段之间的一一对应。记录内容包括接收器原始位姿、人体
肢段方向、五个原始关节角、肘角修正前后数值、pitch/yaw 可观测度以及最终命令。这些数据
为第6章实验设计和第7章结果分析提供统一的数据来源。

\section{本章小结}
\label{sec:ch4_summary}

本章提出并实现了面向重复佩戴的双电磁接收器人体左臂估计方法。该方法没有假设传感器在
每次佩戴后保持相同朝向，而是通过静态质量检查、接收器身份检查、独立前臂旋转轴动作和
七个参考姿势，重新辨识当前安装的坐标对齐、肢段方向轴、大臂位置模型以及肘角修正参数。

本章进一步说明了测量信息与模型自由度之间的关系：大臂单位方向提供两个独立自由度，前臂
完整姿态提供三个独立自由度，从而在一般姿态下支持五自由度状态的逐级几何求解。这五个量
是与 G1 关节语义一致的模型广义坐标，而非完整人体上肢的解剖学真值。

为降低软组织和绑带运动对大臂姿态测量的影响，系统使用仿射位置--方向模型恢复大臂方向；
该模型在几何上把单位方向球面映射为安装相关的椭球轨迹，但并不假设人体肩关节本身为椭球。
前臂方向与旋前/旋后使用相互独立的轴和对齐参数。在线阶段，系统从两个肢段方向计算肩部
pitch、roll、yaw 和几何肘角，通过残余旋转的 swing--twist 分解得到 wrist roll，再将
安装级增益和零偏作用于最终肘部输出。针对侧举
时 pitch 退化和直臂时 yaw/wrist 耦合，本章分别设计了带滞回的标准分支和最后可靠值保持
策略。最后，标定质量门槛、同步检查、位置模型误差和看门狗共同阻止低质量数据进入后续
控制链。

下一章将在本章输出的五自由度人体动作基础上，介绍 G1 关节目标生成、关节和速度限制、
躯干碰撞保护以及真机 Arm SDK 安全桥。

。

\chapter{自适应标定与人体手臂姿态估计}
\label{chap:adaptive_calibration_arm_estimation}

\section{安装标定与姿态估计相关研究}
\label{sec:ch4_related_calibration}

可穿戴传感器报告的是传感器外壳相对参考坐标系的位姿，而人体运动分析需要的是肢段坐标系或
关节变量。两者之间的传感器--肢段关系会随佩戴位置和朝向改变，是可穿戴上肢测量中的核心
误差来源。国际生物力学学会（ISB）给出了肩、肘、腕和手部坐标系及关节运动报告的推荐
定义，为不同测量系统之间统一轴向和角度语义提供了基础~\cite{wu2005isb}。Picerno 等通过
解剖标志点建立上肢惯性传感器与肢段坐标系之间的关系，并使用光学
系统验证肩、肘运动学~\cite{picerno2019upper}。Seel 等利用运动学约束从功能动作中识别关节轴，
说明不必要求传感器外壳与解剖轴精确平行~\cite{seel2014imu}。Zabat 等进一步表明，相比单一
静态姿势，多姿态信息可以减小安装偏差和软组织运动引起的角度误差~\cite{zabat2019multiple}。
近期研究还指出，待测系统与参考系统若采用不同的传感器--肢段定义，会形成明显的系统性关节角
偏差~\cite{adjel2025influence}。

现有标定方法大体可分为静态、功能性和多姿态方法。静态方法操作简单，但对参考姿势执行误差
敏感；功能性方法通过重复运动拟合关节轴，降低了对安装朝向的要求，但依赖充分且近似单轴的
动作激励；多姿态方法利用若干互补方向联合估计安装关系，并可通过各姿态残差发现错误动作或
绑带移动。对肩部三自由度、肘部屈伸和前臂轴向旋转同时估计时，单一姿态不能提供足够约束，
尤其在大臂与前臂接近共线时，肘角和轴向旋转的可观测性会降低。

Cutti 等使用惯性与磁传感器进行肩、肘动态测量的研究也表明，统一的肢段定义、安装标定和
磁环境检查是获得可解释上肢角度的共同前提~\cite{cutti2008ambulatory}。

电磁设备还存在“设备场标定”与“人体安装标定”两个不同层次。前者从已知空间位姿建立电磁
观测与接收器位姿之间的映射，并处理工作空间内的位置相关误差；相关综述与标准化评估研究
均指出，金属畸变、工作距离和空间位置会共同影响电磁跟踪误差
~\cite{kindratenko2000survey,frantz2003accuracy,hummel2006evaluation,franz2014electromagnetic,
maclachlan2025calibration}。后者在设备位姿已经
可用的前提下，估计接收器外壳与人体肢段之间的关系。本文不修改设备固件或制造商场模型，而是
针对每次佩戴执行多姿态安装级标定，并把姿态残差、位置模型残差和重复性作为进入在线估计前的
质量门控。

\section{问题描述与总体思路}
\label{sec:ch4_problem_overview}

本研究使用一套电磁位姿跟踪系统获取操作者左臂运动。电磁发射器固定于操作者躯干，两个
接收器分别固定于左大臂和左前臂。系统以发射器坐标系为共同参考，以约 $50\,\mathrm{Hz}$
的频率输出两个接收器的三维位置和姿态。与基于人体外观的视觉方法相比，电磁测量不依赖
纹理和光照，也不存在人体自身造成的直接视觉遮挡。然而，接收器在每次佩戴后的安装朝向
并不相同，绑带和软组织还会使接收器运动偏离理想的刚体关节模型。因此，原始接收器位姿
不能直接作为机器人关节角使用，必须先完成针对本次安装的标定。这里的“原始位姿”专指
设备直接报告的接收器外壳相对发射器的位置和旋转；它尚未去除外壳安装角，也尚未转换为
人体大臂、前臂方向或关节角。

本章所研究的问题可表述为：给定大臂接收器和前臂接收器在电磁发射器坐标系中的同步位姿，
估计人体左臂的五个运动量
\begin{equation}
  \boldsymbol q_h =
  \begin{bmatrix}
    q_p & q_r & q_y & q_e & q_w
  \end{bmatrix}^{\mathrm T},
  \label{eq:ch4_human_arm_state}
\end{equation}
其中，$q_p$、$q_r$ 和 $q_y$ 分别表示肩部前举/后摆、外展/内收和轴向旋转，$q_e$
表示肘部由约 $90^\circ$ 屈曲姿态向伸直方向的变化量，$q_w$ 表示前臂旋前/旋后。
上述变量沿用 Unitree G1 左臂的关节语义，但其数值首先描述人体动作，随后才在重定向和
安全层中被限制为机器人可执行目标。

需要强调的是，式~\eqref{eq:ch4_human_arm_state} 中的五个量是本研究所选运动学模型的
广义坐标，而不是由传感器直接测得的五个解剖学关节角。人体肩胛骨、锁骨、腕部屈伸以及
关节中心的平移均未被单独建模。本文的目标是构造一组与 G1 左臂关节语义一致、能够由当前
两接收器配置稳定估计的低维运动描述，而不是恢复完整的人体上肢骨骼模型。

系统采用“安装级自适应标定”思路。这里的“自适应”并非在线学习或自适应控制，而是指每次
重新佩戴设备后，系统根据本次安装采集的参考动作重新辨识：
\begin{enumerate}[label=\arabic*)]
  \item 电磁发射器坐标系到人体标准坐标系的映射；
  \item 接收器局部坐标系中的肢段纵轴；
  \item 大臂接收器位置与人体大臂方向之间的映射；
  \item 前臂旋前/旋后的局部旋转轴和零位；
  \item 肘角几何估计的增益和偏置。
\end{enumerate}

在线阶段只使用本次安装生成的参数。任一接收器或发射器被拆卸、重新绑扎或在支架中移动
后，原标定不再有效，必须重新执行标定流程。
换言之，“安装级”表示参数在一次佩戴完成后被辨识，并在该次运行期间保持固定；它不表示
参数随每一帧在线更新。这样既允许系统适应重复佩戴造成的安装差异，又避免在线更新参数时
把操作者的真实动作错误地吸收到标定模型中。

\subsection{坐标系与符号}
\label{subsec:ch4_frames_notation}

本章使用表~\ref{tab:ch4_coordinate_frames} 所列坐标系。
\begin{table}[htbp]
  \centering
  \caption{人体手臂估计使用的坐标系}
  \label{tab:ch4_coordinate_frames}
  \begin{tabularx}{\linewidth}{@{}llX@{}}
    \toprule
    符号 & 英文含义 & 物理含义 \\
    \midrule
    $\{T\}$ & Transmitter & 电磁发射器坐标系 \\
    $\{B\}$ & Body & 人体标准坐标系，与机器人约定一致：$+x$ 向前、$+y$ 向左、$+z$ 向上 \\
    $\{U\}$ & Upper arm & 大臂接收器局部坐标系 \\
    $\{F\}$ & Forearm & 前臂接收器局部坐标系 \\
    \bottomrule
  \end{tabularx}
\end{table}

记 ${}^{T}\boldsymbol p_U$ 和 ${}^{T}\boldsymbol p_F$ 为两个接收器在 $\{T\}$ 中的位置，
${}^{T}\boldsymbol R_U$ 和 ${}^{T}\boldsymbol R_F$ 为将接收器局部向量旋转到 $\{T\}$ 的
旋转矩阵。上标表示量所在的表达坐标系，下标表示被描述的对象；例如
${}^{T}\boldsymbol p_U$ 表示“大臂接收器的位置，用发射器坐标系表达”。
${}^{B}\boldsymbol u$ 和 ${}^{B}\boldsymbol f$ 分别表示人体大臂和前臂在
$\{B\}$ 中的单位方向。本文所称“肢段”是由相邻关节界定、在当前模型中近似视为刚体的
身体部分：大臂肢段连接肩与肘，前臂肢段连接肘与腕。所谓单位方向，是指只保留肢段的朝向
而忽略其长度，即
\begin{equation}
  \lVert{}^{B}\boldsymbol u\rVert_2=1,
  \qquad
  \lVert{}^{B}\boldsymbol f\rVert_2=1.
  \label{eq:ch4_unit_segment_directions}
\end{equation}
将单位向量平移到 $\{B\}$ 的数学原点后，其终点位于单位球面上。向量是自由向量，平移
起点不会改变其方向；因此这里的单位球面只是统一表示所有三维方向的数学空间，并不要求
人体肩关节位于 $\{B\}$ 原点，也不表示接收器在物理上沿半径为 1 的球面运动。

为与 G1 关节模型保持一致，本研究将零位大臂方向和零位前臂方向定义为
\begin{equation}
  \boldsymbol a_{u,0}=\begin{bmatrix}0&0&-1\end{bmatrix}^{\mathrm T},
  \qquad
  \boldsymbol a_{f,0}=\begin{bmatrix}1&0&0\end{bmatrix}^{\mathrm T}.
  \label{eq:ch4_zero_segment_axes}
\end{equation}
因此，零位对应“大臂自然下垂、肘部弯曲约 $90^\circ$、前臂水平向前”的姿态。
这里的“零位方向”是运动学模型在所有广义坐标均为零时规定的参考方向，并非接收器局部
坐标轴，也不等同于医学上的
标准解剖姿势。当所有模型关节变量为零时，大臂和前臂方向分别等于
$\boldsymbol a_{u,0}$ 和 $\boldsymbol a_{f,0}$；其他姿态均由关节旋转作用于这两个参考
方向得到。图~\ref{fig:ch4_frames_zero_pose} 汇总了坐标轴、接收器和零位方向之间的关系。

\begin{figure}[htbp]
  \centering
  \fbox{\parbox{0.90\linewidth}{\centering
    示意图内容：人体左臂处于 DOWN 零位；标出 $\{B\}$ 的 $+x$ 前、$+y$ 左、$+z$ 上，
    发射器坐标系 $\{T\}$，大臂与前臂接收器坐标系 $\{U\}$、$\{F\}$，以及
    $\boldsymbol a_{u,0}=[0,0,-1]^{\mathrm T}$ 和
    $\boldsymbol a_{f,0}=[1,0,0]^{\mathrm T}$。}}
  \caption{双接收器安装、坐标系与运动学零位示意（待替换为正式绘图）}
  \label{fig:ch4_frames_zero_pose}
\end{figure}

\subsection{五自由度运动学模型}
\label{subsec:ch4_five_dof_model}

本研究用肩部三个旋转、肘部一个旋转以及前臂一个轴向旋转描述左臂运动，即
\begin{equation}
  3\ \text{(shoulder)}+1\ \text{(elbow)}+1\ \text{(forearm twist)}
  =5\ \text{DOF}.
  \label{eq:ch4_five_dof_count}
\end{equation}
肩部三个变量 $q_p$（pitch，前举/后摆）、$q_r$（roll，外展/内收）和 $q_y$（yaw，绕
大臂纵轴旋转）共同确定大臂的空间方向及肘平面朝向；$q_e$ 描述肘部由弯曲向伸直的变化，
$q_w$ 描述前臂旋前/旋后。本模型没有包含 G1 腕部的 pitch 和 yaw，也没有单独描述肩胛带
运动，因此应理解为面向动作重定向的五自由度表示，而非完整的人体上肢解剖模型。

肩部和肘部的正向运动学采用
\begin{align}
  {}^{B}\boldsymbol R_S
  &=\boldsymbol R_y(q_p)\boldsymbol R_x(q_r)\boldsymbol R_z(q_y),
  \label{eq:ch4_shoulder_rotation}\\
  {}^{B}\boldsymbol R_E
  &={}^B\boldsymbol R_S\boldsymbol R_y(q_e),
  \label{eq:ch4_forearm_rotation}
\end{align}
其中，${}^{B}\boldsymbol R_S$ 表示肩部旋转，${}^{B}\boldsymbol R_E$ 表示不包含腕部
旋转的前臂坐标系。相应的肢段方向为
\begin{equation}
  {}^{B}\boldsymbol u={}^B\boldsymbol R_S\boldsymbol a_{u,0},
  \qquad
  {}^{B}\boldsymbol f={}^B\boldsymbol R_E\boldsymbol a_{f,0}.
  \label{eq:ch4_segment_directions}
\end{equation}

式~\eqref{eq:ch4_shoulder_rotation} 的乘法顺序是模型定义的一部分，不能任意交换。本文采用
列向量和主动旋转约定，因此右侧旋转先作用：按矩阵作用顺序，零位肢段依次接受
$\boldsymbol R_z(q_y)$、$\boldsymbol R_x(q_r)$ 和 $\boldsymbol R_y(q_p)$。各轴均以本章统一的
矩阵约定解释，而不是在推导中临时更换为另一套欧拉角约定。矩阵乘法不满足交换律，改变顺序会改变同一组
角度对应的空间姿态。本顺序与工程中 G1 左肩三个关节的语义保持一致，而不是声称它是唯一的
人体解剖角定义；解剖学结果的轴序应另按 ISB 推荐报告~\cite{wu2005isb}。

该模型的优点是人体估计量与 G1 五个左臂关节具有一致的语义，避免了先恢复完整人体骨架
再进行第二次逆运动学求解。模型的限制是：仅由两个肢段方向不能在所有姿态下唯一分离肩部
轴向旋转与前臂旋转。第~\ref{sec:ch4_observability_robustness} 节将对这一可观测性问题进行
专门处理。

\subsection{测量信息与五自由度状态的对应关系}
\label{subsec:ch4_measurement_dof_correspondence}

为了说明两个接收器为何能够支持五个运动量的估计，需要区分“数据分量个数”和“独立
自由度个数”。一个三维单位方向属于单位球面
\begin{equation}
  S^2=\left\{\boldsymbol v\in\mathbb R^3\mid
  \boldsymbol v^{\mathrm T}\boldsymbol v=1\right\}.
  \label{eq:ch4_unit_sphere_definition}
\end{equation}
虽然 $\boldsymbol v$ 具有三个坐标分量，但单位长度约束消除了一个独立变量，因此一个方向
只有两个自由度。直观上，这两个自由度类似球面上的经度和纬度；它们能够确定肢段指向，
却不能确定肢段绕自身纵轴旋转了多少。

旋转矩阵属于特殊正交群
\begin{equation}
  SO(3)=\left\{\boldsymbol R\in\mathbb R^{3\times3}\mid
  \boldsymbol R^{\mathrm T}\boldsymbol R=\boldsymbol I,
  \det(\boldsymbol R)=1\right\}.
  \label{eq:ch4_so3_definition}
\end{equation}
旋转矩阵虽然包含九个元素，但正交和单位长度约束使其只有三个独立自由度。也可以将其理解
为“一个轴的空间方向”所需的两个自由度，加上“绕该轴旋转”所需的一个自由度。

“安装标定”与“在线估计”是前后相接的两个阶段。安装标定只在重新佩戴后执行，利用已知
参考姿势求出坐标对齐、局部肢段轴、位置模型和角度修正参数；在线估计则在每一个同步采样
时刻使用这些固定参数，把未知动作转换为五自由度状态，不再重新拟合参数。完成安装标定后，
在线估计实际使用两类几何信息：大臂单位方向
${}^{B}\hat{\boldsymbol u}\in S^2$，以及经过安装对齐的前臂完整姿态
${}^{B}\boldsymbol R_F\in SO(3)$。因此可用独立信息量为
\begin{equation}
  \underbrace{2}_{\text{upper-arm direction}}+
  \underbrace{3}_{\text{forearm attitude}}=5.
  \label{eq:ch4_measurement_dof_count}
\end{equation}
前臂姿态的三个自由度还可以分解为前臂方向 ${}^{B}\hat{\boldsymbol f}$ 的两个自由度和
绕前臂纵轴旋转的一个自由度。由此形成表~\ref{tab:ch4_measurement_to_joint} 所示的逐级
求解关系。

\begin{table}[htbp]
  \centering
  \caption{测量几何量与五自由度状态的对应关系}
  \label{tab:ch4_measurement_to_joint}
  \begin{tabularx}{\linewidth}{@{}lXX@{}}
    \toprule
    待估计量 & 主要几何依据 & 物理含义 \\
    \midrule
    $q_p,q_r$ & 大臂单位方向 ${}^{B}\hat{\boldsymbol u}$ & 大臂指向的两个自由度 \\
    $q_e$ & ${}^{B}\hat{\boldsymbol u}$ 与 ${}^{B}\hat{\boldsymbol f}$ 的夹角
      & 肘部由弯曲向伸直的变化量 \\
    $q_y$ & 前臂方向在大臂法平面内的投影 & 肘平面绕大臂纵轴的朝向 \\
    $q_w$ & 去除肩、肘和固定安装旋转后的前臂残余姿态
      & 绕前臂纵轴的旋前/旋后 \\
    \bottomrule
  \end{tabularx}
\end{table}

上述自由度计数只说明一般姿态下信息量足够，并不保证所有姿态都存在唯一解。当大臂正侧举
时，pitch 方向退化；当肘部完全伸直时，肩部轴向旋转与前臂旋转发生耦合。换言之，该测量
映射在一般姿态下局部可辨识，但在特定几何构型下秩会下降。系统对这些区域的处理将在
第~\ref{sec:ch4_observability_robustness} 节给出。

\section{双接收器同步采集与数据预处理}
\label{sec:ch4_synchronized_acquisition}

\subsection{同步位姿对}
\label{subsec:ch4_synchronized_pose_pairs}

两个 WEMT 接收器分别通过网络返回独立的数据流。如果直接使用两条消息各自最近的样本，
快速运动时可能把不同时刻的大臂和前臂姿态组合在一起，从而产生虚假的肘角和肩部旋转。
因此，采集节点只在以下条件同时满足时发布一对位姿：
\begin{enumerate}[label=\arabic*)]
  \item 两个接收器均处于连接状态且没有设备错误；
  \item 两个接收器都产生了新的硬件帧；
  \item 任一位姿的接收时间不超过 $100\,\mathrm{ms}$；
  \item 两个接收器的到达时间差不超过 $30\,\mathrm{ms}$；
  \item 两个原始四元数均通过有效性检查。
\end{enumerate}

通过检查后，系统给两个 \texttt{geometry\_msgs/PoseStamped} 消息赋予同一个 ROS 时间戳。
该消息由三维位置、四元数姿态、参考坐标系标识和时间戳组成。两条消息分别发布至 ROS 话题
\texttt{/em/upper\_arm\_pose} 和 \texttt{/em/forearm\_pose}。这样，后续标定和在线估计
可按相同时间戳进行严格配对。WEMT 接口内部的平滑窗口被设为 1，即发布未经时间平均的
单帧样本；这里的 1 是窗口中包含的样本数，而不是 1 秒。该设置避免位置和姿态使用不同
滤波造成额外的相位差；统一滤波在完成双流配对后进行。

\subsection{四元数有效性与旋转矩阵}
\label{subsec:ch4_quaternion_validation}

设设备返回四元数为
\begin{equation}
  \boldsymbol q=\begin{bmatrix}q_w&q_x&q_y&q_z\end{bmatrix}^{\mathrm T}.
  \label{eq:ch4_quaternion}
\end{equation}
系统首先检查其模长是否满足
\begin{equation}
  0.95\leq \lVert\boldsymbol q\rVert_2\leq1.05.
  \label{eq:ch4_quaternion_norm_gate}
\end{equation}

未通过检查的数据对被直接丢弃，而不是先强制归一化后继续使用。这是因为持续出现约
0.707 等异常模长时，问题通常来自设备原始数据失效；单纯归一化会掩盖错误并把无效姿态
传入控制链。有效四元数归一化后转换为旋转矩阵：
\begin{equation}
\boldsymbol R(\boldsymbol q)=
\begin{bmatrix}
1-2(q_y^2+q_z^2) & 2(q_xq_y-q_zq_w) & 2(q_xq_z+q_yq_w)\\
2(q_xq_y+q_zq_w) & 1-2(q_x^2+q_z^2) & 2(q_yq_z-q_xq_w)\\
2(q_xq_z-q_yq_w) & 2(q_yq_z+q_xq_w) & 1-2(q_x^2+q_y^2)
\end{bmatrix}.
\label{eq:ch4_quaternion_to_rotation}
\end{equation}

由于当前设备接口只明确四元数分量顺序，而接收器到发射器的主动/被动旋转方向仍需由系统
侧验证，系统在标定时
分别评价原旋转矩阵 $\boldsymbol R$ 和转置矩阵 $\boldsymbol R^{\mathrm T}$。选择标准以前臂
旋转轴在多参考姿势下的方向残差为主，以相对旋转残差作为分数接近时的辅助判据。最终选定
的约定被写入当次标定文件，并在在线估计中保持一致。这一步不是任意改变四元数分量，而是
在互为逆旋转的两种物理解释之间进行一次安装级一致性检验；若后续由设备协议明确方向，仍可
保留该检验作为接线或固件变化的诊断。

\section{安装级自适应标定}
\label{sec:ch4_installation_calibration}

\subsection{标定流程概览}
\label{subsec:ch4_calibration_overview}

完整标定由三个前置质量/动作检查和七个静态参考姿势组成，如图~\ref{fig:ch4_calibration_flow}
所示。
\begin{figure}[htbp]
  \centering
  \fbox{\parbox{0.88\linewidth}{\centering
    图示占位：静态质量检查 $\rightarrow$ 接收器身份检查 $\rightarrow$
    前臂旋转轴标定 $\rightarrow$ 七姿势采集 $\rightarrow$
    参数拟合 $\rightarrow$ 质量门槛 $\rightarrow$ 保存或拒绝。}}
  \caption{安装级自适应标定流程（待替换为正式绘图）}
  \label{fig:ch4_calibration_flow}
\end{figure}

标定开始前，操作者保持自然站立，系统采集 $4\,\mathrm{s}$ 静态数据，用于计算位置和姿态
噪声。随后，操作者在保持大臂基本不动的情况下反复屈伸肘部并转动前臂。若前臂接收器的
姿态变化明显大于大臂接收器，才认为两个接收器的身份与安装位置正确。最后，操作者保持大臂
竖直、肘部弯曲约 $90^\circ$，仅使前臂绕“肘到腕”的纵轴做旋前/旋后，用于辨识前臂接收器
局部坐标系中的旋转轴。

七个静态参考姿势及其理论角度见表~\ref{tab:ch4_reference_poses}。每个姿势离开后重新构建
3 次，每次保持 $2\,\mathrm{s}$。这种“离开---重建---保持”的方式能够发现单次摆姿偶然
正确但重复性不足的问题。

\begin{table}[htbp]
  \centering
  \caption{安装级标定使用的静态参考姿势}
  \label{tab:ch4_reference_poses}
  \begin{tabularx}{\linewidth}{@{}XrrrrX@{}}
    \toprule
    参考姿势 & $q_p$ & $q_r$ & $q_y$ & $q_e$ & 主要作用 \\
    \midrule
    DOWN             & $0^\circ$   & $0^\circ$  & $0^\circ$ & $0^\circ$  & 定义系统零位 \\
    FORWARD\_30      & $-30^\circ$ & $0^\circ$  & $0^\circ$ & $0^\circ$  & 约束大臂位置模型 \\
    FORWARD\_60      & $-60^\circ$ & $0^\circ$  & $0^\circ$ & $0^\circ$  & 约束大臂位置模型 \\
    FORWARD          & $-90^\circ$ & $0^\circ$  & $0^\circ$ & $0^\circ$  & 约束前举方向 \\
    SIDEWAYS         & $0^\circ$   & $90^\circ$ & $0^\circ$ & $90^\circ$ & 约束侧举和伸肘方向 \\
    STRAIGHT\_DOWN   & $0^\circ$   & $0^\circ$  & $0^\circ$ & $90^\circ$ & 约束伸肘姿态 \\
    STRAIGHT\_FORWARD& $-90^\circ$ & $0^\circ$  & $0^\circ$ & $90^\circ$ & 约束主要工作空间 \\
    \bottomrule
  \end{tabularx}
\end{table}

图~\ref{fig:ch4_reference_pose_set} 给出七个参考姿势应呈现的肢段关系。FORWARD\_30 和
FORWARD\_60 的角度不要求操作者用量角器达到绝对精确值；其作用是为前举路径增加两个
中间方向。操作者应优先保证躯干不动、动作平面一致、肘部状态正确以及三次重复可重现，实际
偏差将反映在参考残差和位置模型残差中。

\begin{figure}[htbp]
  \centering
  \fbox{\parbox{0.92\linewidth}{\centering
    示意图内容：依次绘制 DOWN、FORWARD\_30、FORWARD\_60、FORWARD、SIDEWAYS、
    STRAIGHT\_DOWN 和 STRAIGHT\_FORWARD；用实线表示大臂、虚线表示前臂，并标注
    “弯肘约 $90^\circ$”或“伸肘”以及前举/侧举方向。}}
  \caption{七个安装标定参考姿势及其肢段约束（待替换为实拍图或正式绘图）}
  \label{fig:ch4_reference_pose_set}
\end{figure}

FORWARD\_30 和 FORWARD\_60 只用于提高大臂位置模型在前举路径上的可辨识性；前臂方向
拟合使用 DOWN、FORWARD、SIDEWAYS、STRAIGHT\_DOWN 和 STRAIGHT\_FORWARD 五类姿势。

\subsection{静态质量和重复性评价}
\label{subsec:ch4_static_quality_repeatability}

对于一组静态位置样本 $\boldsymbol p_i$，位置均方根波动定义为
\begin{equation}
  e_{p,\mathrm{RMS}}=
  \sqrt{\frac{1}{N}\sum_{i=1}^{N}
  \lVert\boldsymbol p_i-\bar{\boldsymbol p}\rVert_2^2}.
  \label{eq:ch4_static_position_rms}
\end{equation}

姿态中心采用旋转矩阵均值在 $SO(3)$ 上的最近旋转，姿态波动则用每个样本到中心旋转的
测地角计算 RMS。由于旋转不能像普通欧氏向量那样逐元素平均后直接使用，最近旋转投影和
四元数/旋转平均是姿态处理中常用的几何做法~\cite{hartley2013rotation,markley2007averaging}。
当前工程实现中，当任一接收器的位置 RMS 超过 $3\,\mathrm{mm}$，或姿态
RMS 超过 $3^\circ$ 时，系统提示检查绑带、金属物体和跟踪距离。

每个参考姿势的三次重复使用旋转 medoid（样本中心）作为代表，即只在三个实际观测中选择
一个，使它到其余观测的旋转测地距离总和最小。与直接平均相比，medoid 必然是一个真实采集
样本，对单次离群姿态更稳健。若任一重复与代表样本之间的最大角距离超过 $15^\circ$，则只重新采集当前
姿势，已经通过的姿势继续保留。该策略避免一次摆姿错误使整轮标定全部作废，同时保证每个
参考姿势具有最低重复性。

\subsection{前臂局部旋转轴辨识}
\label{subsec:ch4_forearm_twist_axis}

在理想旋前/旋后动作中，前臂在发射器坐标系中的纵轴方向保持近似不变，而接收器绕该轴
转动。实际只有一个前臂接收器，下标 $i$ 表示该接收器在第 $i$ 个采样时刻的观测。设其旋转
矩阵为 ${}^{T}\boldsymbol R_{F,i}$，希望找到接收器外壳坐标系中的局部单位向量
$\boldsymbol a_t$，使
\begin{equation}
  {}^{T}\boldsymbol R_{F,i}\boldsymbol a_t\approx\boldsymbol d_t,
  \qquad i=1,\ldots,N,
  \label{eq:ch4_twist_axis_invariance}
\end{equation}
其中，$\boldsymbol d_t$ 为发射器坐标系中的稳定旋转轴。寻找 $\boldsymbol a_t$ 的目的，是把
任意安装朝向下的“接收器外壳某一轴”辨识为人体前臂纵轴；若轴选择正确，接收器即使绕它
转动，该轴变换到发射器坐标系后的方向仍应近似不变。令
\begin{equation}
  \bar{\boldsymbol R}_F=\frac{1}{N}
  \sum_{i=1}^{N}{}^{T}\boldsymbol R_{F,i},
  \label{eq:ch4_mean_forearm_rotation}
\end{equation}
则 $\boldsymbol a_t$ 可写为
\begin{equation}
  \boldsymbol a_t=
  \arg\max_{\lVert\boldsymbol a\rVert=1}
  \lVert\bar{\boldsymbol R}_F\boldsymbol a\rVert_2,
  \label{eq:ch4_twist_axis_svd}
\end{equation}
即 $\bar{\boldsymbol R}_F$ 最大奇异值对应的右奇异向量。式
~\eqref{eq:ch4_twist_axis_invariance} 给出物理不变性，式
~\eqref{eq:ch4_mean_forearm_rotation} 将全部姿态观测聚合，式
~\eqref{eq:ch4_twist_axis_svd} 则寻找经这些旋转后最不发散的局部方向；最大奇异值对应方向
正是保持合向量长度最大的方向。系统对残差进行基于中位数绝对
偏差的迭代剔除，降低偶发姿态跳变的影响；同时通过样本对之间的相对旋转检查动作范围。
只有当前臂旋转范围至少达到 $55^\circ$、大臂姿态变化不超过 $15^\circ$，且旋转轴拟合
RMS 不超过 $15^\circ$ 时，才继续执行后续标定。

该旋转轴专门服务于腕部旋前/旋后估计。它不与前臂“肘到腕”的方向轴强制共用，因为皮肤
滑动、绑带安装和参考动作误差会使两者的最优拟合结果存在差异。实际实现分别辨识前臂方向
轴与旋转轴，以避免一个轴同时承担两个不同物理任务造成系统性偏差。

\subsection{大臂方向的仿射位置--方向模型}
\label{subsec:ch4_upper_arm_affine_model}

大臂接收器固定在软组织表面，其姿态容易受到皮肤相对肱骨转动的影响。实验开发过程中发现，
直接使用大臂接收器姿态估计肱骨轴向运动会引入较大误差。因此，当前混合模型主要利用大臂
接收器的位置恢复大臂方向。这里的目标不是根据位置直接回归关节角，而是先恢复具有明确
几何意义的单位方向 ${}^{B}\hat{\boldsymbol u}$，再由第~\ref{sec:ch4_online_arm_estimation}
节的逆运动学关系求解关节变量。

若接收器与肩关节构成理想刚体，大臂运动时接收器应在以肩关节为中心的球面上运动。但实际
轨迹受绑带、软组织和发射器安装影响，前举路径往往呈现非球形畸变。为此，本研究采用仿射
位置--方向模型
\begin{equation}
  {}^{T}\boldsymbol p_{U,i}
  =\boldsymbol c+\boldsymbol M\,{}^{B}\boldsymbol u_i^{*}
  +\boldsymbol\varepsilon_i,
  \label{eq:ch4_upper_affine_model}
\end{equation}
其中，$\boldsymbol c\in\mathbb R^3$ 为本次安装下的轨迹中心，
$\boldsymbol M\in\mathbb R^{3\times3}$ 为从人体标准方向到发射器坐标位移的可逆仿射
矩阵，${}^{B}\boldsymbol u_i^{*}$ 是第 $i$ 个参考姿势的理论大臂方向，
$\boldsymbol\varepsilon_i$ 汇集摆姿误差、软组织运动和测量噪声。$\boldsymbol c$ 是用于
位置--方向映射的虚拟中心，不应直接解释为解剖学肩关节中心。

这里称为“混合模型”，是因为大臂方向主要由接收器\emph{位置}恢复，而前臂方向和轴向
旋转由接收器\emph{姿态}恢复。所谓“恢复大臂方向”，是指由测得的位置
${}^{T}\boldsymbol p_U$ 反求单位向量 ${}^{B}\hat{\boldsymbol u}$，而不是估计大臂长度或
肩关节中心。图~\ref{fig:ch4_affine_geometry} 给出式~\eqref{eq:ch4_upper_affine_model} 的
几何含义。

\begin{figure}[htbp]
  \centering
  \fbox{\parbox{0.92\linewidth}{\centering
    示意图内容：以 $\boldsymbol c$ 为椭球中心，单位方向
    ${}^{B}\boldsymbol u_i^{*}$ 经过 $\boldsymbol M$ 映射为从 $\boldsymbol c$ 指向
    ${}^{T}\boldsymbol p_{U,i}$ 的位移；标出实测点、重投影点 $\hat{\boldsymbol p}_U$、
    残差 $\boldsymbol\varepsilon_i$ 与模型误差 $e_M$。左侧可并列绘制单位球，右侧绘制其
    经过 $\boldsymbol M$ 后的椭球像。}}
  \caption{大臂仿射位置--方向模型的几何解释（待替换为正式绘图）}
  \label{fig:ch4_affine_geometry}
\end{figure}

该模型被称为“仿射椭球”，是因为单位方向满足
$({}^{B}\boldsymbol u)^{\mathrm T}{}^{B}\boldsymbol u=1$。忽略噪声并令
$\boldsymbol x={}^{T}\boldsymbol p_U-\boldsymbol c$，由
$\boldsymbol x=\boldsymbol M{}^{B}\boldsymbol u$ 可得
\begin{equation}
  \boldsymbol x^{\mathrm T}\boldsymbol M^{-\mathrm T}
  \boldsymbol M^{-1}\boldsymbol x=1,
  \label{eq:ch4_affine_ellipsoid_implicit}
\end{equation}
这里引入 $\boldsymbol x$ 只是为了简化表达：它是从虚拟中心 $\boldsymbol c$ 指向当前接收器
位置的位移，不是新的测量量。将 ${}^{B}\boldsymbol u=\boldsymbol M^{-1}\boldsymbol x$
代回单位长度约束即可得到式~\eqref{eq:ch4_affine_ellipsoid_implicit}，这正是以
$\boldsymbol c$ 为中心的椭球面二次型。若对矩阵作奇异值分解
\begin{equation}
  \boldsymbol M=\boldsymbol U\boldsymbol\Sigma\boldsymbol V^{\mathrm T},
  \qquad
  \boldsymbol\Sigma=\operatorname{diag}(\sigma_1,\sigma_2,\sigma_3),
  \label{eq:ch4_affine_matrix_svd}
\end{equation}
则 $\boldsymbol U$ 给出椭球主轴方向，$\sigma_1,\sigma_2,\sigma_3$ 给出三个半轴长度。
当三个奇异值相等时，说明三个主方向的尺度相同；若
$\boldsymbol M=r\boldsymbol R$，线性变换只包含统一尺度 $r$ 和旋转 $\boldsymbol R$，模型便
退化为半径为 $r$ 的球面。若奇异值不同，单位球在三个主方向上被拉伸为不同半轴长度，形成
椭球。因此，椭球是单位方向集合经过本次安装相关线性变换后的位置轨迹，而不是对人体肩关节形状的
解剖学假设。实际七个参考姿势只覆盖该椭球面的有限区域，本模型应理解为工作空间内的经验
标定模型，而非完整肩部轨迹的全局物理模型。

所有参考样本通过最小二乘求解
\begin{equation}
  (\hat{\boldsymbol c},\hat{\boldsymbol M})=
  \arg\min_{\boldsymbol c,\boldsymbol M}
  \sum_i\left\|
  {}^{T}\boldsymbol p_{U,i}-\boldsymbol c-
  \boldsymbol M{}^{B}\boldsymbol u_i^{*}
  \right\|_2^2.
  \label{eq:ch4_upper_affine_fit}
\end{equation}
式~\eqref{eq:ch4_upper_affine_fit} 对 $\boldsymbol c$ 和 $\boldsymbol M$ 是线性最小二乘
问题。它求出的 $\hat{\boldsymbol c}$ 和 $\hat{\boldsymbol M}$ 用于在线阶段把新的接收器
位置反算为大臂方向。为了把常量中心项和三维方向项写在同一个线性方程中，将每个理论方向
在末尾附加常数 1，形成增广向量
$[({}^{B}\boldsymbol u_i^{*})^{\mathrm T},1]$。相应设计矩阵必须具有满列秩 4，即四个列方向
都具有独立信息；否则轨迹
中心和全部仿射分量不能被唯一辨识。DOWN、FORWARD 和 SIDEWAYS 三类方向不足以满足这一
条件，因此标定流程额外引入 FORWARD\_30 和 FORWARD\_60，以增加前举路径上的方向变化
并改善模型条件。

在线阶段，根据当前大臂接收器位置计算未归一化方向
\begin{equation}
  \tilde{\boldsymbol u}=
  \hat{\boldsymbol M}^{-1}
  ({}^{T}\boldsymbol p_U-\hat{\boldsymbol c}),
  \label{eq:ch4_upper_direction_raw}
\end{equation}
并归一化得到
\begin{equation}
  {}^{B}\hat{\boldsymbol u}=
  \frac{\tilde{\boldsymbol u}}
  {\lVert\tilde{\boldsymbol u}\rVert_2}.
  \label{eq:ch4_upper_direction_normalized}
\end{equation}
若当前点严格位于标定椭球面上，则
$\lVert\tilde{\boldsymbol u}\rVert_2=1$；实际数据受噪声影响通常偏离曲面，因此仍需归一化。
这一过程可以理解为先用逆仿射变换消除安装相关的方向性畸变，再将结果投影回单位方向空间。

矩阵 $\hat{\boldsymbol M}$ 的奇异值表示沿三个主方向的位移尺度。最小奇异值过小时，说明
某一方向几乎没有被参考姿势激励，求逆时该方向的微小位置噪声会被显著放大。为此定义条件数
\begin{equation}
  \kappa(\hat{\boldsymbol M})=
  \frac{\sigma_{\max}(\hat{\boldsymbol M})}
       {\sigma_{\min}(\hat{\boldsymbol M})}.
  \label{eq:ch4_affine_condition_number}
\end{equation}
条件数越接近 1，三个方向的尺度越均衡；条件数过大则表明参考方向分布不足、摆姿不一致或
拟合接近奇异。当前标定质量门槛要求 $\kappa(\hat{\boldsymbol M})\leq4.0$。

为了检测设备移动或当前动作严重超出标定范围，系统将估计方向重新投影到椭球表面：
\begin{equation}
  \hat{\boldsymbol p}_U=
  \hat{\boldsymbol c}+
  \hat{\boldsymbol M}{}^{B}\hat{\boldsymbol u},
  \label{eq:ch4_upper_position_reprojection}
\end{equation}
并定义位置模型误差
\begin{equation}
  e_M=\left\|{}^{T}\boldsymbol p_U-
  \hat{\boldsymbol p}_U\right\|_2.
  \label{eq:ch4_upper_model_error}
\end{equation}
在线运行时，若 $e_M>50\,\mathrm{mm}$，当前帧被保护性丢弃。该误差不等同于人体大臂
方向的绝对真值误差，而是接收器位置与本次安装模型的一致性指标。小的 $e_M$ 只能说明
当前数据与标定轨迹相容；要评价大臂方向的绝对精度，仍需要独立的光学动作捕捉或其他真值
测量系统。

\subsection{前臂方向与坐标对齐}
\label{subsec:ch4_forearm_direction_alignment}

设 $\boldsymbol a_d$ 为前臂接收器局部坐标系中的“肘到腕”方向，$\boldsymbol A_d$ 为发射器
坐标系到人体标准坐标系的旋转。对于第 $i$ 个参考姿势，预测方向为
\begin{equation}
  {}^{B}\hat{\boldsymbol f}_i=
  \boldsymbol A_d\,{}^{T}\boldsymbol R_{F,i}\boldsymbol a_d.
  \label{eq:ch4_forearm_direction_prediction}
\end{equation}

式~\eqref{eq:ch4_forearm_direction_prediction} 表示两步坐标变换：先用
${}^{T}\boldsymbol R_{F,i}$ 把接收器内部的候选方向 $\boldsymbol a_d$ 旋转到发射器坐标系，
再用 $\boldsymbol A_d$ 将其统一表达在人体坐标系。方向轴和对齐矩阵通过交替优化求解：固定
$\boldsymbol a_d$ 时，$\boldsymbol A_d$ 可归结为
Wahba 旋转对齐问题，即寻找一个旋转，使多组测量方向与对应理论方向的平方误差之和最小
~\cite{wahba1965least}；固定 $\boldsymbol A_d$ 时，可将理论方向反变换到各接收器局部坐标系
并取均值，更新 $\boldsymbol a_d$。算法从六个坐标轴正负方向初始化，保留参考姿势方向 RMS
最小的解。

旋转对齐子问题为
\begin{equation}
  \hat{\boldsymbol A}_d=
  \arg\min_{\boldsymbol A\in SO(3)}
  \sum_i\left\|
  \boldsymbol A\,{}^{T}\boldsymbol R_{F,i}\boldsymbol a_d-
  {}^{B}\boldsymbol f_i^{*}
  \right\|_2^2,
  \label{eq:ch4_forearm_wahba_problem}
\end{equation}
通过互协方差矩阵的奇异值分解求得，并显式约束 $\det(\boldsymbol A)=1$。在线阶段的前臂
单位方向为
\begin{equation}
  {}^{B}\hat{\boldsymbol f}=
  \frac{
    \hat{\boldsymbol A}_d\,{}^{T}\boldsymbol R_F\hat{\boldsymbol a}_d
  }{
    \left\|
      \hat{\boldsymbol A}_d\,{}^{T}\boldsymbol R_F\hat{\boldsymbol a}_d
    \right\|_2
  }.
  \label{eq:ch4_forearm_direction_online}
\end{equation}

除了方向映射，系统还建立独立的姿态对齐矩阵 $\boldsymbol A_t$ 和固定安装旋转
$\boldsymbol R_m$，用于计算 wrist roll。这里的“姿态对齐”是把接收器报告的完整三维旋转
转换到人体模型坐标系；$\boldsymbol R_m$ 则记录接收器外壳相对理论前臂坐标系的固定安装角。
在 DOWN 姿势下，理论前臂姿态作为腕部零位；三次 DOWN 重复的平均旋转确定本次安装的固定
旋转。在线计算时先去除肩、肘造成的理论旋转，再去除 $\boldsymbol R_m$，剩余的绕前臂纵轴
分量才被解释为腕部扭转。由此，外壳安装角和前臂方向误差不会直接被误认为 wrist roll。

\subsection{肘角的安装级仿射修正}
\label{subsec:ch4_elbow_affine_correction}

由两个单位方向可以得到几何肘角
\begin{equation}
  q_{e,g}=\arcsin\left(
  \operatorname{clip}
  ( {}^{B}\hat{\boldsymbol u}^{\mathrm T}
  {}^{B}\hat{\boldsymbol f},-1,1)
  \right).
  \label{eq:ch4_geometric_elbow_angle}
\end{equation}

下标 $g$ 表示 geometric，即由两条肢段方向直接得到的几何角。在本研究的零位约定中，
弯肘约 $90^\circ$ 时两个方向近似正交，因而 $q_{e,g}\approx0$；
手臂伸直时二者近似平行，因而 $q_{e,g}\approx90^\circ$。由于两个方向拟合各自存在固定
残差，几何角可能同时存在零位偏差和幅值缩放误差：固定零偏只能整体平移角度，不能改变从
弯曲端到伸直端的量程；单独改变增益又不能校正零位。因此本研究对每次安装分别拟合一维
仿射修正，即“缩放加平移”模型
\begin{equation}
  q_e=\operatorname{clip}
  (s_e q_{e,g}+b_e,\,0,\,\pi/2),
  \label{eq:ch4_elbow_affine_correction}
\end{equation}
其中增益 $s_e$ 和偏置 $b_e$ 由五类前臂参考姿势的理论肘角和几何肘角进行线性最小二乘
拟合。为拒绝异常标定，当前实现要求 $0.75\leq s_e\leq1.50$，且
$|b_e|\leq20^\circ$。若拟合不可辨识或超出合理范围，系统保留未经修正的几何肘角，同时
输出警告。

\subsection{标定质量门槛与实机资格}
\label{subsec:ch4_calibration_quality_gates}

标定参数并非只要能够求解就可进入真机。系统在保存前检查表~\ref{tab:ch4_quality_gates}
所列工程门槛。
\begin{table}[htbp]
  \centering
  \caption{安装级标定质量门槛}
  \label{tab:ch4_quality_gates}
  \begin{tabularx}{\linewidth}{@{}XllX@{}}
    \toprule
    检查项 & 建议/诊断阈值 & 硬限制 & 不通过时的处理 \\
    \midrule
    参考姿势重复差异 & --- & $15^\circ$ & 重采当前姿势 \\
    大臂方向参考误差 & $10^\circ$ & $15^\circ$ & 警告/拒绝保存 \\
    前臂方向参考误差 & $10^\circ$；真机资格 $20^\circ$ & $25^\circ$ & 警告、仅诊断或拒绝 \\
    前臂腕部零位重复误差 & --- & $15^\circ$ & 拒绝保存 \\
    仿射矩阵条件数 & 越接近 1 越好 & 4.0 & 拒绝保存 \\
    仿射位置模型 RMS/最大误差 & $30/50\,\mathrm{mm}$ & $35/65\,\mathrm{mm}$ & 警告/拒绝保存 \\
    \bottomrule
  \end{tabularx}
\end{table}

表中的“建议/诊断阈值”用于提示边缘质量，“硬限制”决定标定文件能否保存，两者不应简单
理解为“仿真门槛”和“实机门槛”。仿真与 dry-run 仍可使用通过硬限制但质量边缘的标定；
真机资格另由更严格的前臂方向条件控制。若前臂方向误差超过 $20^\circ$ 但仍未超过硬限制
$25^\circ$，标定可以作为诊断资料保存，
但会被标记为 \texttt{dry\_run\_only: true}。该标定只允许用于 dry-run 和离线诊断，
重定向节点拒绝向真机主题发布关节目标。该设计把“算法能够产生数值”和“标定具备实机
资格”明确分开。

这些阈值是本项目的工程安全标准，并非电磁跟踪器或人形机器人领域的通用精度标准。参考
姿势残差描述标定模型与人工构建姿势之间的一致性，不能替代独立视觉系统测得的人体姿态
真值误差。

\section{在线人体手臂姿态估计}
\label{sec:ch4_online_arm_estimation}

安装标定完成后，每一对同步位姿首先被转换为三个中间几何量：大臂单位方向
${}^{B}\hat{\boldsymbol u}$、前臂单位方向 ${}^{B}\hat{\boldsymbol f}$，以及经过腕部姿态
对齐的前臂接收器旋转 ${}^{B}\boldsymbol R_F$。其中，大臂方向由
式~\eqref{eq:ch4_upper_direction_raw}--\eqref{eq:ch4_upper_direction_normalized} 对当前大臂
接收器位置进行逆仿射映射得到；前臂方向由式~\eqref{eq:ch4_forearm_direction_online} 对当前
前臂接收器姿态和已标定局部轴进行变换得到；“腕部姿态对齐”则指用 $\boldsymbol A_t$ 将完整
接收器旋转转换到人体模型坐标系。五个关节量并不是同时进行黑箱优化，而是
按照几何依赖关系逐级求解：先由大臂方向确定 $q_p,q_r$，再由两条肢段方向确定几何肘角
$q_{e,g}$ 和肩部轴向角 $q_y$，最后从前臂残余旋转中提取 $q_w$。几何肘角随后经过本次
安装的增益和偏置修正，形成最终输出 $q_e$。

这种求解顺序保留了每个变量的物理意义，也使不可观测条件能够被显式检测。它与输入信息的
对应关系可概括为
\begin{equation}
  ({}^{T}\boldsymbol p_U,{}^{T}\boldsymbol R_F)
  \longrightarrow
  ({}^{B}\hat{\boldsymbol u},{}^{B}\hat{\boldsymbol f},
   {}^{B}\boldsymbol R_F)
  \longrightarrow
  (q_p,q_r,q_y,q_{e,g},q_w)
  \longrightarrow
  (q_p,q_r,q_y,q_e,q_w).
  \label{eq:ch4_online_estimation_chain}
\end{equation}

\subsection{肩部 pitch 与 roll}
\label{subsec:ch4_shoulder_pitch_roll}

得到大臂单位方向
\begin{equation}
  {}^{B}\hat{\boldsymbol u}=
  \begin{bmatrix}u_x&u_y&u_z\end{bmatrix}^{\mathrm T}
  \label{eq:ch4_upper_direction_components}
\end{equation}
后，由于 $\boldsymbol R_z(q_y)$ 绕零位大臂轴旋转，不改变
$\boldsymbol a_{u,0}=[0,0,-1]^{\mathrm T}$，大臂方向只与 $q_p,q_r$ 有关。展开
式~\eqref{eq:ch4_segment_directions} 可得
\begin{equation}
  {}^{B}\hat{\boldsymbol u}=
  \begin{bmatrix}
    -\sin q_p\cos q_r\\
    \sin q_r\\
    -\cos q_p\cos q_r
  \end{bmatrix}.
  \label{eq:ch4_upper_direction_expanded}
\end{equation}
这里的“零位大臂轴”就是式~\eqref{eq:ch4_zero_segment_axes} 定义的
$\boldsymbol a_{u,0}$。由于该向量位于 $z$ 轴上，绕 $z$ 轴旋转不会改变它，即
$\boldsymbol R_z(q_y)\boldsymbol a_{u,0}=\boldsymbol a_{u,0}$；随后依次左乘
$\boldsymbol R_x(q_r)$ 和 $\boldsymbol R_y(q_p)$，直接进行矩阵乘法即可得到
式~\eqref{eq:ch4_upper_direction_expanded}。

因此，肩部 pitch 和 roll 的主值可写为
\begin{equation}
  q_p=\operatorname{atan2}(-u_x,-u_z),
  \qquad
  q_r=\arcsin(\operatorname{clip}(u_y,-1,1)).
  \label{eq:ch4_pitch_roll_inverse}
\end{equation}

“主值”表示反三角函数在约定区间内返回的代表解，因为同一方向可对应相差 $2\pi$ 的周期
等价角。由式~\eqref{eq:ch4_upper_direction_expanded} 有
$u_y=\sin q_r$，故 $q_r$ 可由反正弦得到；而 $-u_x$ 与 $-u_z$ 分别包含
$\sin q_p\cos q_r$ 和 $\cos q_p\cos q_r$，用双参数反正切可同时保留象限信息并消去共同
因子 $\cos q_r$。$\operatorname{clip}(u_y,-1,1)$ 只用于消除浮点误差造成的轻微越界，
例如数值上出现 $1+10^{-12}$ 时避免 $\arcsin$ 失效，并非裁剪正常人体动作。

由于 \texttt{atan2} 的输出位于 $[-\pi,\pi]$，系统将当前结果展开到最接近上一帧的等价角，
避免跨越 $\pm\pi$ 时产生整周跳变。

当大臂接近正侧方时，$u_x$ 和 $u_z$ 同时接近零，pitch 对测量噪声高度敏感。定义
\begin{equation}
  o_p=\sqrt{u_x^2+u_z^2}=|\cos q_r|
  \label{eq:ch4_pitch_observability}
\end{equation}
为 pitch 可观测度。在线算法采用带滞回的两状态机制：当 $o_p\leq0.50$ 时进入侧举奇异区，
将 pitch 切换到统一的标准解 $q_p=0$；只有当手臂返回水平面以下且 $o_p\geq0.65$ 时才
重新采用测量解。与此同时，roll 使用标准 pitch 坐标系下的 \texttt{atan2} 表达，使其
能够连续通过 $90^\circ$，而不是被 \texttt{asin} 折回到主值范围。该处理保证从 DOWN
直接侧举和从 FORWARD 转向侧举时，同一最终人体姿势映射到同一机器人目标。

由式~\eqref{eq:ch4_upper_direction_expanded} 可见，$u_x$ 和 $u_z$ 都含有因子
$\cos q_r$；当侧举接近 $|q_r|=90^\circ$ 时该因子趋近零，pitch 的分子和分母同时退化。
阈值 0.50 和 0.65 分别约对应 $|q_r|=60^\circ$ 和 $49^\circ$，是当前工程为避免临界点
抖动设置的进入/退出滞回阈值，而不是人体生理常数。进入奇异区后固定 $q_p=0$，相当于从
无穷多个等价 pitch--roll 表示中选择统一标准解；退出阈值更高且要求回到水平面以下，使状态
不会在噪声作用下频繁切换。

\subsection{肘部伸展量}
\label{subsec:ch4_online_elbow_angle}

在运动学模型中，大臂和前臂方向分别为
${}^{B}\hat{\boldsymbol u}={}^B\boldsymbol R_S\boldsymbol a_{u,0}$ 和
${}^{B}\hat{\boldsymbol f}={}^B\boldsymbol R_S
\boldsymbol R_y(q_{e,g})\boldsymbol a_{f,0}$。利用旋转矩阵保持内积的性质，有
\begin{align}
  ({}^{B}\hat{\boldsymbol u})^{\mathrm T}
  {}^{B}\hat{\boldsymbol f}
  &=\boldsymbol a_{u,0}^{\mathrm T}
    \boldsymbol R_y(q_{e,g})\boldsymbol a_{f,0}
  \nonumber\\
  &=\sin q_{e,g}.
  \label{eq:ch4_elbow_dot_product_derivation}
\end{align}
旋转矩阵满足 $\boldsymbol R^{\mathrm T}\boldsymbol R=\boldsymbol I$，因此对任意向量
$\boldsymbol a,\boldsymbol b$ 都有
$(\boldsymbol R\boldsymbol a)^{\mathrm T}(\boldsymbol R\boldsymbol b)
=\boldsymbol a^{\mathrm T}\boldsymbol b$，即共同旋转不改变夹角或内积。在本式中，肩部旋转
${}^{B}\boldsymbol R_S$ 因而被消去；再代入
$\boldsymbol a_{u,0}=[0,0,-1]^{\mathrm T}$、
$\boldsymbol a_{f,0}=[1,0,0]^{\mathrm T}$ 和 $\boldsymbol R_y(q_{e,g})$ 的矩阵形式，内积
恰好等于 $\sin q_{e,g}$。

因此，式~\eqref{eq:ch4_geometric_elbow_angle} 可以直接由两个肢段方向求得，而不依赖肩部
三个角的具体数值。在本研究的约定下，$q_{e,g}=0$ 对应约 $90^\circ$ 屈肘，
$q_{e,g}=\pi/2$ 对应伸直。为避免与临床文献中常用的“屈曲角”混淆，本文将其称为肘部
伸展量。

这种零位来自当前 G1 左臂模型与重定向接口的关节语义：机器人模型零位附近，大臂向下而
前臂向前，二者约成直角；正向 $q_e$ 将前臂逐渐转到与大臂同向。因此本文的 $q_e$ 是
“从约 $90^\circ$ 弯肘零位出发的伸展量”，不等同于临床上通常从完全伸直定义的肘屈曲角。

几何肘角直接反映当前两条估计方向之间的关系。最终发布的肘部变量使用
式~\eqref{eq:ch4_elbow_affine_correction} 进行安装级增益和零偏修正。这里应区分
$q_{e,g}$ 与 $q_e$：前者用于保持几何分解与测量方向一致，后者是面向机器人目标的校正后
输出。
必须保留两者，是因为 wrist 残余姿态分解要求理论前臂方向与当前实测大臂、前臂方向严格
一致，故使用未修正的 $q_{e,g}$；而输出给重定向层的肘角需要补偿安装造成的尺度和零偏，故
使用 $q_e$。若提前用 $q_e$ 构造理论前臂姿态，标定补偿量会被错误地吸收到 wrist twist 中。

\subsection{肩部 yaw}
\label{subsec:ch4_shoulder_yaw}

大臂方向本身不能确定绕大臂轴的旋转。系统使用前臂在垂直于大臂方向平面内的分量估计
肩部 yaw：
\begin{equation}
  \boldsymbol f_{\perp}=
  {}^{B}\hat{\boldsymbol f}-
  \left({}^{B}\hat{\boldsymbol f}^{\mathrm T}
  {}^{B}\hat{\boldsymbol u}\right)
  {}^{B}\hat{\boldsymbol u}.
  \label{eq:ch4_forearm_perpendicular_component}
\end{equation}

若 $\lVert\boldsymbol f_{\perp}\rVert$ 足够大，将其归一化，并与在当前 pitch、roll 下的零
yaw 参考方向比较。几何上，$\boldsymbol f_{\perp}$ 是前臂方向在大臂法平面内的投影，
它确定了“肩--肘--腕”平面绕大臂纵轴的朝向。设零 yaw 参考方向投影并归一化后为
$\boldsymbol f_{0,\perp}$，则
\begin{equation}
  q_y=\operatorname{atan2}
  \left(
  -{}^{B}\hat{\boldsymbol u}^{\mathrm T}
  (\boldsymbol f_{0,\perp}\times\boldsymbol f_{\perp}),
  \boldsymbol f_{0,\perp}^{\mathrm T}\boldsymbol f_{\perp}
  \right).
  \label{eq:ch4_shoulder_yaw_inverse}
\end{equation}

投影的作用是去除前臂沿大臂轴的分量，只保留能够描述肘平面绕大臂轴朝向的分量。归一化
则去除投影长度，只比较方向；但只有投影长度足够大时，归一化才不会放大噪声。零 yaw 参考
方向 $\boldsymbol f_{0,\perp}$ 由当前 $q_p,q_r$ 和 $q_y=0$ 的运动学模型生成，因此比较二者
即可得到实际肘平面相对零 yaw 肘平面的有符号转角。式~\eqref{eq:ch4_shoulder_yaw_inverse}
中，点积给出余弦项，沿大臂轴的叉积投影给出带符号正弦项，$\operatorname{atan2}$ 将二者
合成为保留象限的有符号角。前置负号用于匹配 G1 肩 yaw 的正方向约定。

\begin{figure}[htbp]
  \centering
  \fbox{\parbox{0.92\linewidth}{\centering
    左图：在垂直于大臂方向 ${}^{B}\hat{\boldsymbol u}$ 的平面内，画出前臂投影
    $\boldsymbol f_{\perp}$、零 yaw 参考 $\boldsymbol f_{0,\perp}$ 和有符号角 $q_y$；
    右图：对比弯肘时投影较长、yaw 与 wrist 可分辨，以及直臂时投影趋近零、两种轴向旋转
    耦合的情形。}}
  \caption{肩部 yaw 的平面投影求解及直臂不可观测性（待替换为正式绘图）}
  \label{fig:ch4_yaw_projection_observability}
\end{figure}

与 pitch 相同，yaw 也展开到最接近上一帧的等价角。由于旋转 $q_y$ 与
$q_y+2k\pi$（$k\in\mathbb Z$）表示同一姿态，展开只选择最连续的数值分支，并不改变物理
旋转。

\subsection{前臂旋前/旋后}
\label{subsec:ch4_forearm_pronation_supination}

当前实现使用与两条测量方向（大臂单位方向 ${}^{B}\hat{\boldsymbol u}$ 和前臂单位方向
${}^{B}\hat{\boldsymbol f}$）直接一致的几何肘角 $q_{e,g}$ 构造不含 wrist roll 的理论
前臂旋转
\begin{equation}
  {}^{B}\boldsymbol R_{E,g}=
  \boldsymbol R_y(q_p)\boldsymbol R_x(q_r)
  \boldsymbol R_z(q_y)\boldsymbol R_y(q_{e,g}).
  \label{eq:ch4_forearm_rotation_without_wrist}
\end{equation}
当前接收器在人体标准坐标系中的旋转为
\begin{equation}
  {}^{B}\boldsymbol R_F=
  \boldsymbol A_t\,{}^{T}\boldsymbol R_F.
  \label{eq:ch4_forearm_sensor_body_rotation}
\end{equation}
去除理论前臂旋转和固定安装旋转后，残余旋转为
\begin{equation}
  \boldsymbol R_{res}=
  ({}^{B}\boldsymbol R_{E,g})^{\mathrm T}
  {}^{B}\boldsymbol R_F\boldsymbol R_m^{\mathrm T}.
  \label{eq:ch4_wrist_residual_rotation}
\end{equation}

上述三式按“预测---观测---求残差”的顺序工作。式
~\eqref{eq:ch4_forearm_rotation_without_wrist} 用已估计的肩、肘量预测前臂在未转掌时的姿态；
式~\eqref{eq:ch4_forearm_sensor_body_rotation} 将传感器观测变换到同一人体坐标系；式
~\eqref{eq:ch4_wrist_residual_rotation} 左乘理论姿态的逆、右乘固定安装旋转的逆，从观测中
去除肩部运动、肘部运动和外壳安装角。理想剩余量只包含前臂绕自身纵轴的旋转。本文沿用
机器人关节名将该量称为 wrist roll，但其人体意义是前臂旋前/旋后，并非手腕屈伸。

等价地，前臂接收器姿态模型可写为
\begin{equation}
  {}^{B}\boldsymbol R_F
  \approx{}^{B}\boldsymbol R_{E,g}
  \boldsymbol R_x(q_w)\boldsymbol R_m.
  \label{eq:ch4_forearm_sensor_attitude_model}
\end{equation}
其中，$\boldsymbol R_m$ 表示本次安装中接收器与前臂模型坐标系之间的固定旋转。

理想情况下，$\boldsymbol R_{res}$ 仅包含绕前臂局部 $+x$ 轴的旋转。为抑制小幅残余 swing
误差，系统将 $\boldsymbol R_{res}$ 转为四元数
$\boldsymbol q_{res}=[q_0,q_x,q_y,q_z]^{\mathrm T}$，只保留向量部分在 $+x$ 轴上的投影，
得到
\begin{equation}
  q_w=\operatorname{wrap}_{[-\pi,\pi)}
  \left(2\operatorname{atan2}(q_x,q_0)\right).
  \label{eq:ch4_wrist_twist_angle}
\end{equation}
然后将结果展开到上一帧附近，以维持时间连续性。这一过程等价于围绕前臂纵轴的
swing--twist 分解：swing 改变纵轴指向，twist 则绕该纵轴旋转；只保留 $q_x$ 相当于提取
局部 $+x$ 轴上的 twist 分量。此类分解常用于把三维关节旋转分离为摆动和轴向扭转
~\cite{murthy2019joint}。结果展开的原因与 yaw 相同：四元数给出的主角位于有限区间，相邻
帧若跨过 $\pm\pi$ 会在数值上跳变约 $2\pi$，选择离上一帧最近的等价角可保持连续。
wrist 提取完成后，系统才将 $q_{e,g}$ 修正为最终输出 $q_e$；因此，
肘角的经验增益和零偏不会被误解释为额外的前臂轴向旋转。

\section{不可观测区域与稳健性处理}
\label{sec:ch4_observability_robustness}

\subsection{直臂状态下 yaw 与 wrist 的耦合}
\label{subsec:ch4_yaw_wrist_unobservable}

当肘部接近完全伸直时，大臂和前臂方向近似共线，
$\lVert\boldsymbol f_{\perp}\rVert\rightarrow0$。此时，仅凭两个肢段方向无法区分“肩部
绕大臂轴旋转”和“前臂自身旋转”；同一物理姿态可以对应多组 $(q_y,q_w)$。这是测量模型
的几何不可观测性，而不是数值求解精度不足。

原因是直臂时肩 yaw 轴与前臂 twist 轴近似重合，末端姿态主要取决于二者的组合量。对任意
小角度 $\delta$，参数对
\begin{equation}
  (q_y,q_w)\sim(q_y+\delta,q_w-\delta)
  \label{eq:ch4_yaw_wrist_equivalent_pairs}
\end{equation}
可产生相同或近似相同的总轴向旋转，因而无法由当前两接收器数据唯一判断该旋转应归属于
肩部还是前臂。弯肘后两条旋转轴不再共线，图
~\ref{fig:ch4_yaw_projection_observability} 所示投影重新出现，两个变量才恢复可分辨性。

系统定义
\begin{equation}
  o_y=\lVert\boldsymbol f_{\perp}\rVert_2
  \label{eq:ch4_yaw_observability}
\end{equation}
为 yaw 可观测度。由于 $\boldsymbol u$ 和 $\boldsymbol f$ 均为单位向量，且
$\boldsymbol u^{\mathrm T}\boldsymbol f=\sin q_{e,g}$，可进一步得到
\begin{equation}
  o_y=\sqrt{1-(\boldsymbol u^{\mathrm T}\boldsymbol f)^2}
  =|\cos q_{e,g}|.
  \label{eq:ch4_yaw_observability_elbow_relation}
\end{equation}
因此，这一退化不是经验上偶然出现的数值问题，而是由五自由度模型的几何结构直接决定。
弯肘约 $90^\circ$ 时 $q_{e,g}\approx0$、$o_y\approx1$，肘平面方向清晰；手臂趋于伸直时
$q_{e,g}\rightarrow\pi/2$、$o_y\rightarrow0$，肩 yaw 轴与前臂 twist 轴趋于重合，测量
中只能保留两者的组合旋转而不能唯一分离。

当 $o_y<0.90$ 时，保持最后一次可靠的 $q_y$ 和 $q_w$，避免把微小方向
噪声放大为几十度的轴向旋转。由于本模型中 $o_y$ 随肘部伸直而下降，该阈值会在完全伸直
之前提前冻结轴向角。系统恢复到明显弯肘姿态后，再继续更新 yaw 和 wrist。

因此，直臂前举和直臂侧举可以用于评价大臂方向及肘角，但不适合
把直臂主动转掌作为 wrist 精度的主要证据。wrist 单关节实验应保持肘部弯曲约 $60^\circ$
至 $90^\circ$，使肩部轴向旋转和前臂旋转保持可分辨。

\subsection{周期角展开与时间连续性}
\label{subsec:ch4_temporal_conditioning}

肩部 yaw、wrist roll 和其他周期角首先展开到最接近上一帧的 $2\pi$ 等价分支。反三角函数
通常只返回 $[-\pi,\pi)$ 内的主值，因此真实角度从 $+\pi$ 连续增加时，下一帧数值会变成
接近 $-\pi$；两帧物理上很接近，数值却相差近 $2\pi$。角展开在
$\theta+2k\pi$ 的等价集合中选择最接近上一帧的分支，从而避免该边界产生整周跳变。侧举
奇异区和直臂不可观测区的内部状态也只在满足带滞回的
重新观测条件后更新，防止临界点附近反复切换。

本章输出的是经过可观测性处理的五自由度人体估计，不在此处施加 G1 关节限位、机器人
命令低通滤波或真机速度限制。上述机器人目标调理统一放在第5章说明，从而把“人体状态
估计”和“机器人命令生成”分开。

\subsection{无效数据与看门狗}
\label{subsec:ch4_invalid_data_watchdog}

在线估计过程中，以下情况不会产生新的关节目标：
\begin{itemize}
  \item 接收器数据陈旧或双流时间差过大；
  \item 原始四元数模长异常；
  \item 两个话题的坐标系标识不一致；
  \item 大臂仿射位置模型误差超过 $50\,\mathrm{mm}$；
  \item 标定矩阵无效或当前方向向量接近零。
\end{itemize}

“看门狗”是监测数据或程序是否在规定时间内持续更新的超时机制。若连续
$0.25\,\mathrm{s}$ 没有新的有效同步数据对，重定向节点发出看门狗警告并停止更新
目标。真机桥还具有独立的 UDP 目标超时和控制权释放机制。因此，采集层、重定向层与执行
层分别检查数据新鲜度，避免单点故障持续驱动机器人。

\section{算法流程与实现}
\label{sec:ch4_algorithm_implementation}

一次完整在线估计可归纳为算法 4-1。

\begin{quote}
\textbf{算法 4-1：双电磁接收器人体左臂姿态估计}

\textbf{输入：}同步接收器位姿
$({}^{T}\boldsymbol p_U,{}^{T}\boldsymbol R_U,
{}^{T}\boldsymbol p_F,{}^{T}\boldsymbol R_F)$，当次安装标定参数，上一帧角度。

\textbf{输出：}五自由度人体手臂估计
$\hat{\boldsymbol q}_h=[q_p,q_r,q_y,q_e,q_w]^{\mathrm T}$。

\begin{enumerate}[label=\arabic*.]
  \item 检查数据新鲜度、双流时间差、四元数模长和共同时间戳；
  \item 根据当次标定选择原始或转置旋转约定；
  \item 用大臂接收器位置和仿射椭球模型计算 ${}^{B}\hat{\boldsymbol u}$，并检查模型误差；
  \item 用前臂方向对齐矩阵和局部方向轴计算 ${}^{B}\hat{\boldsymbol f}$；
  \item 根据 ${}^{B}\hat{\boldsymbol u}$ 计算 pitch 和 roll，并处理侧举奇异区；
  \item 根据两条肢段方向的内积计算几何肘角 $q_{e,g}$；
  \item 计算 yaw 可观测度；可观测时求解 yaw，否则保持上一帧 yaw；
  \item 使用 $q_{e,g}$ 去除肩部、肘部和接收器固定安装旋转，提取 wrist twist；yaw
        不可观测时同时保持上一帧 wrist；
  \item 对 $q_{e,g}$ 应用本次安装的增益和零偏，得到最终肘部输出 $q_e$；
  \item 将周期角展开到上一帧附近，输出五自由度人体估计；
  \item 记录原始方向、可观测度、原始角、修正角和质量指标，用于离线评价。
\end{enumerate}
\end{quote}

实现中，双接收器采集、标定和在线估计分别由 \texttt{dual\_em\_tracker\_node}、
\texttt{dual\_calibration\_node} 和 \texttt{dual\_retargeting\_node} 完成。纯几何运算被
独立放置在 \texttt{retargeting\_kinematics.py} 和
\texttt{dual\_retargeting\_kinematics.py} 两个无 ROS 依赖的运动学模块中，便于通过单元
测试验证坐标变换、奇异区分支和角度连续性。每轮运行结束时，系统将全频率诊断数据保存为
NPZ（NumPy compressed archive，NumPy 压缩数组归档）文件。该格式可在一个压缩文件中按
名称保存多个同长度数组，便于保持时间序列字段之间的一一对应。记录内容包括接收器原始位姿、人体
肢段方向、五个原始关节角、肘角修正前后数值、pitch/yaw 可观测度以及最终命令。这些数据
为第6章实验设计和第7章结果分析提供统一的数据来源。

\section{本章小结}
\label{sec:ch4_summary}

本章提出并实现了面向重复佩戴的双电磁接收器人体左臂估计方法。该方法没有假设传感器在
每次佩戴后保持相同朝向，而是通过静态质量检查、接收器身份检查、独立前臂旋转轴动作和
七个参考姿势，重新辨识当前安装的坐标对齐、肢段方向轴、大臂位置模型以及肘角修正参数。

本章进一步说明了测量信息与模型自由度之间的关系：大臂单位方向提供两个独立自由度，前臂
完整姿态提供三个独立自由度，从而在一般姿态下支持五自由度状态的逐级几何求解。这五个量
是与 G1 关节语义一致的模型广义坐标，而非完整人体上肢的解剖学真值。

为降低软组织和绑带运动对大臂姿态测量的影响，系统使用仿射位置--方向模型恢复大臂方向；
该模型在几何上把单位方向球面映射为安装相关的椭球轨迹，但并不假设人体肩关节本身为椭球。
前臂方向与旋前/旋后使用相互独立的轴和对齐参数。在线阶段，系统从两个肢段方向计算肩部
pitch、roll、yaw 和几何肘角，通过残余旋转的 swing--twist 分解得到 wrist roll，再将
安装级增益和零偏作用于最终肘部输出。针对侧举
时 pitch 退化和直臂时 yaw/wrist 耦合，本章分别设计了带滞回的标准分支和最后可靠值保持
策略。最后，标定质量门槛、同步检查、位置模型误差和看门狗共同阻止低质量数据进入后续
控制链。

下一章将在本章输出的五自由度人体动作基础上，介绍 G1 关节目标生成、关节和速度限制、
躯干碰撞保护以及真机 Arm SDK 安全桥。

